%% ============================================================
%%  Exercises -- Homework 2 (assignment family).
%%  Rudin Chapter 2 topology. Entries assembled from
%%  exercises-build/hw2/ (drafted + firewalled by workflow).
%%  Subfile of main.tex.
%% ============================================================
\documentclass[main.tex]{subfiles}
\begin{document}

\beginunittoc{Homework 2}{HOMEWORK 2}{Rudin Chapter 2 --- Limit Points, Compactness, and Perfect, Connected, and Separable Sets}{The Topology of Metric Spaces}{Homework 2 \textendash\ The Topology of Metric Spaces}
\hypertarget{hw-2}{}

\begin{lecturecontents}{Problems in This Set}
  \item \hyperlink{prob:1}{Problem 1\quad Rudin Ch.\,2, Ex.\,1 --- The Empty Set Is a Subset of Every Set}
  \item \hyperlink{prob:2}{Problem 2\quad Rudin Ch.\,2, Ex.\,2 --- The Algebraic Numbers Are Countable}
  \item \hyperlink{prob:3}{Problem 3\quad Rudin Ch.\,2, Ex.\,3 --- Transcendental Numbers Exist}
  \item \hyperlink{prob:4}{Problem 4\quad Rudin Ch.\,2, Ex.\,4 --- Are the Irrationals Countable?}
  \item \hyperlink{prob:5}{Problem 5\quad Rudin Ch.\,2, Ex.\,5 --- A Bounded Set with Exactly Three Limit Points}
  \item \hyperlink{prob:6}{Problem 6\quad Rudin Ch.\,2, Ex.\,6 --- $E'$ Is Closed; $E$ and $\overline E$ Share Limit Points}
  \item \hyperlink{prob:7}{Problem 7\quad Rudin Ch.\,2, Ex.\,7 --- Closures of Unions}
  \item \hyperlink{prob:8}{Problem 8\quad Rudin Ch.\,2, Ex.\,8 --- Limit Points of Open and Closed Sets in $\mathbb R^2$}
  \item \hyperlink{prob:9}{Problem 9\quad Rudin Ch.\,2, Ex.\,9 --- The Interior $E^{\circ}$}
  \item \hyperlink{prob:10}{Problem 10\quad Rudin Ch.\,2, Ex.\,10 --- The Discrete Metric}
  \item \hyperlink{prob:11}{Problem 11\quad Rudin Ch.\,2, Ex.\,11 --- Which Functions Are Metrics on $\mathbb R$?}
  \item \hyperlink{prob:12}{Problem 12\quad Rudin Ch.\,2, Ex.\,12 --- $\{0\}\cup\{1/n\}$ Is Compact, Directly}
  \item \hyperlink{prob:13}{Problem 13\quad Rudin Ch.\,2, Ex.\,13 --- A Compact Set with Countable Limit-Point Set}
  \item \hyperlink{prob:14}{Problem 14\quad Rudin Ch.\,2, Ex.\,14 --- An Open Cover of $(0,1)$ with No Finite Subcover}
  \item \hyperlink{prob:15}{Problem 15\quad Rudin Ch.\,2, Ex.\,15 --- ``Closed'' and ``Bounded'' Do Not Replace ``Compact''}
  \item \hyperlink{prob:16}{Problem 16\quad Rudin Ch.\,2, Ex.\,16 --- A Closed, Bounded, Non-Compact Set in $\mathbb Q$}
  \item \hyperlink{prob:17}{Problem 17\quad Rudin Ch.\,2, Ex.\,17 --- The Digits-$4$-and-$7$ Set}
  \item \hyperlink{prob:18}{Problem 18\quad Rudin Ch.\,2, Ex.\,18 --- A Perfect Set with No Rationals}
  \item \hyperlink{prob:19}{Problem 19\quad Rudin Ch.\,2, Ex.\,19 --- Separated Sets and Connectedness}
  \item \hyperlink{prob:20}{Problem 20\quad Rudin Ch.\,2, Ex.\,20 --- Closures and Interiors of Connected Sets}
  \item \hyperlink{prob:21}{Problem 21\quad Rudin Ch.\,2, Ex.\,21 --- Separated Sets Along a Segment}
  \item \hyperlink{prob:22}{Problem 22\quad Rudin Ch.\,2, Ex.\,22 --- $\mathbb R^k$ Is Separable}
  \item \hyperlink{prob:23}{Problem 23\quad Rudin Ch.\,2, Ex.\,23 --- A Separable Metric Space Has a Countable Base}
  \item \hyperlink{prob:24}{Problem 24\quad Rudin Ch.\,2, Ex.\,24 --- Limit-Point Property $\Rightarrow$ Separable}
  \item \hyperlink{prob:25}{Problem 25\quad Rudin Ch.\,2, Ex.\,25 --- Compact $\Rightarrow$ Countable Base}
  \item \hyperlink{prob:26}{Problem 26\quad Rudin Ch.\,2, Ex.\,26 --- Limit-Point Property $\Rightarrow$ Compact}
  \item \hyperlink{prob:27}{Problem 27\quad Rudin Ch.\,2, Ex.\,27 --- Condensation Points}
  \item \hyperlink{prob:28}{Problem 28\quad Rudin Ch.\,2, Ex.\,28 --- Closed Sets as Perfect Plus Countable}
  \item \hyperlink{prob:29}{Problem 29\quad Rudin Ch.\,2, Ex.\,29 --- Open Sets in $\mathbb R$ as Countable Unions of Segments}
  \item \hyperlink{prob:30}{Problem 30\quad Rudin Ch.\,2, Ex.\,30 --- A Baire-Category Result}
\end{lecturecontents}

% ------------------------------------------------------------
% ------------------------------------------------------------
\prob{1}{Rudin, Chapter 2, Exercise 1}
Prove that the empty set is a subset of every set.

\begin{SolutionBlue}
Let $A$ be any set. The claim $\varnothing\subseteq A$ unfolds, by the definition of subset, into the
implication
\[
x\in\varnothing \quad\Longrightarrow\quad x\in A,
\]
asserted for every $x$. Its hypothesis $x\in\varnothing$ is false for every $x$, since $\varnothing$
has no elements. An implication with a false hypothesis is true, so the implication holds for every
$x$, with no instance left to check. Hence $\varnothing\subseteq A$, and since $A$ was arbitrary, the
empty set is a subset of every set.
\end{SolutionBlue}

\begin{Reflection}
The word that fixes the strategy is \emph{subset}. A claim ``$\varnothing\subseteq A$'' is not a claim
about a single object but a universal implication---``for every $x$, if $x\in\varnothing$ then $x\in
A$''---which is exactly how a subset statement is built (\xrefL{0.3.6}), and the standard way to prove
one is to take an arbitrary element of the left-hand set and show it lies in the right
(\xrefL{0.5.10}). The second clue is the set on the left: it is $\varnothing$, the one set with no
elements at all (\xrefL{0.3.4}). The reflex those two clues should trigger is to look at the
\emph{hypothesis} of the implication rather than its conclusion, because that hypothesis is the thing
the empty set makes impossible.

It helps to believe it first on a small case. Asking whether $\varnothing\subseteq\{1,2\}$ means asking
whether every element of $\varnothing$ lies in $\{1,2\}$; there is no element of $\varnothing$ to
inspect, so there is nothing that could go wrong, and the answer is yes by default. The same picture
works for any target set, which is already the whole theorem.

The argument is short, and each move rests on the one before:
\begin{enumerate}[leftmargin=1.7em,itemsep=3pt,topsep=2pt,label=\textnormal{(\arabic*)}]
  \item fix an arbitrary set $A$ and unfold $\varnothing\subseteq A$ into the implication ``$x\in
        \varnothing\Rightarrow x\in A$'' for every $x$, the only thing a subset claim means
        (\xrefL{0.3.6});
  \item read the hypothesis $x\in\varnothing$: it is false for every $x$, because the empty set has no
        elements (\xrefL{0.3.4});
  \item conclude the implication is true, since an implication with a false hypothesis is true---this is
        vacuous truth (\xrefL{0.1.6}, on the definition of implication \xrefL{0.1.5}), and it is the
        single load-bearing step;
  \item since $A$ was named with no special property, the conclusion holds for every set.
\end{enumerate}

The rigor check is to see that there is genuinely no gap rather than a missing case. The proof never
produces an element of $\varnothing$---it cannot, and that is the point; it argues that the implication
to be verified has no false instance, so it is true with nothing to verify. Drop the feature that
$\varnothing$ is empty and the reasoning collapses at once: if the left-hand set had even one element
outside $A$, the implication would have a true hypothesis and a false conclusion, and the subset claim
would fail. So emptiness is exactly the hypothesis doing the work, which is why the result records a
basic property of $\varnothing$ alone (\xrefL{0.3.8}).

The transferable move is how to handle any universal statement quantified over a set that might be
empty: it is automatically true there, because a ``for all'' over nothing has no counterexample. This
is the same vacuous-truth reflex that lets a proof safely open ``take an arbitrary element of $E$''
even when $E$ could turn out to be empty (\xrefL{0.3.7})---the empty case needs no separate argument,
it is handled for free.
\end{Reflection}

% ------------------------------------------------------------
% ------------------------------------------------------------
\prob{2}{Rudin, Chapter 2, Exercise 2}
A complex number $z$ is said to be \emph{algebraic} if there are integers $a_0,\dots,a_n$, not all
zero, such that
\[
a_0z^n+a_1z^{n-1}+\cdots+a_{n-1}z+a_n=0.
\]
Prove that the set of all algebraic numbers is countable. \emph{Hint:} For every positive integer $N$
there are only finitely many equations with
\[
n+|a_0|+|a_1|+\cdots+|a_n|=N.
\]

\begin{SolutionBlue}
To each integer polynomial $a_0z^n+a_1z^{n-1}+\cdots+a_n$ with coefficients not all zero, assign the
\emph{height}
\[
N=n+|a_0|+|a_1|+\cdots+|a_n|\ge 1.
\]

Fix a positive integer $N$. A polynomial of height $N$ has degree $n\le N$, and every coefficient
satisfies $|a_i|\le N$; so there are at most $N+1$ choices for $n$ and at most $(2N+1)^{n+1}$ choices
for the coefficients. Hence only finitely many integer polynomials have height $N$. Each such
polynomial is nonzero of degree $\le N$, so it has at most $N$ complex roots. Let $A_N$ be the set of
all roots of all height-$N$ polynomials; being a finite union of finite sets, $A_N$ is finite.

Every algebraic number is a root of some integer polynomial, which has some height $N\ge 1$, so it
lies in the corresponding $A_N$; conversely every element of every $A_N$ is algebraic. Writing
$\mathbb{A}$ for the set of algebraic numbers,
\[
\mathbb{A}=\bigcup_{N=1}^{\infty}A_N,
\]
a countable union of finite sets, which is at most countable. Finally $\mathbb{A}$ is infinite, since
each integer $k$ is algebraic as the root of $z-k$, and $\mathbb{Z}$ is countably infinite
(\xrefL{2.1.4}). An infinite, at-most-countable set is countable, so $\mathbb{A}$ is countable.
\end{SolutionBlue}

\begin{Reflection}
The phrase \emph{prove that this set is countable} fixes the strategy before anything else: in our
notes the only practical way to certify countability is to display the set as a countable union of
finite or countable pieces---the quotable \factref{countable-union}{``a countable union of countable
sets is countable''}, supported by \factref{subset-countable}{``a subset of a countable set is finite
or countable''} (\xrefL{2.1.9}). What that pushes the work onto is the \emph{slicing}: into which
finite or countable pieces should the algebraic numbers be cut? Trying to list the roots directly is
hopeless, because there is no natural order on $\mathbb{C}$ and no first algebraic number to start
from.

The hint names the slicing and is worth taking apart, since it teaches the whole technique. An
algebraic number is not a free-standing object; it is manufactured by a polynomial with integer
coefficients, and a polynomial is finite integer data---a degree together with a finite list of
coefficients. Whenever objects are built from finite integer data, the move is to attach to that data
a single integer \emph{size} that is finite-to-one, meaning only finitely many objects share each
value. The height $N=n+\sum_i|a_i|$ is exactly such a size, because fixing $N$ simultaneously caps the
degree ($n\le N$) and every coefficient ($|a_i|\le N$); that simultaneous cap is the reason only
finitely many equations have a given height, and it is the one thing to verify rather than assert.

Believe it on a small case first. Height $N=2$ forces $n+\sum|a_i|=2$ with the coefficients not all
zero; the polynomials are $\pm z$, $\pm 2$ (degree $0$), and a few low-degree pieces, a finite
list one can write out by hand, and each contributes only its handful of roots. Stacking these finite layers $N=1,2,3,\dots$
builds the whole set from the bottom up.

The proof then runs in three short moves, each closed by a named result:
\begin{enumerate}[leftmargin=1.7em,itemsep=3pt,topsep=2pt,label=\textnormal{(\arabic*)}]
  \item each height gives only finitely many polynomials, and a nonzero polynomial of degree $n$ has at
        most $n$ complex roots, so the height-$N$ roots form a finite set $A_N$---this is where the hint
        is spent, and the root bound is the only fact from outside topology that the argument needs;
  \item gathering over all heights writes $\mathbb{A}=\bigcup_{N\ge 1}A_N$, a countable union of finite
        sets, hence at most countable by \factref{countable-union}{the countable-union fact};
  \item ``at most countable'' becomes ``countable'' only after observing the set is infinite, and the
        cheapest witness is that every integer $k$ is the root of $z-k$, so $\mathbb{A}\supseteq\mathbb{Z}$
        is infinite (\xrefL{2.1.4}).
\end{enumerate}

The rigor check is mostly that third step, which is the easy one to drop: a countable union of finite
sets is only guaranteed to be \emph{at most} countable, so without an explicit infinite subset one has
proved less than ``countable.'' The other hidden hypothesis is that the polynomials are taken nonzero
(``not all zero''); the zero polynomial vanishes everywhere and would drag all of $\mathbb{C}$ into a
single layer, destroying the finiteness of $A_N$---which is exactly why the definition excludes it.

The transferable move is the height itself. Anything assembled from finite integer data---polynomials,
finite words over a finite alphabet, finite sequences of integers---can be sorted into finite layers by
a well-chosen integer size with finite fibres, after which the countable-union fact finishes the count.
This same template returns in \hyperlink{prob:3}{Problem~3}, where the countability proved here is
played against the uncountability of $\mathbb{R}$ (\factref{r-uncountable}{``$\mathbb{R}$ is
uncountable''}) to force the existence of transcendental reals.
\end{Reflection}

% ------------------------------------------------------------
% ------------------------------------------------------------
\prob{3}{Rudin, Chapter 2, Exercise 3}
Prove that there exist real numbers which are not algebraic.

\begin{SolutionBlue}
A complex number is \emph{algebraic} when it is a root of some nonzero integer polynomial; write
$\mathbb A$ for the set of all algebraic numbers and $\mathbb A_{\mathbb R}=\mathbb A\cap\mathbb R$ for
the real ones.

The set $\mathbb A$ is countable: this is \hyperlink{prob:F}{Problem~F} (Rudin, Chapter~2,
Exercise~2), which sorts the integer polynomials by the height $N=n+|a_0|+\cdots+|a_n|$, finds only
finitely many of each height and at most $N$ roots apiece, and assembles $\mathbb A$ as a countable
union of finite sets. Hence $\mathbb A_{\mathbb R}$, a subset of $\mathbb A$, is finite or countable
(\factref{subset-countable}{subset of a countable set}; \xrefL{2.1.9}).

Suppose, toward a contradiction, that every real number were algebraic. Then
$\mathbb R=\mathbb A_{\mathbb R}$ would be at most countable. But $\mathbb R$ is uncountable
(\factref{r-uncountable}{$\mathbb R$ is uncountable}; \xrefL{2.1.13}), and an uncountable set cannot
equal an at-most-countable one. This contradiction shows some real number lies outside $\mathbb A$,
i.e.\ there exist real numbers that are not algebraic.
\end{SolutionBlue}

\begin{Reflection}
The statement asks only that a non-algebraic real \emph{exist}, naming no candidate and demanding no
formula---and a pure existence claim with no constructive thread is the classic signal for a counting
argument by contradiction (\xrefL{0.5.5}): rather than exhibit a transcendental number, show that the
algebraic ones are too few to fill $\mathbb R$. The two halves of that strategy are already on the
page. ``Algebraic'' was the subject of the previous problem, where it was shown to be a
\emph{countable} property, and ``real numbers'' carries the one fact about $\mathbb R$ that beats
countability head-on---its uncountability (\factref{r-uncountable}{$\mathbb R$ is uncountable},
\xrefL{2.1.13}). A countable family inside an uncountable set cannot exhaust it, so something must be
left over.

Believe it by sizes before proving it. The algebraic numbers can be listed
$\alpha_1,\alpha_2,\alpha_3,\dots$ (every rational, $\sqrt2$, the golden ratio, and so on), while the
reals provably cannot be put in any such list; a list cannot contain everything off a longer-than-any-
list line, so reals outside the list survive. Those survivors are exactly the transcendental numbers.

The argument is three short moves, each leaning on a named prior result:
\begin{enumerate}[leftmargin=1.7em,itemsep=3pt,topsep=2pt,label=\textnormal{(\arabic*)}]
  \item $\mathbb A$ is countable---this is not reproved here but imported wholesale from
        \hyperlink{prob:F}{Problem~F}, whose height argument writes $\mathbb A$ as a
        \factref{countable-union}{countable union} of finite root-sets;
  \item the real algebraic numbers $\mathbb A_{\mathbb R}=\mathbb A\cap\mathbb R$ form a subset of a
        countable set, hence are at most countable
        (\factref{subset-countable}{subset of a countable set}, \xrefL{2.1.9});
  \item assuming every real were algebraic makes $\mathbb R=\mathbb A_{\mathbb R}$ at most countable,
        which collides with the uncountability of $\mathbb R$ (\xrefL{2.1.13})---the contradiction that
        forces a non-algebraic real.
\end{enumerate}

The reasoning is honest because the two ingredients are genuinely independent and neither is assumed of
the conclusion: countability of $\mathbb A$ comes from organising polynomials, uncountability of
$\mathbb R$ from Cantor's diagonal (\xrefL{2.1.12}), and nothing here presumes a transcendental number
in advance---it is produced, not posited. Every hypothesis earns its place: drop the countability of
$\mathbb A$ and the comparison is empty, drop the uncountability of $\mathbb R$ and a countable line
could be all-algebraic (as the rational line nearly is). Restricting to $\mathbb A_{\mathbb R}$ rather
than all of $\mathbb A$ matters too, since the comparison must happen inside $\mathbb R$.

The move to keep is the cardinality existence proof: to show an object with some property exists,
prove the objects \emph{lacking} it are at most countable and then drop them into an uncountable
ambient set---the overflow is your witness. It is the same accounting that makes $\mathbb Q$ thin in
$\mathbb R$ (\xrefL{2.1.14}); here it conjures transcendentals without ever building one, which is why
exhibiting a specific transcendental (Liouville, $e$, $\pi$) is a strictly harder and separate task.
\end{Reflection}

% ------------------------------------------------------------
% ------------------------------------------------------------
\prob{4}{Rudin, Chapter 2, Exercise 4}
Is the set of all irrational real numbers countable?

\begin{Solution}
No. The set $\mathbb R\setminus\mathbb Q$ of irrational reals is uncountable.

The argument rests on two facts: $\mathbb Q$ is countable (\factref{nzq-countable}{List of Facts};
\xrefL{2.1.7}) and $\mathbb R$ is uncountable (\factref{r-uncountable}{List of Facts};
\xrefL{2.1.13}). Suppose, toward a contradiction, that $\mathbb R\setminus\mathbb Q$ were countable.
Every real is rational or irrational, so
\[
\mathbb R=\mathbb Q\cup(\mathbb R\setminus\mathbb Q)
\]
would be a union of two countable sets, hence countable
(\factref{countable-union}{``a countable union of countable sets is countable''}). That contradicts
the uncountability of $\mathbb R$, so $\mathbb R\setminus\mathbb Q$ is uncountable.

It remains to see why $\mathbb R$ is uncountable, the one input above worth reproving. It is enough to
show the subset $(0,1)$ cannot be listed, since a countable $\mathbb R$ would force the subset
$(0,1)$ to be at most countable (\factref{subset-countable}{``a subset of a countable set is finite
or countable''}). Suppose $(0,1)$ were listed as $x_1,x_2,x_3,\dots$, each written in decimal as
$x_n=0.d_{n1}d_{n2}d_{n3}\cdots$. Build $y=0.e_1e_2e_3\cdots$ along the diagonal by
\[
e_n=\begin{cases} 4 & \text{if } d_{nn}\neq 4,\\[2pt] 7 & \text{if } d_{nn}=4.\end{cases}
\]
Every digit of $y$ is a $4$ or a $7$, so $y$ avoids the ambiguous tails $0.\!\cdots\!999\cdots$ and
$0.\!\cdots\!000\cdots$ and thus has a single decimal expansion. That expansion differs from $x_n$ in
the $n$th digit, so $y\neq x_n$ for every $n$; yet $y\in(0,1)$, so the list omits a point of
$(0,1)$. The list was supposed to be exhaustive, a contradiction. Hence $(0,1)$ is uncountable, and so
is $\mathbb R\supseteq(0,1)$.
\end{Solution}

\begin{Reflection}
The question asks whether the irrationals can be threaded onto a single list $a_1,a_2,a_3,\dots$ the
way $\mathbb Z$ can---whether there are only countably many of them. The word \emph{countable} names
the lecture object (\xrefL{2.1.2}), and the honest first reaction is that listing the irrationals
\emph{directly} is hopeless: they carry no first element and no successor rule, nothing to enumerate
along. That obstruction is itself the clue. When a set resists a head-on count, the move is to place
it inside a set you understand and compare sizes, and the irrationals come pre-packaged as the part of
$\mathbb R$ left over from $\mathbb Q$.

Believe the splitting before proving it. Each real is rational or irrational and never both, so
$\mathbb R=\mathbb Q\cup(\mathbb R\setminus\mathbb Q)$ partitions the line into the listable part and
the part in question. If the leftover were also listable, interleaving the two lists would enumerate
all of $\mathbb R$---and that is exactly what cannot happen.

The proof is a short contradiction (\xrefL{0.5.5}) resting on two quotable facts and one diagonal
core:
\begin{enumerate}[leftmargin=1.7em,itemsep=3pt,topsep=2pt,label=\textnormal{(\arabic*)}]
  \item assume $\mathbb R\setminus\mathbb Q$ is countable; since $\mathbb Q$ is countable as well
        (\factref{nzq-countable}{List of Facts}; \xrefL{2.1.7}), the union
        $\mathbb Q\cup(\mathbb R\setminus\mathbb Q)$ is a union of two countable sets, hence countable
        by \factref{countable-union}{``a countable union of countable sets is countable''};
  \item that union is all of $\mathbb R$, so $\mathbb R$ would be countable, against
        \factref{r-uncountable}{``$\mathbb R$ is uncountable''} (\xrefL{2.1.13})---the collision that
        closes the contradiction;
  \item the only input not to be taken on faith is the uncountability of $\mathbb R$, secured by
        Cantor's diagonal argument on $(0,1)$ (\xrefL{2.1.12}), the decimal cousin of the binary
        statement \factref{binary-uncountable}{``the $0$--$1$ sequences are uncountable''}
        (\xrefL{2.1.10}).
\end{enumerate}

The diagonal step is the one to actually understand. Given any candidate list of $(0,1)$ written as
decimals, read down the diagonal and change each digit; the number you build disagrees with the $n$th
entry in the $n$th place, so it lies in $(0,1)$ yet sits on no row---the list was incomplete. The lone
subtlety is that a few reals wear two decimal coats ($0.4999\cdots=0.5000\cdots$), which would let a
``new'' number secretly equal a listed one; restricting the manufactured digits to $4$ and $7$ dodges
both ambiguous tails, so the constructed $y$ has a unique expansion and the disagreement is real. Drop
that precaution and the argument genuinely has a gap, which is why it is built in rather than waved
past.

The rigor turns on the asymmetry between the two facts: countability is preserved by the union of
\emph{two} pieces, but only because finitely (indeed countably) many countable sets stay countable;
were $\mathbb R\setminus\mathbb Q$ countable, $\mathbb R$ would inherit that, and the diagonal forbids
it. Nothing is circular---the uncountability of $\mathbb R$ is proved independently of any claim about
the irrationals.

The transferable move is the whole point: to size a set that defies direct enumeration, write the
ambient set as that set together with a countable remainder, and let a known uncountability do the
work. The same accounting shows the transcendentals are uncountable (the algebraic numbers being only
countable, as in \hyperlink{prob:F}{Problem~F}); a countable exception cannot dent an uncountable
whole, so the strange numbers are the overwhelming majority, even when not one is easy to name.
\end{Reflection}

% ------------------------------------------------------------
% ------------------------------------------------------------
\prob{5}{Rudin, Chapter 2, Exercise 5}
Construct a bounded set of real numbers with exactly three limit points.

\begin{Solution}
Glue together three shifted copies of the sequence $\tfrac1n$, aimed at the separated targets $0$, $2$,
$4$. With $n$ ranging over the positive integers, set
\[
E=\Bigl\{\tfrac1n\Bigr\}\cup\Bigl\{2+\tfrac1n\Bigr\}\cup\Bigl\{4+\tfrac1n\Bigr\}.
\]
The verdict: $E$ is bounded, and its set of cluster points is exactly $\{0,2,4\}$.

The set is bounded. From $0<\tfrac1n\le 1$, $\;2<2+\tfrac1n\le 3$, and $4<4+\tfrac1n\le 5$, every element
lies in $[0,5]$, so $E\subseteq[0,5]$ is bounded (\xrefL{4.3.4}).

Each of $0$, $2$, $4$ is a cluster point. Fix $c\in\{0,2,4\}$ and any radius $r>0$. By the Archimedean
property (\factref{archimedean}{Archimedean}) there is an $n$ with $\tfrac1n<r$; the point $c+\tfrac1n$ lies
in $E$, differs from $c$, and satisfies $|(c+\tfrac1n)-c|=\tfrac1n<r$. Every ball about $c$ thus contains a
point of $E$ other than $c$, which is the definition of a cluster point (\xrefL{3.2.1}).

No other point is a cluster point. Take $p\notin\{0,2,4\}$ and set
\[
\delta=\min\{\,|p|,\,|p-2|,\,|p-4|\,\}>0,
\]
the distance from $p$ to the nearest target. Each of the three families converges to its target, so for all
but finitely many $n$ the points $\tfrac1n$, $2+\tfrac1n$, $4+\tfrac1n$ lie within $\delta/2$ of $0$, $2$,
$4$ respectively, hence outside the ball $B(p,\delta/2)$. That ball therefore meets $E$ in only finitely
many points, and a cluster point would force infinitely many points of $E$ into every one of its balls
(\factref{cluster-infinite}{a cluster point is approached infinitely often}). So $p$ is not a cluster point.

Thus $E$ is bounded with cluster-point set precisely $\{0,2,4\}$.
\end{Solution}

\begin{Reflection}
The phrase that fixes the strategy is \emph{limit point}---a cluster point, a place a set piles up so that
every ball around it catches another member of the set (\xrefL{3.2.1}). The request is constructive: build
one set that stays inside a bounded interval, accumulates at three spots, and accumulates nowhere else. So
the question is really ``what is the cheapest way to manufacture a single accumulation point, and how do I
place three of them without their debris colliding?''

The cheapest accumulation point is a convergent sequence that never reaches its limit. The numbers $\tfrac1n$
crowd toward $0$ without ever hitting it, so $0$ is a cluster point of $\{\tfrac1n\}$ and---because the terms
march away to the right---its only one (\xrefL{3.2.5}). Believe that picture first: a string of dots
bunching up against a wall they never touch. To get three cluster points, take three such strings and aim
them at $0$, $2$, $4$; keeping the targets two apart, and the whole set inside $[0,5]$, both bounds the set
and stops the three families from interfering with one another.

The two halves of the proof are unequal in difficulty, and reading which is which is the lesson:
\begin{enumerate}[leftmargin=1.7em,itemsep=3pt,topsep=2pt,label=\textnormal{(\arabic*)}]
  \item that $0$, $2$, $4$ \emph{are} cluster points is the easy half---it falls straight out of the
        convergences, with the Archimedean property (\factref{archimedean}{Archimedean}) supplying, for any
        radius $r$, an index $n$ with $\tfrac1n<r$, so the point $c+\tfrac1n\in E$ sits in the ball about $c$
        and differs from $c$;
  \item that there are \emph{no others} is the real work, and the right tool is that a cluster point pulls
        infinitely many points of the set into every ball around it (\factref{cluster-infinite}{a cluster
        point is approached infinitely often}); so to acquit a candidate $p\notin\{0,2,4\}$ it is enough to
        trap it in one ball holding only finitely many points of $E$, which the separation gap
        $\delta=\min\{|p|,|p-2|,|p-4|\}$ provides, since all but finitely many points of each family have
        already gathered next to their target.
\end{enumerate}

The rigor turns on that gap being strictly positive, which is exactly why the targets had to be kept apart:
$\delta>0$ holds precisely because $p$ avoids all three targets, and it is what guarantees only a finite
remnant of each family can intrude. Boundedness is independent and never in doubt, since $E\subseteq[0,5]$
outright (\xrefL{4.3.4}). The two directions together---``these three are'' and ``nothing else is''---are
both needed; stopping after the first would only show $\{0,2,4\}$ is contained in the cluster set, not equal
to it.

The move generalises into a template. For a bounded set with exactly $k$ cluster points, union $k$ copies of
a sequence converging to $0$, each shifted onto a distinct, well-separated target; the targets become the
cluster points, and the infinitely-many-points criterion (\factref{cluster-infinite}{a cluster point is
approached infinitely often}) keeps any stray accumulation from appearing. The same idea, pushed to countably many targets that themselves accumulate, builds far richer
cluster sets---it is the engine behind \hyperlink{prob:13}{Problem~13}.
\end{Reflection}

% ------------------------------------------------------------
% ------------------------------------------------------------
\prob{6}{Rudin, Chapter 2, Exercise 6}
Let $E'$ be the set of all limit points of a set $E$. Prove that $E'$ is closed. Prove that $E$ and
$\overline{E}$ have the same limit points. (Recall that $\overline{E}=E\cup E'$.) Do $E$ and $E'$
always have the same limit points?

\begin{SolutionBlue}
Work in a metric space $(X,d)$, and write $B_r(p)=\{x:d(p,x)<r\}$ for the open ball. A point $p$ is a
limit point of a set $S$ when every $B_r(p)$ contains a point of $S$ other than $p$ (\xrefL{3.2.1});
$S$ is closed when it contains all of its limit points (\xrefL{3.2.4}).

The set $E'$ is closed. Let $p$ be a limit point of $E'$ and fix any radius $r>0$; the goal is to
produce a point of $E$ inside $B_r(p)$ different from $p$. Since $p$ is a limit point of $E'$, there is
a point $q\in E'$ with $q\neq p$ and $d(p,q)<r$. Put $s=\min\{\,r-d(p,q),\,d(p,q)\,\}>0$. Since $q$ is
itself a limit point of $E$, there is a point $x\in E$ with $x\neq q$ and $d(q,x)<s$. Two estimates now
finish it. By the triangle inequality
\[
  d(p,x)\le d(p,q)+d(q,x)<d(p,q)+\bigl(r-d(p,q)\bigr)=r,
\]
so $x\in B_r(p)$; and again by the triangle inequality, in the reverse form,
\[
  d(p,x)\ge d(p,q)-d(q,x)>d(p,q)-d(p,q)=0,
\]
so $x\neq p$. Every ball about $p$ thus contains a point of $E$ other than $p$, so $p\in E'$. Hence
$E'$ contains all its limit points and is closed.

The sets $E$ and $\overline{E}$ have the same limit points. Since $E\subseteq\overline{E}$, any ball
that meets $E$ in a point other than $p$ also meets $\overline{E}$ there, so every limit point of $E$
is a limit point of $\overline{E}$. For the reverse, let $p$ be a limit point of $\overline{E}$ and fix
$r>0$. Choose $q\in\overline{E}$ with $q\neq p$ and $d(p,q)<r$. By $\overline{E}=E\cup E'$ there are two
cases. If $q\in E$, then $B_r(p)$ already contains the point $q\in E$ with $q\neq p$. If instead $q\in
E'$, then $q$ is a limit point of $E$, and the same two-step estimate as above---taking $x\in E$ with
$x\neq q$ and $d(q,x)<\min\{\,r-d(p,q),\,d(p,q)\,\}$---yields a point $x\in E$ with $x\neq p$ and
$d(p,x)<r$. In either case $B_r(p)$ contains a point of $E$ other than $p$, so $p$ is a limit point of
$E$. Therefore $E$ and $\overline{E}$ have exactly the same limit points.

The sets $E$ and $E'$ need not have the same limit points. In $\mathbb{R}$ take
\[
  E=\Bigl\{\tfrac1n:n\in\mathbb{Z}^{+}\Bigr\}.
\]
Its only limit point is $0$, so $E'=\{0\}$. Then $0$ is a limit point of $E$ but not of $E'$: the
one-point set $E'=\{0\}$ has no limit point at all, since the ball $B_1(0)$ contains no point of $E'$
different from $0$. So $E$ has a limit point that $E'$ does not, and the two sets do not always share
their limit points.
\end{SolutionBlue}

\begin{Reflection}
The whole problem is built from one definition---a limit point of $S$ is a point every ball around
which catches a member of $S$ other than itself (\xrefL{3.2.1})---and the recurring move is to chain
two of these statements together. ``Prove $E'$ is closed'' is a request to show $E'$ contains its own
limit points (\xrefL{3.2.4}), and the only data you start with is that $p$ is close to points $q$ that
are in turn close to points of $E$; passing the closeness along, in two hops $p\rightsquigarrow
q\rightsquigarrow x$, is the entire idea. The hard part is purely bookkeeping on the radii, and the
quotable fact that a limit point attracts \emph{infinitely many} members
(\factref{cluster-infinite}{cluster points pull in infinitely many points}) is what guarantees each hop
has somewhere to land even after you shrink the tolerance.

A picture makes the radius juggling believable before the inequalities are written. Imagine $p$ with a
ball of radius $r$; pick $q\in E'$ strictly inside, at distance $d(p,q)<r$. Two things could spoil the
landing point $x\in E$ near $q$: it could be pushed out past radius $r$, or it could land exactly on
$p$. Shrinking the search radius around $q$ to $s=\min\{r-d(p,q),\,d(p,q)\}$ closes off both:
$r-d(p,q)$ is the slack left before the outer wall, and $d(p,q)$ is the room between $q$ and $p$, so an
$x$ within $s$ of $q$ stays inside the $r$-ball yet cannot reach $p$.

The proof then runs in clean steps, each resting on the one before:
\begin{enumerate}[leftmargin=1.7em,itemsep=3pt,topsep=2pt,label=\textnormal{(\arabic*)}]
  \item from $p\in(E')'$ and the definition of a limit point (\xrefL{3.2.1}), extract $q\in E'$ with
        $q\neq p$ and $d(p,q)<r$---the first hop;
  \item set the guarded radius $s=\min\{r-d(p,q),\,d(p,q)\}$, positive because $d(p,q)$ lies strictly
        between $0$ and $r$, and use $q\in E'$ to extract $x\in E$ with $x\neq q$ and $d(q,x)<s$---the
        second hop;
  \item the forward triangle inequality $d(p,x)\le d(p,q)+d(q,x)<r$ keeps $x$ inside $B_r(p)$, and the
        reverse form $d(p,x)\ge d(p,q)-d(q,x)>0$ keeps $x$ off $p$, so $B_r(p)$ holds a point of $E$
        unequal to $p$, which is exactly $p\in E'$ (\xrefL{3.2.4}).
\end{enumerate}

The equality of limit points for $E$ and $\overline E$ reuses the very same machine. One inclusion is
free from monotonicity---a bigger set has at least as many limit points---and the other splits on
$\overline E=E\cup E'$ (\xrefL{4.4.1}): a nearby point of $\overline E$ is either already in $E$, and
then there is nothing to do, or it is a limit point of $E$, and then the two-hop estimate of step (3)
manufactures the needed point of $E$. This is also why $E'$ being closed is not a throwaway first
part---it is the engine that makes $\overline E=E\cup E'$ closed at all
(\factref{closure-closed}{$E\cup E'$ is closed}), since adjoining a closed set of limit points to $E$
cannot create new ones to escape through.

Two checks keep the argument honest. The guard $s$ must be \emph{strictly} positive, which is why both
inequalities $0<d(p,q)<r$ are needed---the first from $q\neq p$, the second from $d(p,q)<r$---and
dropping either one lets $x$ collapse onto $p$ or spill outside the ball. And the case split on
$\overline E=E\cup E'$ is genuinely two cases: skipping the easy $q\in E$ branch would leave a gap, not
a shortcut. There is no circularity, since closedness of $E'$ is established first and only then fed
into the closure statement.

The final question is the one that separates the pattern from its limits. Closing a set never adds
limit points, but \emph{replacing} a set by its limit points can lose them: $E=\{1/n\}$ piles up at
$0$, so $0\in E'$, yet $E'=\{0\}$ is a lone point with no pile-up of its own. The takeaway is that
$E\mapsto\overline E$ is idempotent---$E$ and $\overline E$ carry exactly the same limit points, so
$\overline{\overline E}=\overline E$---while $E\mapsto E'$ is a genuine collapse that can strip a set
down to isolated points: here it sends $\{1/n\}$ to $\{0\}$ and then to $\varnothing$, so $E''=\varnothing
\neq E'$. Passing to $\overline E$ is the safe closure operation and passing to $E'$ is not.
\end{Reflection}

% ------------------------------------------------------------
% ------------------------------------------------------------
\prob{7}{Rudin, Chapter 2, Exercise 7}
Let $A_{1}, A_{2}, A_{3}, \ldots$ be subsets of a metric space.
\begin{enumerate}[leftmargin=1.7em,itemsep=2pt,topsep=2pt,label=\textnormal{(\alph*)}]
  \item If $B_{n} = \bigcup_{i=1}^{n} A_{i}$, prove that $\overline{B}_{n} = \bigcup_{i=1}^{n} \overline{A}_{i}$, for $n = 1, 2, 3, \ldots$.
  \item If $B = \bigcup_{i=1}^{\infty} A_{i}$, prove that $\overline{B} \supset \bigcup_{i=1}^{\infty} \overline{A}_{i}$.
\end{enumerate}
Show, by an example, that this inclusion can be proper.

\begin{Solution}
Throughout, the closure $\overline{E}=E\cup E'$ is the smallest closed set containing $E$, and it is
monotone: if $A\subseteq B$, then $\overline A$ is contained in the closed set $\overline B$, so
$\overline A\subseteq\overline B$.

For (a), fix $n$. Each $A_{i}\subseteq B_{n}$ for $i\le n$, so monotonicity gives
$\overline{A}_{i}\subseteq\overline{B}_{n}$, and hence
$\bigcup_{i=1}^{n}\overline{A}_{i}\subseteq\overline{B}_{n}$. For the reverse inclusion, each
$\overline{A}_{i}$ is closed (\xrefL{4.4.1}), so $\bigcup_{i=1}^{n}\overline{A}_{i}$ is a \emph{finite}
union of closed sets and is therefore closed (\factref{inter-closed}{finite unions of closed sets are
closed}; \xrefL{3.2.10}). It contains $B_{n}=\bigcup_{i=1}^{n}A_{i}$, since $A_{i}\subseteq\overline{A}_{i}$
for each $i$. As $\overline{B}_{n}$ is the smallest closed set containing $B_{n}$, it follows that
$\overline{B}_{n}\subseteq\bigcup_{i=1}^{n}\overline{A}_{i}$. The two inclusions give
\[
\overline{B}_{n}=\bigcup_{i=1}^{n}\overline{A}_{i}.
\]

For (b), every $A_{i}\subseteq B$, so monotonicity gives $\overline{A}_{i}\subseteq\overline{B}$ for each
$i$, and therefore
\[
\bigcup_{i=1}^{\infty}\overline{A}_{i}\subseteq\overline{B}.
\]

The inclusion in (b) can be proper. Take $A_{i}=\{1/i\}\subseteq\mathbb{R}$ for $i=1,2,3,\dots$. Each
singleton is closed (a finite set has no cluster points, so it contains them all vacuously, \xrefL{3.2.4}),
so $\overline{A}_{i}=A_{i}$ and
\[
\bigcup_{i=1}^{\infty}\overline{A}_{i}=\Bigl\{\tfrac{1}{i}:i\ge 1\Bigr\}.
\]
The point $0$ is a cluster point of $B=\bigcup_{i=1}^{\infty}A_{i}=\{1/i:i\ge 1\}$: given $\varepsilon>0$,
the Archimedean property (\factref{archimedean}{Archimedean property}) supplies an integer $i>1/\varepsilon$,
and then $0<1/i<\varepsilon$, so the ball $B_{\varepsilon}(0)$ contains the point $1/i\neq 0$ of $B$. Hence
$0\in\overline{B}$, while $0\notin\bigcup_{i=1}^{\infty}\overline{A}_{i}$, and the inclusion is strict.
\end{Solution}

\begin{Reflection}
The question is whether one of the two operations in sight commutes with the other: does closing a union
equal the union of the closures? The word \emph{closure} points at the one structural fact that drives the
whole problem, that $\overline{E}=E\cup E'$ is the \emph{smallest} closed set containing $E$ (\xrefL{4.4.1}).
Once closure is read that way---as a minimal closed hull rather than ``$E$ with its limit points glued
on''---the easy half falls out of monotonicity and the hard half turns into a question about whether a union
of closed sets stays closed, which is exactly where ``finite'' versus ``infinite'' will decide the answer.

Monotonicity is worth isolating because it is used in every part and is itself just the minimality of the
hull: if $A\subseteq B$, then $\overline B$ is a closed set that already contains $A$, so the smallest such
set $\overline A$ cannot be larger. That single observation gives $\bigcup\overline{A}_{i}\subseteq
\overline{\bigcup A_{i}}$ for free, in both the finite and the infinite case---each piece $A_{i}$ sits inside
the union, so its hull sits inside the union's hull. The content of the problem is never this direction; it
is the reverse one in (a), and its failure in (b).

The reverse inclusion in (a) rests on one load-bearing step, which is where you should slow down:
\begin{enumerate}[leftmargin=1.7em,itemsep=3pt,topsep=2pt,label=\textnormal{(\arabic*)}]
  \item each $\overline{A}_{i}$ is closed (\xrefL{4.4.1}), and a \emph{finite} union of closed sets is again
        closed (\factref{inter-closed}{finite unions of closed sets are closed}; \xrefL{3.2.10}), so
        $\bigcup_{i=1}^{n}\overline{A}_{i}$ is itself a closed set;
  \item this closed set contains $B_{n}$, because $A_{i}\subseteq\overline{A}_{i}$ pushes the whole union
        $B_{n}=\bigcup_{i\le n}A_{i}$ inside it;
  \item $\overline{B}_{n}$ is by definition the \emph{smallest} closed set containing $B_{n}$ (\xrefL{4.4.1}),
        so it can only be smaller than the closed set produced in (1), giving
        $\overline{B}_{n}\subseteq\bigcup_{i=1}^{n}\overline{A}_{i}$ and, with monotonicity, equality.
\end{enumerate}
The entire argument is the minimality of the hull played against a closed set you have built by hand; nothing
is circular, since the built set is closed for an independent reason.

Step (1) is exactly where the infinite case breaks, and seeing why is the point of the counterexample. An
\emph{infinite} union of closed sets need not be closed---the same phenomenon as an infinite intersection of
open sets failing to be open (\xrefL{3.1.3})---so for the infinite union there is no closed set built from the
$\overline{A}_{i}$ to play against, and $\overline{B}$ is free to contain points lying in no single
$\overline{A}_{i}$. The choice $A_{i}=\{1/i\}$ is the cleanest witness: each set is a lone closed point, equal
to its own closure, yet the points $1/i$ together pile up at $0$. That $0$ is a cluster point of $B$ is the
Archimedean property in disguise (\factref{archimedean}{$\inf\{1/n\}=0$})---every ball about $0$ swallows a
tail of the $1/i$ (\factref{cluster-infinite}{a cluster point is approached by infinitely many points of the
set})---so $0\in\overline{B}$, while $0$ belongs to none of the $\overline{A}_{i}=\{1/i\}$.

The transferable reading is that closure commutes with \emph{finite} unions and only finite ones, and the gap
is always carried by cluster points that no single piece $A_{i}$ can reach---points produced only by drawing
terms from infinitely many of the sets at once. Whenever a topological identity holds for finite unions but
you suspect it fails in the limit, look for the operation's stability theorem (here, finite unions of closed
sets are closed, \factref{inter-closed}{finite unions of closed sets are closed}) and probe exactly the
hypothesis ``finite'': the counterexample lives at the limit point the finite version was quietly forbidding.
\end{Reflection}

% ------------------------------------------------------------
% ------------------------------------------------------------
\prob{8}{Rudin, Chapter 2, Exercise 8}
Is every point of every open set $E\subset\mathbb R^{2}$ a limit point of $E$? Answer the same
question for closed sets in $\mathbb R^{2}$.

\begin{Solution}
Yes for open sets; no for closed sets.

Let $E\subseteq\mathbb R^{2}$ be open and fix $p\in E$. Openness makes $p$ an interior point, so some
ball $\Beps(p)\subseteq E$ with $\eps>0$ (\xrefL{2.3.3}). To see $p$ is a limit point, take any radius
$s>0$, set $\rho=\min(\eps,s)>0$, and let $q=p+\bigl(\tfrac{\rho}{2},0\bigr)$. Then
$|q-p|=\tfrac{\rho}{2}<\rho\le\eps$, so $q\in\Beps(p)\subseteq E$; and $\tfrac{\rho}{2}<\rho\le s$, so
$q\in B_{s}(p)$; and $q\neq p$. Every ball about $p$ thus contains a point of $E$ other than $p$, which
is exactly the definition of a limit point (\xrefL{3.2.1}). Hence every point of an open set in
$\mathbb R^{2}$ is a limit point of it.

For closed sets the answer is no. Fix any $p\in\mathbb R^{2}$ and take $E=\{p\}$. This $E$ is closed:
being finite it has no limit points---a limit point would force infinitely many points of $E$ into
each of its neighbourhoods (\factref{cluster-infinite}{a cluster point captures infinitely many points
of the set}, \xrefL{3.2.1})---so $E$ vacuously contains all of them and is closed (\xrefL{3.2.4}). Yet
$p$ is \emph{not} a limit point of $E$: the ball $B_{1}(p)$ holds no point of $E$ different from $p$,
so $p$ is an isolated point of $E$ (\xrefL{3.2.2}). Thus $p$ is a point of a closed set that fails to
be a limit point of it.
\end{Solution}

\begin{Reflection}
The two answers split because ``open'' and ``closed'' constrain genuinely different things, and the way
to read the problem is to translate each into the language of limit points. A limit point is a
\emph{crowded} point: every ball around it, no matter how small, still catches another point of the set
(\xrefL{3.2.1}). So the real question is whether merely belonging to the set forces that crowding---and
the equivalent notion, an \emph{isolated} point (\xrefL{3.2.2}), a member with a ball around it meeting
the set in itself alone, is precisely the witness that crowding can fail. The contrast the problem is
fishing for is whether open or closed sets are allowed isolated points.

Believe the open case through the picture first: a point of an open set sits at the centre of a whole
disk lying inside the set (that is what \xrefL{2.3.3} means), and in the plane no disk is a single
point---slide a hair in any direction and you are still inside it. So neighbours of $p$ inside the set
are unavoidable. The closed case is believed through the cleanest counterexample, a lone point
$\{p\}$: there is nothing nearby to crowd it, and closedness has no complaint, since a finite set has
no limit points to omit.

The open half runs in three short moves:
\begin{enumerate}[leftmargin=1.7em,itemsep=3pt,topsep=2pt,label=\textnormal{(\arabic*)}]
  \item openness gives a ball $\Beps(p)\subseteq E$, because each point of an open set is interior to
        it (\xrefL{2.3.3})---this is the only hypothesis used, and it is the whole supply of nearby
        points;
  \item to test a limit point you must beat an \emph{arbitrary} radius $s$, so you shrink to
        $\rho=\min(\eps,s)$ to stay inside both the witnessing ball and the test ball at once;
  \item the explicit neighbour $q=p+(\tfrac{\rho}{2},0)$ lands in $E$ and in $B_{s}(p)$ and differs
        from $p$, meeting the definition (\xrefL{3.2.1}) on the nose.
\end{enumerate}
Move (2) is where rigour lives: a limit point must be checked against \emph{every} radius, not one
convenient one, and taking the minimum is the standard way to honour that ``for all $s$'' without case
work.

The closed half is a single counterexample, and its honesty rests on why $\{p\}$ is closed: a finite
set has no limit points at all, since a limit point would drag infinitely many members into every
neighbourhood (\factref{cluster-infinite}{the infinitely-many-points criterion}, \xrefL{3.2.1}), so the
``contains all its limit points'' condition (\xrefL{3.2.4}) holds vacuously. One example is enough to
answer a ``does this always hold?'' question in the negative, and one is all that is needed.

Each verdict turns on isolated points, and that is the transferable line to keep (\xrefL{3.2.7}): an
open set in $\mathbb R^{k}$ has none---nothing about $\mathbb R^{2}$ was special, every $\mathbb R^{k}$
behaves the same---so every one of its points is a limit point; a closed set may consist entirely of
isolated points. When a problem wants a point of a set that is \emph{not} a limit point, reach for a
closed set with isolated points: a finite set, or a discrete lattice like $\mathbb Z^{2}$.
\end{Reflection}

% ------------------------------------------------------------
% ------------------------------------------------------------
\prob{9}{Rudin, Chapter 2, Exercise 9}
Let $E^{\circ}$ denote the set of all interior points of a set $E$. ($E^{\circ}$ is called the
\emph{interior} of $E$.)
\begin{enumerate}[leftmargin=1.7em,itemsep=2pt,topsep=2pt,label=\textnormal{(\alph*)}]
  \item Prove that $E^{\circ}$ is always open.
  \item Prove that $E$ is open if and only if $E^{\circ}=E$.
  \item If $G\subset E$ and $G$ is open, prove that $G\subset E^{\circ}$.
  \item Prove that the complement of $E^{\circ}$ is the closure of the complement of $E$.
  \item Do $E$ and $\overline{E}$ always have the same interiors?
  \item Do $E$ and $E^{\circ}$ always have the same closures?
\end{enumerate}

\begin{Solution}
Throughout, $E$ is a subset of a metric space $X$ and every complement is taken in $X$. A point $p$ is
\emph{interior} to $E$ when some ball $N_{r}(p)\subseteq E$, and $E^{\circ}$ is the set of such points
(\xrefL{2.3.3}); by definition $E^{\circ}\subseteq E$.

\smallskip\noindent\textit{(a)}\enspace Let $p\in E^{\circ}$, so $N_{r}(p)\subseteq E$ for some
$r>0$. The ball $N_{r}(p)$ is itself open (\xrefL{2.3.5}; \factref{balls-open}{the open balls are
open}), so each $q\in N_{r}(p)$ has a ball $N_{s}(q)\subseteq N_{r}(p)\subseteq E$, making $q$ an
interior point of $E$. Thus $N_{r}(p)\subseteq E^{\circ}$, so $p$ is interior to $E^{\circ}$. Every
point of $E^{\circ}$ is interior to it, hence $E^{\circ}$ is open.

\smallskip\noindent\textit{(b)}\enspace If $E^{\circ}=E$, then $E$ is open by (a). Conversely, if $E$
is open, every $p\in E$ has a ball $N_{r}(p)\subseteq E$, so $p\in E^{\circ}$; with the standing
inclusion $E^{\circ}\subseteq E$ this gives $E^{\circ}=E$.

\smallskip\noindent\textit{(c)}\enspace Let $p\in G$. Since $G$ is open, some ball $N_{r}(p)\subseteq
G$, and $G\subseteq E$ gives $N_{r}(p)\subseteq E$; thus $p\in E^{\circ}$. Hence $G\subseteq
E^{\circ}$.

\smallskip\noindent\textit{(d)}\enspace We prove $(E^{\circ})^{c}=\overline{E^{c}}$ by showing the two
sides have the same points. A point $p$ fails to lie in $E^{\circ}$ exactly when \emph{no} ball about
$p$ is contained in $E$, that is, when \emph{every} ball about $p$ contains a point of $E^{c}$. On the
other side, $\overline{E^{c}}=E^{c}\cup(E^{c})'$ (\xrefL{4.4.1}), so $p\in\overline{E^{c}}$ means $p\in
E^{c}$ or $p$ is a cluster point of $E^{c}$ (\xrefL{3.2.1})---and either way every ball about $p$ meets
$E^{c}$ (if $p\in E^{c}$, the point $p$ itself is in the ball). Conversely, if every ball about $p$
meets $E^{c}$: should $p\notin E^{c}$, each such meeting point differs from $p$, making $p$ a cluster
point of $E^{c}$, so $p\in\overline{E^{c}}$; and if $p\in E^{c}$ then $p\in\overline{E^{c}}$ outright.
The two conditions coincide, so $(E^{\circ})^{c}=\overline{E^{c}}$.

\smallskip\noindent\textit{(e)}\enspace No. In $\mathbb R$ take $E=(-1,0)\cup(0,1)$. Then $E$ is open,
so $E^{\circ}=E$ by (b); but $\overline{E}=[-1,1]$, whose interior is $(-1,1)$. Since $0\in(-1,1)\setminus
E$, the interiors of $E$ and $\overline{E}$ differ.

\smallskip\noindent\textit{(f)}\enspace No. In $\mathbb R$ take $E=\{0\}$. A one-point set has no
cluster point, so it is closed (\xrefL{3.2.4}) and $\overline{E}=\{0\}$; but $0$ is not interior to
$E$, so $E^{\circ}=\varnothing$ and $\overline{E^{\circ}}=\varnothing$. The closures of $E$ and
$E^{\circ}$ differ.
\end{Solution}

\begin{Reflection}
The single word that organises the whole problem is \emph{interior}. An interior point is one with
room to spare---a ball entirely inside $E$ (\xrefL{2.3.3})---and that ``ball inside'' clause is the only
machinery the first three parts ever touch. Recognising this is what tells you that (a), (b), (c) are
not three separate tasks but three readings of one definition: (a) says $E^{\circ}$ is open, (c) says it
swallows every open subset of $E$, and together they identify $E^{\circ}$ as the \emph{largest} open
set inside $E$, the inside-out twin of the closure (\xrefL{4.4.1}), which is the smallest closed set
containing $E$. Part (b) is then the fixed-point statement: a set is open exactly when it already equals
its own core.

Believe that picture on an example before proving it. For $E=[0,1]$ the interior is the open core
$(0,1)$ and the closure is $[0,1]$; the endpoints are interior to neither but cluster points of both.
Interior carves \emph{off} the boundary, closure glues it \emph{on}, and part (d) is the precise
statement that these two operations are mirror images through complementation.

Each step rests on something named:
\begin{enumerate}[leftmargin=1.7em,itemsep=3pt,topsep=2pt,label=\textnormal{(\arabic*)}]
  \item for (a), the leap from ``$p$ has a ball in $E$'' to ``every \emph{nearby} point does too'' is
        powered by the fact that a ball is open (\xrefL{2.3.5}; \factref{balls-open}{open balls are
        open}): inside $N_{r}(p)$ each $q$ gets its own sub-ball, so $N_{r}(p)\subseteq E^{\circ}$ and
        openness propagates from the one point to the whole ball;
  \item (b) is a biconditional, so it is two implications (\xrefL{0.5.9})---one direction is just (a),
        the other reads ``$E$ open'' through the definition of interior, and the standing inclusion
        $E^{\circ}\subseteq E$ supplies the half you never have to argue;
  \item (c) chains two containments, $N_{r}(p)\subseteq G$ from openness of $G$ and $G\subseteq E$ from
        hypothesis, so the witnessing ball lands in $E$ and certifies $p\in E^{\circ}$ (\xrefL{0.3.6});
  \item (d) is the crux, an equality of sets proved by matching membership conditions (\xrefL{0.5.10}):
        negating ``some ball lies in $E$'' yields ``every ball meets $E^{c}$,'' and that is exactly the
        neighbourhood test for the closure $E^{c}\cup(E^{c})'$ (\xrefL{4.4.1}, \xrefL{3.2.1}). The lone
        subtlety is the point $p$ itself---when $p\in E^{c}$ a ball meets $E^{c}$ at $p$ without $p$
        being a cluster point, which is why the closure is $E^{c}\cup(E^{c})'$ and not $(E^{c})'$ alone.
\end{enumerate}

The rigour check is to confirm each direction of (d) is reversible and that the two counterexamples
genuinely separate the operations. They do, in opposite directions: closing can \emph{grow} the
interior, as $(-1,0)\cup(0,1)$ (its own interior) closes to $[-1,1]$ whose interior $(-1,1)$ has
swallowed the missing $0$; while taking the interior can \emph{erase} a set, as $\{0\}$ (closed, equal
to its closure) has empty interior whose closure is empty. So $E$, $E^{\circ}$, and $\overline{E}$ are
three honestly different sets, and neither rounding operation is stable under the other.

The move to carry away is duality by complement: part (d) says $(E^{\circ})^{c}=\overline{E^{c}}$, so
every theorem about interiors converts---swap $E\leftrightarrow E^{c}$ and interior$\leftrightarrow$
closure---into a theorem about closures, and you prove only one of each pair. It is the same
bookkeeping that turns ``open iff complement closed'' (\factref{open-iff-closed}{open $\iff$ complement
closed}, \xrefL{3.2.8}) into a two-for-one, and the reason interior and closure always travel together.
\end{Reflection}

% ------------------------------------------------------------
% ------------------------------------------------------------
\prob{10}{Rudin, Chapter 2, Exercise 10}
Let $X$ be an infinite set. For $p\in X$ and $q\in X$, define
\[
d(p,q)=
\begin{cases}
1 & (\text{if } p\neq q)\\
0 & (\text{if } p=q).
\end{cases}
\]
Prove that this is a metric. Which subsets of the resulting metric space are open? Which are closed?
Which are compact?

\begin{Solution}
The map $d$ is a metric. Every subset of $X$ is both open and closed. The compact subsets are exactly
the finite ones.

First we check that $d$ is a metric. From the definition $d(p,q)\ge 0$, with $d(p,q)=0$ precisely when
$p=q$, and $d(p,q)=d(q,p)$ since ``$p\neq q$'' is symmetric in $p$ and $q$. For the triangle inequality, fix
$p,q,r\in X$. If $p=q$ the left side $d(p,q)=0\le d(p,r)+d(r,q)$. If $p\neq q$, then $r$ cannot equal
both $p$ and $q$, so at least one of $d(p,r),d(r,q)$ equals $1$, giving $d(p,q)=1\le d(p,r)+d(r,q)$.
Hence all four axioms hold and $d$ is a metric (\xrefL{2.2.1}); this is the discrete metric of
\xrefL{2.2.4}.

Everything topological is read off the open balls, so we compute them next. For any $p\in X$ and radius
$0<\varepsilon\le 1$,
\[
B_\varepsilon(p)=\{q\in X:d(p,q)<\varepsilon\}=\{p\},
\]
since the only point within distance $1$ of $p$ is $p$ itself; for $\varepsilon>1$ the ball is all of
$X$. Everything below rests on the first case.

Every subset is open. Let $E\subseteq X$ and $p\in E$. The ball $B_{1/2}(p)=\{p\}\subseteq E$
witnesses $p$ as an interior point, so $E$ is open (\xrefL{2.3.3}). (Equivalently, each singleton
$\{p\}=B_{1/2}(p)$ is open by \factref{balls-open}{``open balls are open''}, and $E=\bigcup_{p\in E}\{p\}$
is a union of open sets, hence open by \factref{union-open}{the union rule}.)

Every subset is closed as well. Each subset $E$ has open complement $X\setminus E$, so $E$ is closed
(\xrefL{3.2.8}; \factref{open-iff-closed}{``open $\iff$ complement closed''}). Every subset is therefore
clopen.

Finally, the compact subsets are exactly the finite ones. A finite set $\{x_1,\dots,x_m\}$ is compact:
given any open cover, choose for each $x_j$ one cover member containing it, and these $m$ members form a
finite subcover (\xrefL{3.4.1}). Conversely, suppose $E\subseteq X$ is infinite. The family
$\{\,B_{1/2}(p):p\in E\,\}=\{\,\{p\}:p\in E\,\}$ is an open cover of $E$ in which each member contains
exactly one point of $E$. Any finite subfamily covers only finitely many points, so it omits a point of
the infinite set $E$; no finite subcover exists, and $E$ is not compact. Thus a subset is compact iff it
is finite.
\end{Solution}

\begin{Reflection}
The statement hands you an explicit two-valued formula for $d$ rather than a named space, so the first
job is to recognise the object: distances are only $0$ or $1$, which is the discrete metric of our notes
(\xrefL{2.2.4}). The word \emph{metric} then points at the four axioms (\xrefL{2.2.1}) as the thing to
verify, and the three follow-up questions---open, closed, compact---each name a definition to apply
(\xrefL{2.3.3}, \xrefL{3.2.4}, \xrefL{3.4.1}). The strategic clue is that in a metric space every
notion of nearness is read off the open balls, so before answering anything you compute the balls; here
they are so coarse that one calculation settles all three questions at once.

Believe the coarseness first. Picture the points of $X$ as scattered so that any two distinct ones are a
full unit apart---a cloud of isolated specks, no two ever closer. A ball of radius $\tfrac12$ then
reaches nobody but its own centre, and a ball of radius $2$ swallows the entire space; there is no
intermediate scale. That single picture is the whole problem.

The argument runs in four moves, each resting on the one before it or on a cited definition:
\begin{enumerate}[leftmargin=1.7em,itemsep=3pt,topsep=2pt,label=\textnormal{(\arabic*)}]
  \item the metric axioms (\xrefL{2.2.1}) reduce to one real check, the triangle inequality, and that
        is settled by a counting remark: a third point $r$ cannot coincide with two distinct points
        $p,q$ at once, so one of the legs $d(p,r),d(r,q)$ is already $1$ and dominates the unit left
        side;
  \item the crux computation $B_{1/2}(p)=\{p\}$ comes straight from the definition of the ball
        (\xrefL{2.3.1}), because no point other than $p$ lies within distance $1$;
  \item openness of every $E$ is then immediate---each point $p\in E$ has the witnessing ball
        $B_{1/2}(p)=\{p\}\subseteq E$, so $p$ is interior (\xrefL{2.3.3}); equivalently every singleton
        is an open ball (\factref{balls-open}{open balls are open}) and $E$ is their union
        (\factref{union-open}{unions of open sets are open});
  \item closedness needs no new work: with every set open, every set has open complement, so every set
        is closed (\xrefL{3.2.8}; \factref{open-iff-closed}{open $\iff$ complement closed}). The same
        ball calculation gives the direct reading too---no point is a cluster point of anything, since a
        ball around $p$ contains no point of a set other than $p$ itself (\xrefL{3.2.1};
        \factref{cluster-infinite}{cluster points are approached infinitely often}), so $E$ contains its
        empty set of cluster points vacuously.
\end{enumerate}

Compactness is the lone notion the coarse metric does \emph{not} trivialise, and seeing why sharpens the
whole space. The verdict ``compact $=$ finite'' is proved in both directions and both halves are
load-bearing: finiteness alone yields a finite subcover, while infiniteness is exploited by the cover of
singletons, where every member is forced to do the work of a single point so that nothing finite can
suffice. This is also where the infinitude of $X$ in the hypothesis finally matters---it guarantees
$X$ itself is a closed and bounded (every distance is $\le 1$) set that fails to be compact.

That failure is the reason to keep the example. Heine--Borel (\xrefL{4.3.5};
\factref{heine-borel}{``compact $\iff$ closed and bounded in $\mathbb R^k$''}) makes ``closed and
bounded'' equal to ``compact'' only inside $\mathbb R^k$; the discrete metric on an infinite set is the
clean witness that the equivalence is special to Euclidean space, since here closed-and-bounded sets are
everywhere yet only finite sets are compact. So when you need a metric in which some Euclidean intuition
breaks---a bounded set with no convergent subsequence, an infinite set with no cluster point, a space
where every set is clopen and hence totally disconnected (\xrefL{4.4.2})---the discrete metric is the
first place to reach. The transferable move is the one that opened the solution: in any metric space,
compute the open balls before answering a single topological question, because in a metric space every
property of nearness is encoded in them.
\end{Reflection}

% ------------------------------------------------------------
% ------------------------------------------------------------
\prob{11}{Rudin, Chapter 2, Exercise 11}
For $x\in\mathbb R^1$ and $y\in\mathbb R^1$, define
\[
\begin{aligned}
d_{1}(x,y)&=(x-y)^{2}, & d_{2}(x,y)&=\sqrt{\,|x-y|\,}, & d_{3}(x,y)&=|x^{2}-y^{2}|,\\
d_{4}(x,y)&=|x-2y|, & d_{5}(x,y)&=\frac{|x-y|}{1+|x-y|}. & &
\end{aligned}
\]
Determine, for each of these, whether it is a metric or not.

\begin{SolutionBlue}
The metrics are $d_{2}$ and $d_{5}$; the functions $d_{1},d_{3},d_{4}$ each fail one axiom of a metric
(\xrefL{2.2.1}). A metric on $\mathbb R$ must be nonnegative and vanish exactly on the diagonal
$x=y$, be symmetric, and satisfy the triangle inequality.

The function $d_{1}$ fails the triangle inequality: with $x=0$, $z=2$, and the intermediate point
$y=1$,
\[
d_{1}(0,2)=4>2=1+1=d_{1}(0,1)+d_{1}(1,2).
\]
The function $d_{3}$ fails definiteness, since $d_{3}(1,-1)=|1^{2}-(-1)^{2}|=0$ although $1\neq-1$;
distinct points sit at distance $0$. The function $d_{4}$ fails already at $d(x,x)=0$, since
$d_{4}(1,1)=|1-2\cdot 1|=1\neq 0$. One broken axiom is enough to disqualify each.

For $d_{2}$, nonnegativity, the vanishing $d_{2}(x,x)=\sqrt{0}=0$ together with $d_{2}(x,y)>0$ when
$x\neq y$, and symmetry are inherited directly from $|x-y|$, since the square root is nonnegative,
vanishes only at $0$, and ignores order. For the triangle inequality, fix $x,y,z\in\mathbb R$ and start
from the ordinary one, $|x-z|\le|x-y|+|y-z|$. The square root is increasing, so
\[
d_{2}(x,z)=\sqrt{|x-z|}\le\sqrt{\,|x-y|+|y-z|\,}\le\sqrt{|x-y|}+\sqrt{|y-z|}=d_{2}(x,y)+d_{2}(y,z),
\]
the last step being $\sqrt{a+b}\le\sqrt a+\sqrt b$ for $a,b\ge 0$, which holds because squaring the
right-hand side gives $a+b+2\sqrt{ab}\ge a+b$. Thus $d_{2}$ is a metric.

For $d_{5}$, write $f(t)=\dfrac{t}{1+t}$ for $t\ge 0$, so that $d_{5}(x,y)=f(|x-y|)$. Again
nonnegativity, definiteness ($f(t)=0$ iff $t=0$), and symmetry come from $|x-y|$. The map $f$ is
increasing on $[0,\infty)$, since $f(t)=1-\frac{1}{1+t}$ grows as $t$ does, and it is subadditive: for
$a,b\ge 0$,
\[
f(a+b)=\frac{a+b}{1+a+b}=\frac{a}{1+a+b}+\frac{b}{1+a+b}
\le\frac{a}{1+a}+\frac{b}{1+b}=f(a)+f(b),
\]
because $1+a+b\ge 1+a$ and $1+a+b\ge 1+b$ make each fraction only larger when the denominator shrinks.
Combining monotonicity with the ordinary triangle inequality,
\[
d_{5}(x,z)=f(|x-z|)\le f\bigl(|x-y|+|y-z|\bigr)\le f(|x-y|)+f(|y-z|)=d_{5}(x,y)+d_{5}(y,z).
\]
Hence $d_{5}$ is a metric.
\end{SolutionBlue}

\begin{Reflection}
The single word governing every part is \emph{metric}: each candidate is to be tested against the four
clauses of the definition (\xrefL{2.2.1})---nonnegativity, vanishing exactly on the diagonal, symmetry,
and the triangle inequality---and the verdict is whichever clause first holds or breaks. Because the
task splits cleanly into ``produce a counterexample'' for the failures and ``verify all four clauses''
for the successes, this is a proof by cases (\xrefL{0.5.7}) with two different burdens, and reading the
burden correctly is the whole skill. To \emph{refute} ``$d$ is a metric'' you negate a universally
quantified statement (\xrefL{0.2.5}), so a single triple $x,y,z$ suffices; to \emph{confirm} it you owe
all four clauses for all inputs.

That asymmetry sets the order of work: hunt first for a cheap failure, and only do the labour of
verification on what survives. The cheapest failures to scan for are a nonzero $d(x,x)$, two distinct
points at distance $0$, and a three-point violation of the triangle inequality---which is exactly how
$d_{4}$, $d_{3}$, and $d_{1}$ fall:
\begin{enumerate}[leftmargin=1.7em,itemsep=3pt,topsep=2pt,label=\textnormal{(\arabic*)}]
  \item $d_{4}$ dies at the first clause, since $d_{4}(1,1)=|1-2|=1\neq 0$: the formula $|x-2y|$ treats
        its two slots unequally, so a point is not even at distance $0$ from itself (it also fails
        symmetry, but one broken axiom is all the negation needs);
  \item $d_{3}$ dies at definiteness, since $x\mapsto x^{2}$ collapses $1$ and $-1$ to the same value,
        giving $d_{3}(1,-1)=0$ with $1\neq-1$---a metric must \emph{separate} distinct points, and any
        non-injective inner transformation forfeits that;
  \item $d_{1}$ satisfies the first three clauses but breaks the triangle inequality, because squaring
        inflates a long step out of proportion to two short ones: $d_{1}(0,2)=4$ exceeds
        $d_{1}(0,1)+d_{1}(1,2)=2$, the standard ``too convex'' failure.
\end{enumerate}

The two survivors share a structure worth naming: each is $g(|x-y|)$ for a function $g$ that fixes the
ordinary distance on $\mathbb R$. The first three clauses then transfer for free---$g(0)=0$ keeps the
diagonal, $g(t)>0$ for $t>0$ keeps separation, and $|x-y|$ is already symmetric. Only the triangle
inequality can be lost, and it is preserved precisely when $g$ is increasing and subadditive
($g(a+b)\le g(a)+g(b)$): feed the ordinary inequality $|x-z|\le|x-y|+|y-z|$ through the increasing $g$,
then split with subadditivity, exactly the two-line chain the solution runs for both. For $g=\sqrt{\cdot}$
subadditivity is $\sqrt{a+b}\le\sqrt a+\sqrt b$; for $g(t)=t/(1+t)$ it is the bounded-distance trick that
recurs throughout the subject, since this $g$ caps every distance below $1$ yet still defines the same
notion of nearness as $|x-y|$.

The cases are exhaustive---five functions, each assigned exactly one fate, with at least one explicit
axiom failing in every rejection and all four verified in every acceptance---and nothing is circular,
as the verifications lean only on the ordinary metric on $\mathbb R$ (\xrefL{2.2.3}) and on elementary
inequalities. The transferable move is this dichotomy in how you read ``is it a metric?'': a refutation
is one well-chosen triple, while a confirmation is four clauses earned in full, and the fastest path is
to chase the cheap counterexample before committing to the verification. The deeper lesson, the one the
survivors teach, is that an increasing subadditive reshaping of a metric is again a metric---the source
of $t/(1+t)$ as the standard way to make any metric bounded without disturbing its topology.
\end{Reflection}

% ------------------------------------------------------------
% ------------------------------------------------------------
\prob{12}{Rudin, Chapter 2, Exercise 12}
Let $K\subset\mathbb R^1$ consist of $0$ and the numbers $1/n$, for $n=1,2,3,\ldots$. Prove that $K$ is
compact directly from the definition (without using the Heine--Borel theorem).

\begin{Solution}
Write $K=\{0\}\cup\{1/n:n\ge 1\}$, and let $\{G_\alpha\}$ be an arbitrary open cover of $K$; we must
extract a finite subcover.

The point $0$ lies in some member $G_{\alpha_0}$ of the cover. Since $G_{\alpha_0}$ is open and $0$ is
one of its points, $0$ is an interior point (\xrefL{2.3.3}), so some ball $(-r,r)\subseteq G_{\alpha_0}$
with $r>0$. By the Archimedean property (\xrefL{1.7.1}) there is an integer $N\ge 1$ with $1/N<r$. For
every $n\ge N$ then $0<1/n\le 1/N<r$, so $1/n\in(-r,r)\subseteq G_{\alpha_0}$; and $0\in G_{\alpha_0}$ as
well. Thus the single set $G_{\alpha_0}$ already contains the entire tail
$\{0\}\cup\{1/n:n\ge N\}$.

What remains uncovered lies in the finite set $\{1,1/2,\ldots,1/(N-1)\}$. For each such point choose one
member of the cover containing it, say $G_{\alpha_1},\ldots,G_{\alpha_{N-1}}$. The finite collection
\[
G_{\alpha_0},\,G_{\alpha_1},\,\ldots,\,G_{\alpha_{N-1}}
\]
then covers all of $K$: the tail by $G_{\alpha_0}$, each leftover by its chosen set. Every open cover of
$K$ admits a finite subcover, so $K$ is compact by definition (\xrefL{3.4.1}).
\end{Solution}

\begin{Reflection}
The phrase ``directly from the definition (without using Heine--Borel)'' is an explicit instruction
about which tool to reach for. The shortcut would be to remark that $K\subseteq[0,1]$ is closed and
bounded and quote \factref{heine-borel}{Heine--Borel} (\xrefL{4.3.5}); the problem forbids exactly that,
leaving only the raw definition of compactness---every open cover has a finite subcover (\xrefL{3.4.1}).
So the task is fixed in advance: start from an \emph{arbitrary} open cover $\{G_\alpha\}$ and produce a
finite subcollection that still covers, with the obstacle that $K$ is infinite while a subcover may use
only finitely many sets.

The second clue is the shape of $K$ itself: a sequence $1/n$ together with the one point $0$ it
approaches. That $0$ is the lone accumulation point, and reading it as a limit is the whole idea, since
$1/n\to 0$ means every neighbourhood of $0$ swallows all but finitely many terms
(\factref{limit-tail}{the tail lands in every ball about the limit}, \xrefL{5.1.3}). A cover set that
catches $0$ is forced, by openness, to contain a whole interval around $0$, and that interval drags in
the entire tail for free. Picture it on the line: pin an interval at $0$ and only the few largest terms
$1,1/2,\ldots$ poke out to its right; everything else is already inside.

That picture is the proof, step by step:
\begin{enumerate}[leftmargin=1.7em,itemsep=3pt,topsep=2pt,label=\textnormal{(\arabic*)}]
  \item some $G_{\alpha_0}$ contains $0$ because $\{G_\alpha\}$ covers $K$ and $0\in K$, and openness
        makes $0$ an interior point (\xrefL{2.3.3}), so a ball $(-r,r)\subseteq G_{\alpha_0}$;
  \item the Archimedean property (\xrefL{1.7.1}; \factref{archimedean}{$\inf\{1/n\}=0$}) supplies an
        $N$ with $1/N<r$, and then $1/n\le 1/N<r$ for all $n\ge N$, so this one set $G_{\alpha_0}$ holds
        the entire infinite tail $\{0\}\cup\{1/n:n\ge N\}$;
  \item only the finitely many points $1,1/2,\ldots,1/(N-1)$ are left, and finitely many points are
        harmless---assign one cover member to each;
  \item $G_{\alpha_0}$ together with those finitely many sets is a finite subcover, which is the
        definition (\xrefL{3.4.1}) satisfied.
\end{enumerate}

Each hypothesis earns its place, and dropping the point $0$ shows why. The bare set $\{1/n:n\ge 1\}$ is
\emph{not} compact: cover it by tiny disjoint intervals, one snug around each $1/n$, and no finite
subfamily reaches the points near $0$ (\xrefL{3.4.2}). What rescues $K$ is precisely that $0$ has been
included, so a single cover set is compelled to absorb the tail; without it nothing forces that
absorption. The Archimedean step is equally load-bearing---it is what turns ``$1/n$ gets small'' into a
concrete cutoff $N$ that splits $K$ cleanly into a covered tail and a finite head.

The transferable move is the head--tail split for compactness from scratch: find the one point that
catches infinitely many elements, let the open set around it inhale the whole tail (this is where
$x_n\to p$ does the work), and mop up the finite remainder one set at a time. It is the model argument
for ``a convergent sequence together with its limit is compact,'' and it is exactly the engine behind
the more elaborate construction in \hyperlink{prob:13}{Problem~13}.
\end{Reflection}

% ------------------------------------------------------------
% ------------------------------------------------------------
\prob{13}{Rudin, Chapter 2, Exercise 13}
Construct a compact set of real numbers whose limit points form a countable set.

\begin{SolutionBlue}
Push the construction of \hyperlink{prob:5}{Problem~5} one level deeper: instead of three target points,
take the countably many targets $1,\tfrac12,\tfrac13,\dots$, attach a sequence accumulating at each, and add
the point $0$ at which the targets themselves accumulate. With $n,k$ ranging over the positive integers, set
\[
K=\{0\}\;\cup\;\Bigl\{\tfrac1n:n\ge 1\Bigr\}\;\cup\;\Bigl\{\tfrac1n+\tfrac{1}{nk}:n,k\ge 1\Bigr\},
\qquad
L=\{0\}\cup\Bigl\{\tfrac1n:n\ge 1\Bigr\}.
\]
The verdict: $K$ is compact, and its set of cluster points is exactly $L$, which is countable.

The set is bounded. For every $n,k\ge 1$ we have $0<\tfrac1n\le 1$ and $0<\tfrac1n+\tfrac1{nk}\le 1+1=2$,
the maximum occurring at $n=k=1$; so $K\subseteq[0,2]$ is bounded (\xrefL{4.3.4}).

Every point of $L$ is a cluster point of $K$. Fix $n$ and any radius $r>0$. By the Archimedean property
(\factref{archimedean}{Archimedean}) there is a $k$ with $\tfrac{1}{nk}<r$; the point $\tfrac1n+\tfrac1{nk}$
lies in $K$, differs from $\tfrac1n$, and sits within $r$ of $\tfrac1n$, so $\tfrac1n$ is a cluster point.
For $0$: given $r>0$, choose $n$ with $\tfrac1n<r$; then $\tfrac1n\in K$, $\tfrac1n\neq 0$, and $\tfrac1n$ is
within $r$ of $0$, so $0$ is a cluster point as well (\xrefL{3.2.1}).

No point outside $L$ is a cluster point. The key is that $L$ is closed: its only cluster point is $0$, since
$\tfrac1n\to 0$ while the terms march away to the right (\xrefL{3.2.5}), and $0\in L$. Take $p\notin L$; in
particular $p\neq 0$. If $p<0$ then the ball $B(p,|p|)$ misses $K\subseteq[0,2]$ entirely, so $p$ is trivially
not a cluster point; assume henceforth $p>0$. Because
$p\notin L$ and $L$ is closed, $p$ is not a cluster point of $L$, so some ball $B(p,\rho)$ with $\rho>0$ is
disjoint from $L$; shrinking $\rho$, arrange in addition that $\rho<\tfrac{p}{2}$, so that
$B(p,\rho)\subseteq(\tfrac{p}{2},\,p+\rho)$ stays a fixed distance $\tfrac{p}{2}>0$ to the right of $0$. I claim
$B(p,\rho)$ contains only finitely many points of $K$. Each remaining family $\{\tfrac1n+\tfrac1{nk}:k\ge1\}$
converges to $\tfrac1n$ as $k\to\infty$, and $\tfrac1n+\tfrac1{nk}\le\tfrac2n$, so once $\tfrac2n\le\tfrac{p}{2}$ the
\emph{entire} family for that $n$ lies inside $(0,\tfrac{p}{2}]$, hence to the left of $B(p,\rho)$ and disjoint
from it. Only the finitely many indices $n$ with $\tfrac2n>\tfrac{p}{2}$ survive; for each such
$n$ the sequence $\tfrac1n+\tfrac1{nk}\to\tfrac1n\notin B(p,\rho)$, so only finitely many of its terms can
lie in $B(p,\rho)$. A finite union of finite sets is finite, so $B(p,\rho)\cap K$ is finite, and a cluster
point would force infinitely many points of $K$ into every ball around it
(\factref{cluster-infinite}{a cluster point is approached infinitely often}). Hence $p$ is not a cluster
point of $K$.

The cluster points of $K$ are thus exactly the points of $L=\{0\}\cup\{\tfrac1n\}$, a subset of the countable
set $\mathbb Q$, hence countable (\factref{subset-countable}{a subset of a countable set is countable}).
Finally $K$ contains all its cluster points, so $K$ is closed; being closed and bounded in $\mathbb R$, it is
compact by Heine--Borel (\xrefL{4.3.5}; \factref{heine-borel}{closed and bounded $\Rightarrow$ compact}).
\end{SolutionBlue}

\begin{Reflection}
Two demands sit in the statement and they pull against each other, which is the first thing to read.
\emph{Compact} in $\mathbb R$ means closed and bounded (\factref{heine-borel}{Heine--Borel}, \xrefL{4.3.5}),
and ``closed'' insists the set already \emph{owns} every one of its cluster points (\xrefL{3.2.5}); yet the
problem wants the cluster set to be \emph{countably infinite}---an honest, unbounded-in-cardinality supply of
accumulation points. So the design problem is to engineer infinitely many places where the set piles up
(\xrefL{3.2.1}), keep them all inside a bounded interval, and then make sure not a single accumulation point
has been left out, or closedness fails and compactness with it.

The cheapest accumulation point, as in \hyperlink{prob:5}{Problem~5}, is a sequence converging to a limit it
never reaches---$\tfrac1n$ bunching against $0$. There the trick was three such strings aimed at three
separated targets. The leap here is to let the \emph{targets} be a convergent sequence in their own right:
aim one little string at each of $1,\tfrac12,\tfrac13,\dots$, by hanging $\tfrac1n+\tfrac1{nk}$ ($k\to\infty$)
off the point $\tfrac1n$. Believe the picture before proving anything: clusters of dots bunch against each
$\tfrac1n$, and the $\tfrac1n$ themselves bunch against $0$---accumulation at two scales, a fractal-flavoured
staircase. That second scale is exactly why $0$ must be thrown in by hand; it is a cluster point of the
targets, so it is a cluster point of $K$, and an unowned cluster point is precisely what would make $K$
non-closed.

The proof then has four moves, each resting on the one before or on a quotable fact:
\begin{enumerate}[leftmargin=1.7em,itemsep=3pt,topsep=2pt,label=\textnormal{(\arabic*)}]
  \item boundedness is read straight off the formula---$\tfrac1n+\tfrac1{nk}\le 2$ with equality at
        $n=k=1$, so $K\subseteq[0,2]$ (\xrefL{4.3.4})---and it never comes into doubt again;
  \item that every $\tfrac1n$ and $0$ \emph{is} a cluster point is the easy direction, falling out of the
        two convergences, with the Archimedean property (\factref{archimedean}{Archimedean}) producing, for
        any radius, an index that lands a genuine point of $K$ inside the ball;
  \item that \emph{nothing else} is a cluster point is the real labour, and the lever is that a cluster
        point drags infinitely many points of $K$ into every ball around it
        (\factref{cluster-infinite}{a cluster point is approached infinitely often}); so a candidate
        $p\notin L$ is acquitted the moment one ball around it is shown to meet $K$ only finitely
        often---which works because $L$ is closed, giving a ball about $p$ that avoids $L$ and $0$, after
        which the bound $\tfrac1n+\tfrac1{nk}\le\tfrac2n$ sweeps all but finitely many whole families down
        near $0$ and out of the ball;
  \item with the cluster set pinned to $L=\{0\}\cup\{\tfrac1n\}\subseteq\mathbb Q$, countability is immediate
        (\factref{subset-countable}{a subset of a countable set is countable}), and $K\supseteq L$ means $K$
        owns its cluster points, so $K$ is closed and therefore compact (\xrefL{4.3.5}).
\end{enumerate}

The rigor hinges on two places where the construction could quietly leak. First, the secondary limit $0$:
omit it and $K$ has a cluster point it does not contain, so $K$ fails to be closed and Heine--Borel never
fires---this is the single most common way the problem is botched. Second, the uniform bound
$\tfrac1n+\tfrac1{nk}\le\tfrac2n$ is what collapses an \emph{entire} family down toward $0$ at once, clear of a
fixed ball about $p$; without controlling the whole family for large $n$ one cannot reduce to finitely many
indices, and the ``finitely many points'' argument stalls. Both halves of ``the cluster set is exactly $L$'' are needed:
showing $L$ is contained would not stop some rogue point from also being a cluster point, and showing nothing
outside $L$ qualifies would not confirm the $\tfrac1n$ actually do.

The transferable move is the engineering recipe and its closure tax. To prescribe the cluster set of a
compact set, choose your target points, hang a convergent sequence on each to force accumulation exactly
there, and then---this is the part the bookkeeping demands---include every cluster point of the targets
themselves, so that the finished set owns all its accumulation points and is closed
(\factref{closure-closed}{adding the cluster points closes a set}). Countably many targets accumulating at a
single new point give a countable cluster set; the same idea iterated builds the Cantor--Bendixson towers
that classify how complicated a closed set's accumulation structure can be.
\end{Reflection}

% ------------------------------------------------------------
% ------------------------------------------------------------
\prob{14}{Rudin, Chapter 2, Exercise 14}
Give an example of an open cover of the segment $(0,1)$ which has no finite subcover.

\begin{SolutionBlue}
For each integer $n\ge 2$ put
\[
G_n=\Bigl(\tfrac1n,\,1\Bigr),
\]
an open interval, and consider the family $\{G_n:n\ge 2\}$.

This is an open cover of $(0,1)$. Given $x\in(0,1)$, the Archimedean property supplies an integer
$n$ with $n>1/x$; then $\tfrac1n<x<1$, so $x\in G_n$. Hence $(0,1)\subseteq\bigcup_{n\ge 2}G_n$, and
since each $G_n\subseteq(0,1)$ the union is exactly $(0,1)$.

No finite subcollection covers $(0,1)$. Suppose $G_{n_1},\dots,G_{n_k}$ are finitely many of these
sets, and let $N=\max\{n_1,\dots,n_k\}$. The intervals are nested, $G_n\subseteq G_{N}$ whenever
$n\le N$, so
\[
\bigcup_{j=1}^{k}G_{n_j}=G_N=\Bigl(\tfrac1N,\,1\Bigr).
\]
This set misses the point $\tfrac1N\in(0,1)$ (in fact every point of $(0,\tfrac1N]$), so it does not
cover $(0,1)$. Thus the open cover $\{G_n\}$ admits no finite subcover.
\end{SolutionBlue}

\begin{Reflection}
The phrase ``open cover \ldots\ with no finite subcover'' is the literal negation of the definition of
compactness (\xrefL{3.4.1}): a set is compact when \emph{every} open cover has a finite subcover, so
to certify a set is not compact you produce one cover that defeats every finite selection. The task is
therefore not to study $(0,1)$ in the abstract but to engineer a single bad cover, and the shape of
$(0,1)$---an open interval missing its endpoints---tells you where to aim. The trouble must sit at a
missing endpoint, since that is the only thing distinguishing $(0,1)$ from the compact $[0,1]$; the
notes flag exactly this in their two standard non-compact examples (\xrefL{3.4.2}).

Picture it before proving it. Slide a window $(\tfrac1n,1)$ leftward as $n$ grows: each window swallows
more of the interval, and any fixed $x>0$ is eventually inside one. But the left edges $\tfrac1n$ march
toward $0$ without ever arriving, so no \emph{single} window---and hence no finite batch of them, which
reaches only as far left as the smallest edge---can reach all the way down to $0$. The cover chases the
absent endpoint $0$ and never catches it.

The argument has two halves, and each rests on something concrete:
\begin{enumerate}[leftmargin=1.7em,itemsep=3pt,topsep=2pt,label=\textnormal{(\arabic*)}]
  \item the family covers $(0,1)$ because every $x>0$ has some $\tfrac1n<x$, which is precisely the
        Archimedean property (\factref{archimedean}{$\inf\{1/n\}=0$}); this is the only place the
        ambient field does any work;
  \item no finite subfamily covers, because finitely many indices have a largest one $N$, the intervals
        are nested so their union collapses to the single largest interval $(\tfrac1N,1)$, and that
        interval visibly omits $\tfrac1N$.
\end{enumerate}

The rigor turns entirely on the word \emph{finite} in step (2): an infinite subfamily---say all of
them---does cover, so the obstruction is not the cover itself but the impossibility of choosing finitely
many edges that reach $0$. Take a finite set of indices, extract its maximum, and the nesting hands you
a single offending interval; no maximum exists for the full infinite family, which is exactly why the
whole cover succeeds where every finite part fails. That contrast is the entire content. (One can also
confirm non-compactness through Heine--Borel: $(0,1)$ is bounded but not closed, since $0$ is a cluster
point lying outside it, and \factref{heine-borel}{a compact subset of $\mathbb R^k$ is closed and
bounded} would force closedness. But the exercise asks for the explicit cover, which exhibits the
failure directly rather than appealing to a theorem.)

The transferable move is the recipe for ``show $X$ is not compact'': find the point the space is missing
and build a cover whose members each stop just short of it, indexed so that any finite subcollection has
a last member and therefore still falls short. The same template handles $(0,1/n)$ shrinking to a
missing $0$ or $[n,\infty)$ running off to infinity, the two failures the notes single out
(\xrefL{3.4.2}); a finite subfamily always has an extremal index, and that is what you turn against the
cover.
\end{Reflection}

% ------------------------------------------------------------
% ------------------------------------------------------------
\prob{15}{Rudin, Chapter 2, Exercise 15}
Show that Theorem~2.36 and its Corollary become false (in $\mathbb R^{1}$, for example) if the word
``compact'' is replaced by ``closed'' or by ``bounded.''

\begin{SolutionBlue}
Rudin's Theorem~2.36 is our finite-intersection property (\xrefL{4.2.1}): a family of compact sets
whose finite subfamilies all have nonempty intersection has nonempty total intersection. Its Corollary
is our nested statement (\xrefL{4.2.2}, the quotable \factref{nested-compact}{nested nonempty compact
sets meet}): a decreasing chain $K_1\supseteq K_2\supseteq\cdots$ of nonempty compact sets has
$\bigcap_n K_n\neq\varnothing$. Each fails over $\mathbb R^{1}$ once ``compact'' is weakened to
``closed'' or to ``bounded,'' and a single nested family settles both the theorem and its corollary at
once, since a nested chain automatically has the finite-intersection property.

Weaken ``compact'' to ``closed'' first. For $n\in\mathbb N$ with $n\ge 1$ put
\[
  F_n=[n,\infty).
\]
Each $F_n$ is nonempty and closed (its complement $(-\infty,n)$ is open, \xrefL{3.2.8}), and the chain
is nested: $F_1\supseteq F_2\supseteq\cdots$. Any finite subfamily intersects in its smallest member,
\[
  F_{n_1}\cap\cdots\cap F_{n_k}=[N,\infty)\neq\varnothing,\qquad N=\max\{n_1,\dots,n_k\},
\]
so the family has the finite-intersection property. Yet $\bigcap_{n\ge 1}F_n=\varnothing$: a point in
the intersection would be a real number $x\ge n$ for every $n$, which the Archimedean property forbids
(\factref{archimedean}{$\mathbb R$ is Archimedean}, \xrefL{1.7.1}). Thus closed, nonempty, nested sets
with the finite-intersection property can have empty intersection.

Now weaken ``compact'' to ``bounded'' instead. For $n\in\mathbb N$ with $n\ge 1$ put
\[
  B_n=\Bigl(0,\tfrac1n\Bigr).
\]
Each $B_n$ is nonempty and bounded (contained in $[0,1]$, \xrefL{4.3.4}), and the chain is nested:
$B_1\supseteq B_2\supseteq\cdots$. Any finite subfamily intersects in its smallest member,
\[
  B_{n_1}\cap\cdots\cap B_{n_k}=\Bigl(0,\tfrac1N\Bigr)\neq\varnothing,\qquad N=\max\{n_1,\dots,n_k\},
\]
so again the finite-intersection property holds. Yet $\bigcap_{n\ge 1}B_n=\varnothing$: a point in the
intersection would be a real $x$ with $0<x<\tfrac1n$ for every $n$, contradicting that
$\inf\{1/n\}=0$, the Archimedean property once more (\factref{archimedean}{$\mathbb R$ is Archimedean},
\xrefL{1.7.1}). Thus bounded, nonempty, nested sets with the finite-intersection property can have
empty intersection.

In both families the sets are nested, so each is simultaneously a counterexample to Theorem~2.36 (via
the finite-intersection property) and to its nested Corollary. What each family drops is one half of
``closed and bounded,'' which by Heine--Borel (\factref{heine-borel}{compact in $\mathbb R^k$ means
closed and bounded}, \xrefL{4.3.5}) is exactly compactness: $\{[n,\infty)\}$ keeps closed but discards
bounded, and $\{(0,1/n)\}$ keeps bounded but discards closed.
\end{SolutionBlue}

\begin{Reflection}
The instruction is unusually direct---``show the theorem becomes false if you weaken the hypothesis''---so the reading is not which tool to invoke but which hypothesis to attack. Theorem~2.36, our finite-intersection property (\xrefL{4.2.1}), and its Corollary, the nested statement
(\xrefL{4.2.2}, the quotable \factref{nested-compact}{nested nonempty compact sets meet}), both turn a
local promise (every finite subfamily meets) into a global one (the whole family meets), and the only
hypothesis paying for that leap is compactness. So the task is to find a family that satisfies
everything \emph{except} compactness and still fails to meet---one family per weakening the problem
names, ``closed'' and ``bounded.''

The right picture comes from asking what compactness buys here. In $\mathbb R^{1}$ Heine--Borel
(\factref{heine-borel}{closed and bounded}, \xrefL{4.3.5}) factors compactness into two independent
conditions, so to break the theorem you keep one and sabotage the other, and there are exactly two ways
a nested chain of nonempty sets can drain to nothing: it can escape to infinity, or it can shrink onto
a point it never reaches. Escaping to infinity, $[n,\infty)$, stays closed but is unbounded; shrinking
to a missing endpoint, $(0,1/n)$, stays bounded but is not closed. Each keeps exactly the half the
other discards.

Watching the closed family run shows where each requirement is spent:
\begin{enumerate}[leftmargin=1.7em,itemsep=3pt,topsep=2pt,label=\textnormal{(\arabic*)}]
  \item the sets $F_n=[n,\infty)$ are closed because each complement $(-\infty,n)$ is open
        (\xrefL{3.2.8}), and they are nested by construction, $F_1\supseteq F_2\supseteq\cdots$;
  \item nestedness hands over the finite-intersection property for free---a finite subfamily meets in
        its smallest member $[N,\infty)$, which is nonempty---so the hypothesis of Theorem~2.36 is
        genuinely satisfied and not quietly skipped;
  \item the intersection is empty because no real number exceeds every natural number, which is the
        Archimedean property (\factref{archimedean}{$\mathbb R$ is Archimedean}, \xrefL{1.7.1})---the
        single fact that powers the collapse, and the thing compactness would have prevented by
        forbidding the escape to infinity.
\end{enumerate}
The bounded family $B_n=(0,1/n)$ runs identically, with the same Archimedean fact appearing as
$\inf\{1/n\}=0$: the chain shrinks toward $0$, but $0$ was deleted from every $B_n$, so nothing
survives.

The rigour to check is that the finite-intersection property really holds---it does, and nestedness is
why, which is also the economy of the example: a nested chain is automatically a finite-intersection
family, so one family disproves both Theorem~2.36 and its nested Corollary, no second construction
needed. It is worth seeing that neither half alone is enough and each is load-bearing: drop bounded and
$[n,\infty)$ kills it; drop closed and $(0,1/n)$ kills it; restore both and Heine--Borel makes the sets
compact, at which point \xrefL{4.2.2} applies and the intersection must be nonempty. There is no
circularity---the counterexamples never assume the conclusion, they exhibit it failing.

The transferable move is the diagnostic habit behind ``show this becomes false.'' When a theorem
bundles two conditions into one named hypothesis, a sharp counterexample keeps one condition and breaks
the other, isolating which half does the work; here Heine--Borel is the bundle, and the two escapes a
nested real chain can make---off to infinity, or onto a deleted limit---are precisely the two ways
compactness can be missing while the finite-intersection property survives.
\end{Reflection}

% ------------------------------------------------------------
% ------------------------------------------------------------
\prob{16}{Rudin, Chapter 2, Exercise 16}
Regard $\mathbb Q$, the set of all rational numbers, as a metric space, with $d(p,q)=|p-q|$. Let $E$
be the set of all $p\in\mathbb Q$ such that $2<p^2<3$. Show that $E$ is closed and bounded in
$\mathbb Q$, but that $E$ is not compact. Is $E$ open in $\mathbb Q$?

\begin{Solution}
$E$ is bounded and is both closed and open in $\mathbb Q$, yet it is not compact; so the answer to the
last question is yes.

No rational squares to $2$ or to $3$ (\xrefL{1.4.1}; the same parity/divisibility argument settles
$3$), so for $p\in\mathbb Q$ the condition $2<p^2<3$ is exactly $\sqrt2<|p|<\sqrt3$. Writing
\[
U=(-\sqrt3,-\sqrt2)\cup(\sqrt2,\sqrt3),\qquad
C=[-\sqrt3,-\sqrt2]\cup[\sqrt2,\sqrt3],
\]
both subsets of $\mathbb R$, we therefore have $E=\mathbb Q\cap U$ \emph{and} $E=\mathbb Q\cap C$: since
$\pm\sqrt2,\pm\sqrt3\notin\mathbb Q$, adjoining or deleting those four endpoints changes nothing inside
$\mathbb Q$.

The set $E$ is bounded, since every $p\in E$ has $p^2<3<4$, so $|p|<2$; thus $E\subseteq B_2(0)$ in
$\mathbb Q$.

It is open in $\mathbb Q$, being the trace $E=\mathbb Q\cap U$ of the set $U$, which is open in
$\mathbb R$; the subspace criterion that the open sets of $\mathbb Q$ are exactly the intersections of
$\mathbb Q$ with open sets of $\mathbb R$ then applies (\xrefL{3.3.1}).

It is also closed in $\mathbb Q$. Since $C$ is closed in $\mathbb R$, its complement $\mathbb R\setminus
C$ is open (\xrefL{3.2.8}), so the trace $\mathbb Q\setminus E=\mathbb Q\cap(\mathbb R\setminus C)$ is
open in $\mathbb Q$ by the same subspace criterion, and its complement $E$ is therefore closed in
$\mathbb Q$ (\xrefL{3.2.8}).

The set $E$ is not compact. For each integer $n\ge 1$ set
\[
G_n=\Bigl\{q\in\mathbb Q:\ q^2<3-\tfrac1n\Bigr\}=\mathbb Q\cap\bigl(-\sqrt{3-\tfrac1n},\ \sqrt{3-\tfrac1n}\bigr),
\]
open in $\mathbb Q$ as the trace of an open interval. These cover $E$: if $q\in E$ then $3-q^2>0$, so by
the Archimedean property (\xrefL{1.7.1}; \factref{archimedean}{List of Facts}) $\tfrac1n<3-q^2$ for some
$n$, whence $q\in G_n$. The $G_n$ increase with $n$, so any finite subfamily lies inside a single
$G_N$, the one of largest index. But $G_N$ omits every $q\in E$ with $\sqrt{3-\tfrac1N}\le|q|<\sqrt3$,
and such a rational exists: since $\mathbb Q$ is dense in $\mathbb R$ (\xrefL{1.7.2};
\factref{q-dense}{List of Facts}), the interval $\bigl(\sqrt{3-\tfrac1N},\sqrt3\bigr)$ contains some
$q\in\mathbb Q$, and then $2<q^2<3$ puts $q\in E\setminus G_N$. No finite subfamily covers $E$, so $E$
is not compact (\xrefL{3.4.1}).
\end{Solution}

\begin{Reflection}
The phrase that should stop you is ``regard $\mathbb Q$ \emph{as a metric space}'': the problem is
quietly warning that every topological word---open, closed, compact, bounded---is to be read inside
$\mathbb Q$, not inside $\mathbb R$. In $\mathbb R^k$ those words are tightly linked by Heine--Borel,
which equates compactness with closed-and-bounded (\xrefL{4.3.5}; \factref{heine-borel}{List of
Facts}); seeing a set that is closed and bounded yet \emph{not} compact is the standard signal that the
ambient space has been changed, and that the link Heine--Borel forges is a property of $\mathbb R^k$
rather than of sets in the abstract. So the real content is to take each property relative to
$\mathbb Q$ and watch what survives.

Believe it through one picture before proving anything. In $\mathbb R$ the set lives in two open
intervals with the irrational endpoints $\pm\sqrt2,\pm\sqrt3$; passing to $\mathbb Q$ deletes those
endpoints from the universe entirely. A set whose only boundary points have been removed from the space
has no boundary left to own or to miss, so it comes out both open and closed---and a sequence in $E$
creeping up toward $\sqrt3$ is trying to converge to a point that simply is not there, which is exactly
how compactness will die.

The cleanest route is to stop fighting the metric pointwise and use the subspace description. The
engine is that $E$ is squeezed between irrational thresholds, so it equals the trace on $\mathbb Q$ of
\emph{both} the open set $U$ and the closed set $C$ that differ only at those absent endpoints:
\begin{enumerate}[leftmargin=1.7em,itemsep=3pt,topsep=2pt,label=\textnormal{(\arabic*)}]
  \item because no rational squares to $2$ or $3$ (\xrefL{1.4.1}), the rational solutions of $2<p^2<3$
        are the same whether the endpoints are included or not, giving the two readings $E=\mathbb
        Q\cap U=\mathbb Q\cap C$ on which everything else rests;
  \item boundedness is immediate from $p^2<3<4$, so $|p|<2$ for all $p\in E$, with no reference to the
        space at all;
  \item openness in $\mathbb Q$ is the trace of the open $U$, by the subspace criterion that the open
        sets of $\mathbb Q$ are exactly the intersections $\mathbb Q\cap V$ with $V$ open in $\mathbb R$
        (\xrefL{3.3.1});
  \item closedness in $\mathbb Q$ follows by taking complements: $C$ closed makes $\mathbb R\setminus C$
        open (\xrefL{3.2.8}), its trace $\mathbb Q\setminus E$ is open in $\mathbb Q$ by the same
        criterion, and so $E$, its complement, is closed (\factref{open-iff-closed}{List of Facts});
  \item noncompactness is exhibited by hand: the increasing cover $G_n=\{q\in\mathbb Q:q^2<3-\tfrac1n\}$
        catches each point of $E$ once $\tfrac1n$ drops below $3-q^2$, which the Archimedean property
        guarantees (\xrefL{1.7.1}), yet a finite subfamily reduces to the largest $G_N$, and density of
        $\mathbb Q$ (\xrefL{1.7.2}) plants a rational in $\bigl(\sqrt{3-\tfrac1N},\sqrt3\bigr)$ that
        $G_N$ leaves uncovered, so no finite subfamily suffices (\xrefL{3.4.1}).
\end{enumerate}

The argument is honest about where each space-dependent fact is spent. Openness and closedness both
ride entirely on the endpoints being irrational---were $\sqrt3$ rational, $C$ and $U$ would no longer
have the same trace and $E$ would lose one of the two properties. The noncompactness needs both the
Archimedean property (every point is eventually covered) and density (the finite stage always leaves a
rational stranded near the missing $\sqrt3$); drop density and one could not produce the witness
$q\in E\setminus G_N$. There is no circularity: the cover is built explicitly and the failure is shown
directly from the definition of compactness, not from any converse of Heine--Borel.

The move to carry away is the subspace dictionary itself---a set is open (closed) in $\mathbb Q$
exactly when it is the trace of an open (closed) set of $\mathbb R$, so questions about the induced
metric become questions about $\mathbb R$ that you already know how to answer. And the cautionary half
is the lasting one: ``closed and bounded'' is only as strong as the space it lives in. Heine--Borel
(\xrefL{4.3.5}) is a theorem about $\mathbb R^k$; over $\mathbb Q$ a set can be closed, bounded, and
still riddled with missing limits, and those holes---here the irrational $\pm\sqrt2,\pm\sqrt3$---are
precisely what a finite subcover cannot corner.
\end{Reflection}

% ------------------------------------------------------------
% ------------------------------------------------------------
\prob{17}{Rudin, Chapter 2, Exercise 17}
Let $E$ be the set of all $x\in[0,1]$ whose decimal expansion contains only the digits $4$ and $7$.
Is $E$ countable? Is $E$ dense in $[0,1]$? Is $E$ compact? Is $E$ perfect?

\begin{SolutionBlue}
$E$ is uncountable, not dense in $[0,1]$, compact, and perfect.

Each point of $E$ is determined by a choice of one digit from $\{4,7\}$ in every decimal place, so write
\[
x=0.a_1a_2a_3\cdots=\sum_{n=1}^{\infty}a_n\,10^{-n},\qquad a_n\in\{4,7\}.
\]
Distinct digit sequences give distinct points: if $(a_n)$ and $(b_n)$ first differ at place $m$, with
$a_m>b_m$ say, then $a_m-b_m=3$, while the tails contribute at most
$\sum_{n>m}3\cdot 10^{-n}=10^{-m}/3<3\cdot 10^{-m}$ to either number, so the two values cannot
coincide. Hence the map $(a_n)\mapsto x$ is a bijection of $\{4,7\}^{\mathbb N}$ onto $E$. The set of
all such sequences is uncountable (\factref{binary-uncountable}{\xrefL{2.1.10}}, relabelling $\{4,7\}$
as $\{0,1\}$), so $E$ is uncountable.

$E$ is not dense in $[0,1]$. Every $x\in E$ has first digit $4$ or $7$, so
$x\ge 0.4\overline{4}=4/9$; thus $E\cap(0,4/9)=\varnothing$. A point like $1/9=0.111\cdots$ has a
neighbourhood, say $(0,4/9)$, that meets no point of $E$, so $\overline E\ne[0,1]$.

$E$ is closed. For $m\ge 1$ and each finite block $(a_1,\dots,a_m)\in\{4,7\}^m$ form the length-$10^{-m}$
closed interval
\[
I_{(a_1,\dots,a_m)}=\Bigl[\,s,\ s+10^{-m}\Bigr],\qquad s=\sum_{j=1}^{m}a_j\,10^{-j},
\]
and let $F_m=\bigcup_{(a_1,\dots,a_m)\in\{4,7\}^m}I_{(a_1,\dots,a_m)}$ be the union of the
$2^m$ of them. Each $F_m$ is a finite union of closed sets, hence closed
(\factref{inter-closed}{\xrefL{3.2.10}}). I claim
\[
E=\bigcap_{m=1}^{\infty}F_m.
\]
If $x\in E$ then for every $m$ its first $m$ digits lie in $\{4,7\}$, so $x$ lies in the corresponding
interval and $x\in F_m$. Conversely suppose $x\in\bigcap_m F_m$; for each $m$ the interval of $F_m$
containing $x$ pins the first $m$ decimal digits of $x$ to values in $\{4,7\}$, and these choices are
consistent as $m$ grows, so $x$ has a decimal expansion using only $4$ and $7$ and $x\in E$. Being an
intersection of closed sets, $E$ is closed (\factref{inter-closed}{\xrefL{3.2.10}}).

$E$ is compact. It is closed, and $E\subseteq[0,1]$ is bounded, so by Heine--Borel a closed bounded
subset of $\mathbb R$ is compact (\factref{heine-borel}{\xrefL{4.3.5}}).

$E$ is perfect: it is closed, and every point is a limit point of $E$. Fix
$x=0.a_1a_2\cdots\in E$ and any radius $\eps>0$. Choose $m$ with $3\cdot 10^{-m}<\eps$, and let $x_m\in
E$ be the point obtained by switching the $m$th digit ($4\leftrightarrow7$) and keeping the rest. Then
$x_m\in E$, $x_m\ne x$, and $|x_m-x|=3\cdot 10^{-m}<\eps$, so every ball about $x$ contains a point of
$E$ other than $x$. Thus $x$ is a cluster point of $E$ (\xrefL{3.2.1}). A closed set all of whose points
are cluster points is perfect, so $E$ is perfect.
\end{SolutionBlue}

\begin{Reflection}
The statement is a checklist---countable, dense, compact, perfect---and the way to read it is to notice
what the \emph{description} of $E$ controls: each point is a free choice of one of \emph{two} symbols in
\emph{every} decimal place. ``Two choices in each of countably many slots'' is the exact signature of the
binary sequences, whose uncountability is the headline theorem of the counting lecture
(\factref{binary-uncountable}{\xrefL{2.1.10}}); so the moment you see ``only the digits $4$ and $7$'' you
should expect $E$ to be in bijection with $\{0,1\}^{\mathbb N}$ and hence uncountable. The remaining three
questions are topological, and the word that unlocks all of them is the one hiding in ``compact'' and
``perfect'': both are built on \emph{closed} (\xrefL{3.2.4}), so the real engine of the problem is a
single structural fact---that $E$ is closed---reused three times.

Believe the picture first. Pinning the first digit to $\{4,7\}$ confines $x$ to two short intervals
inside $[0.4,0.8]$; pinning the first two digits splits each into two shorter ones, and so on. What you
are watching is a base-ten Cantor set: at stage $m$ you keep $2^m$ intervals of length $10^{-m}$ and
discard everything between them, and $E$ is what survives forever. That image makes every verdict
visible---it is a dust of uncountably many points, it sits entirely above $4/9$ so it cannot be dense,
and it has no isolated specks because each surviving interval always contains two of the next stage's.

The four verdicts then come out in order, each leaning on a named result:
\begin{enumerate}[leftmargin=1.7em,itemsep=3pt,topsep=2pt,label=\textnormal{(\arabic*)}]
  \item the digit-to-point map is a bijection onto $E$, the one delicate point being that distinct
        sequences give distinct numbers---a digit gap of $3$ at the first disagreement outweighs the
        entire tail---so $E$ inherits uncountability from $\{4,7\}^{\mathbb N}$
        (\factref{binary-uncountable}{\xrefL{2.1.10}});
  \item density fails for a concrete reason, not by vague smallness: every point of $E$ is at least
        $0.4\overline 4=4/9$, so the whole interval $(0,4/9)$ is missed and points like $1/9$ have a
        neighbourhood disjoint from $E$, whence $\overline E\ne[0,1]$;
  \item closedness is the crux, and the move is to write $E$ as the nested intersection
        $\bigcap_m F_m$ of the stage-$m$ unions of closed intervals---an intersection of closed sets is
        closed (\factref{inter-closed}{\xrefL{3.2.10}}), and the equality holds because lying in every
        $F_m$ forces every finite initial digit-block into $\{4,7\}$, which is exactly membership in $E$;
  \item compactness is then immediate from Heine--Borel, since $E$ is closed and bounded in $\mathbb R$
        (\factref{heine-borel}{\xrefL{4.3.5}}), and perfectness adds only the limit-point check: flip a
        far-out digit to produce a distinct point of $E$ within $3\cdot 10^{-m}$, so every point is a
        cluster point (\xrefL{3.2.1}).
\end{enumerate}

Two checks keep the argument honest. ``Not dense'' is not the same as ``small''---$E$ is uncountable yet
fails density, because density is about \emph{reaching every neighbourhood}, not about cardinality; the
gap below $4/9$ is what kills it. And perfect demands \emph{both} halves: closed and no isolated points.
Closedness alone is not perfect (a single point is closed), and ``every point a limit point'' alone is
not perfect either (the open interval $(0,1)$ qualifies but is not closed); $E$ passes both, the
digit-flip supplying nearby companions and the intersection presentation supplying closedness, so no step
is circular.

The transferable move is the digit-restriction construction itself. Whenever a set is carved out by
forcing each entry of an infinite expansion into a fixed finite menu, read off its properties from the
menu: two-or-more choices per slot make it uncountable, a forbidden leading value makes it miss an
interval, and the stage-by-stage nested-interval picture makes it a closed---and in fact perfect---Cantor-type
set. The very same recipe, with zeros wedged between the free digits, builds a perfect set of irrationals
in \hyperlink{prob:18}{Problem~18}.
\end{Reflection}

% ------------------------------------------------------------
% ------------------------------------------------------------
\prob{18}{Rudin, Chapter 2, Exercise 18}
Is there a nonempty perfect set in $\mathbb R^{1}$ which contains no rational number?

\begin{SolutionBlue}
Yes. The verdict is built from a single device: take the digit-restriction trick of
\hyperlink{prob:17}{Problem~17}, which already yields a perfect set, and space the free digits out so
that ever-longer blocks of forced zeros sit between them---those zero blocks are what bar the rationals.
Write $T_n=\tfrac{n(n+1)}2$ for the $n$-th triangular number, so $T_1,T_2,T_3,\dots=1,3,6,10,\dots$, and
put
\[
P=\Bigl\{\textstyle\sum_{n=1}^{\infty}a_n\,10^{-T_n}\;:\;a_n\in\{4,7\}\ \text{for every }n\Bigr\}.
\]
A point of $P$ is the real number whose decimal expansion carries a $4$ or a $7$ in each triangular place
$T_n$ and a $0$ in every other place. The series converges absolutely (its terms are dominated by
$7\cdot 10^{-T_n}$), so $P$ is a well-defined nonempty subset of $[0,1]$; we show it is perfect---closed
with no isolated point---and disjoint from $\mathbb Q$.

$P$ is closed. For each $m\ge 1$ and each choice $(a_1,\dots,a_m)\in\{4,7\}^m$ of the first $m$ free
digits, the points of $P$ whose leading free digits are exactly $a_1,\dots,a_m$ all share the partial sum
$s=\sum_{n=1}^{m}a_n10^{-T_n}$, and the unwritten tail $\sum_{n>m}a_n10^{-T_n}$ is a nonnegative number at
most $\sum_{k>T_m}9\cdot10^{-k}=10^{-T_m}$; so each such point lies in the closed interval
$[\,s,\;s+10^{-T_m}\,]$. Let $C_m$ be the union of these $2^m$ closed intervals, one per choice of
$(a_1,\dots,a_m)$. Each $C_m$ is a finite union of closed sets, hence closed (\xrefL{3.2.10};
\factref{inter-closed}{a finite union of closed sets is closed}). Now $P=\bigcap_{m\ge 1}C_m$: every point
of $P$ lies in every $C_m$ by the bound just given, while a real number in all the $C_m$ has, for each
$m$, its first $m$ free digits drawn from $\{4,7\}$ and zeros elsewhere, so all its free digits are in
$\{4,7\}$ and it belongs to $P$. An arbitrary intersection of closed sets is closed (\xrefL{3.2.10};
\factref{inter-closed}{any intersection of closed sets is closed}), so $P$ is closed.

Every point of $P$ is a cluster point of $P$. Fix $x=\sum_{n\ge1}a_n10^{-T_n}\in P$ and any radius
$r>0$. Choose $m$ with $3\cdot 10^{-T_m}<r$, and let $x_m$ be the point of $P$ obtained by flipping the
single digit $a_m$ ($4\leftrightarrow 7$) and leaving every other $a_n$ unchanged. Then $x_m\in P$,
$x_m\neq x$, and
\[
|x_m-x|=|a_m-a'_m|\,10^{-T_m}=3\cdot 10^{-T_m}<r,
\]
so the ball $(x-r,x+r)$ contains a point of $P$ other than $x$. Thus $x$ is a cluster point of $P$
(\xrefL{3.2.1}). With $P$ closed and every point a cluster point, $P$ is perfect; and $P\neq\varnothing$
(for instance $\sum_n 4\cdot10^{-T_n}\in P$).

No point of $P$ is rational. A rational has an eventually periodic decimal expansion: dividing by a
denominator $q$ leaves only $q$ possible remainders, so the remainders---and hence the quotient digits---
must eventually cycle. Suppose $x\in P$ were rational, with its expansion eventually periodic of period
$\ell$ from some place onward. The gaps between consecutive triangular places grow without bound,
$T_{n+1}-T_n=n+1\to\infty$, so beyond any point of the expansion there is a run of more than $\ell$
consecutive zeros (the stretch strictly between two far-apart free places). A periodic block that contains
$\ell$ consecutive zeros is the all-zero block, forcing the expansion to be eventually $0$; but $P$ has a
nonzero digit $4$ or $7$ in every triangular place $T_n$, arbitrarily far out. This contradiction shows
$x$ is irrational. Hence $P$ is a nonempty perfect set containing no rational number.
\end{SolutionBlue}

\begin{Reflection}
The word \emph{perfect} is the whole instruction set. In our vocabulary it unpacks into two demands that
must be met together: the set is closed---it owns every one of its cluster points (\xrefL{3.2.4})---and it
has no isolated point, so every point is itself a cluster point (\xrefL{3.2.1}, \xrefL{3.2.2}). The extra
clause ``contains no rational'' is a number-theoretic constraint laid on top, and the tension between the
two is what makes the problem more than a definition-chase: rationals are \emph{dense}
(\factref{q-dense}{$\mathbb Q$ is dense in $\mathbb R$}, \xrefL{1.7.2}), so any interval is riddled with
them, which is exactly why a perfect set with room to spread---had it contained an interval---could never
dodge them. The set must be perfect yet nowhere interval-like, a Cantor-type dust.

That points straight at the construction from \hyperlink{prob:17}{Problem~17}, where restricting decimal
digits to $\{4,7\}$ produced a perfect set with empty interior. Believe that picture first: prescribing a
digit at each place and giving it two free choices builds a set that, at every point, has a near-twin got
by flipping one far-out digit---so no point stands alone---while leaving gaps no interval can fit through.
The single new idea here is to thin the free places out, parking them at the triangular positions
$T_n=\tfrac{n(n+1)}2$ so that the runs of forced zeros between them lengthen without bound. Those growing
zero runs are not decoration; they are the engine that defeats rationality.

The proof has three independent obligations, each closed by a named result or by the step before it:
\begin{enumerate}[leftmargin=1.7em,itemsep=3pt,topsep=2pt,label=\textnormal{(\arabic*)}]
  \item closedness comes from exhibiting $P$ as an intersection $\bigcap_m C_m$, where $C_m$ is the finite
        union of the $2^m$ closed ``decimal intervals'' pinned down by the first $m$ free digits; a finite
        union of closed sets is closed and so is an arbitrary intersection of them
        (\xrefL{3.2.10}; \factref{inter-closed}{intersections of closed sets are closed}), and the
        bookkeeping that makes $P=\bigcap_m C_m$ exact is that membership in every $C_m$ is the same as
        having all free digits in $\{4,7\}$;
  \item that every point is a cluster point is the digit-flip move: flipping the $m$-th free digit perturbs
        the value by exactly $3\cdot 10^{-T_m}$, which the freedom in $m$ drives below any radius $r$, so
        every ball around $x$ catches a different point of $P$---the definition of a cluster point
        (\xrefL{3.2.1}), and equivalently the failure of $x$ to be isolated (\xrefL{3.2.2});
  \item irrationality rests on the one fact about rationals worth quoting here---their decimals are
        eventually periodic, because long division by $q$ admits only $q$ remainders---played against the
        widening zero gaps: past any place there sits a zero run longer than the alleged period, a period
        meeting a full zero run is all zeros, yet $P$ keeps planting a $4$ or $7$ arbitrarily far out, and
        the two cannot coexist.
\end{enumerate}

The rigor turns on where the triangular spacing earns its keep, and it is worth checking by asking what a
\emph{constant} spacing would cost. Problem~17's set, with a free digit in every place, is perfect too, but
it is stuffed with rationals---$0.4444\dots=\tfrac49$ for one---because a fixed periodic digit pattern is
itself a periodic decimal. Only by letting the gaps grow, $T_{n+1}-T_n=n+1\to\infty$, do we guarantee a
zero run beating \emph{any} prescribed period, which is precisely the lever the eventual-periodicity test
needs; drop the growth and the argument collapses. Closedness and the no-isolated-point property, by
contrast, never used the spacing at all---they survive verbatim from Problem~17---so the design cleanly
separates ``perfect'' (handled by the digit menu) from ``no rationals'' (handled by the spacing). Nothing
is circular: closedness is built from the closed-set algebra, the cluster condition from an explicit
nearby point, and irrationality from a property of $\mathbb Q$ alone.

The transferable lesson is how decimal (or base-$b$) digit menus turn topological and arithmetic wishes
into independent dials. A two-symbol alphabet at chosen places forces perfection and uncountability
(\factref{binary-uncountable}{two-valued digit sequences are uncountable}, \xrefL{2.1.10}); spreading those
places out destroys eventual periodicity and so banishes the rationals; thinning them further still
controls measure and dimension. When a problem asks for a set that is large in one sense (perfect,
uncountable) yet evasive in another (rational-free, interval-free), reach for a Cantor-style digit
construction and tune the placement of the free digits to whatever the second condition forbids.
\end{Reflection}

% ------------------------------------------------------------
% ------------------------------------------------------------
\prob{19}{Rudin, Chapter 2, Exercise 19}
Call two subsets $A,B$ of a metric space $X$ \emph{separated} if $A\cap\overline{B}=\varnothing$ and
$\overline{A}\cap B=\varnothing$ --- neither meets the closure of the other.
\begin{enumerate}[leftmargin=1.7em,itemsep=2pt,topsep=2pt,label=\textnormal{(\alph*)}]
  \item If $A$ and $B$ are disjoint closed sets in some metric space $X$, prove that they are
        separated.
  \item Prove the same for disjoint open sets.
  \item Fix $p\in X$, $\delta>0$, define $A$ to be the set of all $q\in X$ for which $d(p,q)<\delta$,
        define $B$ similarly, with $>$ in place of $<$. Prove that $A$ and $B$ are separated.
  \item Prove that every connected metric space with at least two points is uncountable.
        \emph{Hint:} Use (c).
\end{enumerate}

\begin{SolutionBlue}
Throughout, $\overline{E}=E\cup E'$ is the closure (\xrefL{4.4.1}), so a point lies in $\overline{E}$
exactly when it is in $E$ or is a cluster point of $E$ --- equivalently, when every ball about it meets
$E$ (\xrefL{3.2.1}).

\smallskip\noindent\textit{(a)}\enspace Let $A,B$ be disjoint and closed. A closed set equals its own
closure (\xrefL{3.2.4}), so $\overline{A}=A$ and $\overline{B}=B$. Then
$A\cap\overline{B}=A\cap B=\varnothing$ and $\overline{A}\cap B=A\cap B=\varnothing$, so $A$ and $B$ are
separated.

\smallskip\noindent\textit{(b)}\enspace Let $A,B$ be disjoint and open. Take any $x\in A$. Since $A$ is
open, some ball $N_{r}(x)\subseteq A$, and $A\cap B=\varnothing$ forces $N_{r}(x)\cap B=\varnothing$; a
ball about $x$ thus misses $B$, so $x$ is neither in $B$ nor a cluster point of $B$, i.e.\
$x\notin\overline{B}$. As $x\in A$ was arbitrary, $A\cap\overline{B}=\varnothing$. Interchanging the
roles of $A$ and $B$ (both are open) gives $\overline{A}\cap B=\varnothing$. So $A$ and $B$ are
separated.

\smallskip\noindent\textit{(c)}\enspace With
\[
A=\{q\in X:d(p,q)<\delta\},\qquad B=\{q\in X:d(p,q)>\delta\},
\]
fix $q\in A$ and set $r=\delta-d(p,q)>0$. If $d(q,y)<r$, the triangle inequality (\xrefL{2.2.1}) gives
\[
d(p,y)\le d(p,q)+d(q,y)<d(p,q)+r=\delta,
\]
so $N_{r}(q)\subseteq A$ and meets $B$ nowhere; hence $q\notin\overline{B}$, and
$A\cap\overline{B}=\varnothing$. Symmetrically, fix $q\in B$ and set $r=d(p,q)-\delta>0$. If
$d(q,y)<r$, the triangle inequality gives
\[
d(p,y)\ge d(p,q)-d(q,y)>d(p,q)-r=\delta,
\]
so $N_{r}(q)\subseteq B$ misses $A$, whence $q\notin\overline{A}$ and $\overline{A}\cap B=\varnothing$.
Thus $A$ and $B$ are separated.

\smallskip\noindent\textit{(d)}\enspace Let $X$ be connected with two distinct points $p,q$, and put
$D=d(p,q)>0$. For each $t\in(0,D)$ set
\[
A_{t}=\{x\in X:d(p,x)<t\},\qquad B_{t}=\{x\in X:d(p,x)>t\}.
\]
The set $A_{t}=N_{t}(p)$ is an open ball, hence open (\xrefL{2.3.5}; \factref{balls-open}{open balls
are open}); and $B_{t}$ is open by the argument of (c), since each of its points carries a ball staying
in $B_{t}$. The two are disjoint, and both are nonempty: $p\in A_{t}$ (as $0<t$), while $q\in B_{t}$
(as $d(p,q)=D>t$).

I claim some $x\in X$ has $d(p,x)=t$. If not, then no point lands on the sphere $d(p,x)=t$, so every
point of $X$ falls into $A_{t}$ or $B_{t}$, giving $X=A_{t}\cup B_{t}$ as a union of two nonempty,
disjoint, open sets --- which is precisely a disconnection of $X$ (\xrefL{4.4.2}), against the
hypothesis that $X$ is connected. So a point $x_{t}$ with $d(p,x_{t})=t$ exists.

The assignment $t\mapsto x_{t}$ is injective: $d(p,x_{t})=t$ recovers $t$ from $x_{t}$, so distinct
parameters give distinct points. Hence $(0,D)$ injects into $X$. The interval $(0,D)$ is uncountable
--- the map $s\mapsto Ds$ carries $(0,1)$ bijectively onto it, and $(0,1)$ is uncountable
(\factref{r-uncountable}{$\mathbb R$ is uncountable}; \xrefL{2.1.12}). An injection of an uncountable
set into $X$ forbids $X$ from being countable, so $X$ is uncountable.
\end{SolutionBlue}

\begin{Reflection}
The load-bearing word is \emph{separated}, and the statement is generous enough to hand you its
meaning: $A$ and $B$ are separated when $A\cap\overline{B}=\varnothing$ and $\overline{A}\cap
B=\varnothing$. That this is strictly stronger than disjoint is the thing to register first --- it
forbids each set even from \emph{touching} the other's closure, so the entire problem is about the
closure $\overline{E}=E\cup E'$ (\xrefL{4.4.1}) and its neighbourhood test: $x\in\overline{E}$ exactly
when every ball about $x$ meets $E$ (\xrefL{3.2.1}). Once you read ``separated'' as ``a ball walls each
point of one set off from the other,'' parts (a)--(c) become three settings in which such a wall is
free, and the closure is the only machinery any of them touches.

The cleanest picture is part (c): the open ball $d(p,x)<\delta$ and its far exterior $d(p,x)>\delta$
sit on opposite sides of the sphere $d(p,x)=\delta$, and the sphere itself is the buffer that keeps
either set off the other's closure. The shell of thickness given by ``how far $q$ is from the
threshold'' is exactly the radius that fits a ball staying on $q$'s own side.

Each part rests on a single named fact, applied to closures:
\begin{enumerate}[leftmargin=1.7em,itemsep=3pt,topsep=2pt,label=\textnormal{(\arabic*)}]
  \item closed disjoint sets are separated because a closed set is its own closure (\xrefL{3.2.4}), so
        $A\cap\overline{B}$ and $\overline{A}\cap B$ collapse to the given $A\cap B=\varnothing$ ---
        nothing to prove beyond reading off the definition;
  \item disjoint open sets are separated because openness gives each $x\in A$ a ball $N_{r}(x)\subseteq
        A$, which is disjoint from $B$, so $x$ flunks the neighbourhood test for $\overline{B}$ and
        $A\cap\overline{B}=\varnothing$; the same with $A,B$ swapped finishes it;
  \item the two sides of a metric sphere are separated by a quantitative version of the same move:
        from $q\in A$, the slack $r=\delta-d(p,q)$ feeds the triangle inequality (\xrefL{2.2.1}) to put
        $N_{r}(q)$ wholly inside $A$, hence off $\overline{B}$; the mirror estimate
        $d(p,y)\ge d(p,q)-d(q,y)$ handles $q\in B$;
  \item connectedness turns (c) into a counting bound: for each threshold $t\in(0,D)$ the open ball
        $A_{t}$ (open by \factref{balls-open}{open balls are open}, \xrefL{2.3.5}) and the open
        exterior $B_{t}$ are nonempty and disjoint, so if no point sat on the sphere $d(p,x)=t$ they
        would disconnect $X$ (\xrefL{4.4.2}); connectedness therefore forces a witness $x_{t}$ at every
        radius, and $t\mapsto x_{t}$ injects the uncountable interval $(0,D)$ into $X$.
\end{enumerate}

The rigor of (d) turns on two checkpoints. First, $B_{t}$ must genuinely be open, and that is not a
freebie --- it is exactly part (c) reused, which is why the problem stages (c) before (d). Second, the
injection has to be honest: $d(p,x_{t})=t$ recovers $t$, so distinct radii give distinct points, and
$(0,D)$ is uncountable as an affine copy of $(0,1)$ (\factref{r-uncountable}{$\mathbb R$ is
uncountable}; \xrefL{2.1.12}). Both hypotheses in (d) are load-bearing: drop ``at least two points''
and $D=0$ leaves no interval to inject; drop ``connected'' and the conclusion fails outright, since
$\mathbb Q$ as a metric space is countable yet has at least two points --- it is its
\emph{disconnectedness}, a missing irrational threshold splitting it at every $t$, that lets it stay
small.

The move to carry away is that connectedness is a no-gaps condition you exploit by manufacturing a
candidate gap and daring the space to have it. Any clopen splitting is forbidden, so each nonempty
disjoint pair of open sets you can build must fail to cover $X$, and the points it misses are forced to
exist; sweeping the threshold across an interval converts ``the space has no gaps'' into ``the space
has at least a continuum of points.'' The same template --- present a would-be disconnection, let
connectedness veto it, and harvest the witness --- is the standard way the clopen definition
(\xrefL{4.4.2}) does real work.
\end{Reflection}

% ------------------------------------------------------------
% ------------------------------------------------------------
\prob{20}{Rudin, Chapter 2, Exercise 20}
Are closures and interiors of connected sets always connected? (Look at subsets of $\mathbb R^2$.)

\begin{SolutionBlue}
Closures of connected sets are always connected; interiors need not be.

For the closure, let $E$ be connected and suppose, toward a contradiction, that $\overline E=A\cup B$
with $A,B$ nonempty and \emph{separated}, meaning $A\cap\overline B=\varnothing$ and
$\overline A\cap B=\varnothing$. Intersecting with $E$ gives $E=(E\cap A)\cup(E\cap B)$, and these two
pieces are separated, since a subset of a separated set stays separated: $E\cap A\subseteq A$ and
$E\cap B\subseteq B$, so $(E\cap A)\cap\overline{E\cap B}\subseteq A\cap\overline B=\varnothing$ and
likewise on the other side. Because $E$ is connected, one piece is empty; say $E\cap A=\varnothing$,
so $E\subseteq B$ and hence $\overline E\subseteq\overline B$. But $A\subseteq\overline E$, so
$A\subseteq\overline B$, forcing $A=A\cap\overline B=\varnothing$ against $A$ nonempty. The case
$E\cap B=\varnothing$ is identical. Thus no such separation exists and $\overline E$ is connected.

For interiors, one counterexample suffices. In $\mathbb R^2$ take the two closed unit disks tangent at
the origin,
\[
  E=\{(x,y):(x+1)^2+y^2\le 1\}\cup\{(x,y):(x-1)^2+y^2\le 1\}.
\]
Each disk is convex, hence connected, and the two disks share the point $(0,0)$, so their union $E$ is
connected (a union of two connected sets with a common point cannot be separated, the same argument as
for the closure). Its interior, however, is the union of the two \emph{open} disks,
\[
  E^{\circ}=\{(x,y):(x+1)^2+y^2< 1\}\cup\{(x,y):(x-1)^2+y^2< 1\},
\]
because the tangency point $(0,0)$ lies on both bounding circles and so is not interior to either disk.
These two open disks are nonempty, disjoint, and open, hence separated: a disjoint open set cannot meet
the closure of another, since each of its points has a neighbourhood (the set itself) missing the other
(this is exactly \hyperlink{prob:19}{Problem~19}(b)). Therefore $E^{\circ}$ is disconnected, and
interiors of connected sets need not be connected.
\end{SolutionBlue}

\begin{Reflection}
The phrase \emph{connected set} fixes the entire toolkit. Connectedness is defined negatively---a set is
connected when it admits no splitting into two nonempty separated pieces (\xrefL{4.4.2})---and a
negative definition is a standing invitation to argue by contradiction (\xrefL{0.5.5}): assume the
separation you are forbidden and break it. The other two words, \emph{closure} and \emph{interior}, name
the operations $\overline E=E\cup E'$ and $E^{\circ}$ from \xrefL{4.4.1}, and the bracketed steer ``look
at subsets of $\mathbb R^2$'' is itself a clue: in $\mathbb R$ the connected sets are exactly the
intervals and both operations keep you among intervals, so any honest counterexample must live in a
space with more room---the plane.

Before proving anything, picture why the two halves should differ. Adding boundary points to a connected
blob cannot tear it apart---the new points cling to the set they came from---so the closure should stay
in one piece. But carving out the boundary can delete the very point that fused two lobes together: two
disks kissing at a single point form one connected set, yet erasing that shared point to pass to the
interior leaves two disjoint open disks drifting free. That single picture is the whole content of the
problem.

The closure half runs through four linked steps:
\begin{enumerate}[leftmargin=1.7em,itemsep=3pt,topsep=2pt,label=\textnormal{(\arabic*)}]
  \item assume a forbidden separation $\overline E=A\cup B$ with $A,B$ nonempty separated, which is the
        negation of ``$\overline E$ connected'' (\xrefL{4.4.2}) and the only foothold a negative
        hypothesis offers;
  \item restrict it to $E$, writing $E=(E\cap A)\cup(E\cap B)$, and check the pieces stay separated
        because a subset of a separated set is separated---shrinking a set only shrinks its closure, so
        the empty intersections survive;
  \item invoke the connectedness of $E$ itself to kill one piece, say $E\cap A=\varnothing$, which pins
        $E\subseteq B$ and therefore $\overline E\subseteq\overline B$ since closure is monotone
        (\xrefL{4.4.1});
  \item spring the trap: $A\subseteq\overline E\subseteq\overline B$ collides with the separation
        condition $A\cap\overline B=\varnothing$, forcing the nonempty $A$ to be empty.
\end{enumerate}
The hinge is step (2)'s monotonicity together with the defining asymmetry of separation---separated is
strictly more than disjoint, demanding $A\cap\overline B=\varnothing$, and it is precisely that closure
appearing in the condition that the argument exploits.

The interior half is a single well-chosen example, and its rigor lives in two checks. First, $E$ really
is connected: two convex (hence connected) disks meeting at $(0,0)$ join into one piece, again because
two connected sets sharing a point admit no separation. Second, $E^{\circ}$ really is disconnected: the
tangency point sits on both circles, so it is interior to neither disk, and the interior collapses to two
disjoint open disks; disjoint open sets are separated, which is \hyperlink{prob:19}{Problem~19}(b), so the
splitting is genuine. The example is sharp---move the disks apart and $E$ disconnects; overlap them and
$E^{\circ}$ reconnects; tangency is the knife's edge where the closure stays whole but the interior does
not. And nothing is circular: the closure direction never assumed what it proved, building the
contradiction only from monotonicity of closure and the meaning of separation.

The transferable move is to treat ``is this operation connectedness-preserving?'' as two opposite
questions. Adding closure points cannot create a separation, so closure (and in fact any set squeezed
between $E$ and $\overline E$) inherits connectedness; deleting points to form an interior can sever a
one-point bridge, so it does not. Whenever an operation only \emph{enlarges} toward the closure, suspect
permanence and chase a contradiction through monotonicity; whenever it can \emph{remove} a load-bearing
point, suspect failure and hunt in $\mathbb R^2$ for the tangency that the operation cuts.
\end{Reflection}

% ------------------------------------------------------------
% ------------------------------------------------------------
\prob{21}{Rudin, Chapter 2, Exercise 21}
Let $A$ and $B$ be separated subsets of some $\mathbb R^{k}$, suppose $\mathbf a\in A$, $\mathbf b\in B$,
and define
\[
\mathbf p(t)=(1-t)\mathbf a+t\mathbf b
\]
for $t\in\mathbb R^{1}$. Put $A_{0}=\mathbf p^{-1}(A)$, $B_{0}=\mathbf p^{-1}(B)$. [Thus $t\in A_{0}$ if
and only if $\mathbf p(t)\in A$.]
\begin{enumerate}
  \item Prove that $A_{0}$ and $B_{0}$ are separated subsets of $\mathbb R^{1}$.
  \item Prove that there exists $t_{0}\in(0,1)$ such that $\mathbf p(t_{0})\notin A\cup B$.
  \item Prove that every convex subset of $\mathbb R^{k}$ is connected.
\end{enumerate}

\begin{SolutionBlue}
Two sets $S,T$ are \emph{separated} when neither meets the other's closure: $\overline S\cap T=\varnothing$
and $S\cap\overline T=\varnothing$. The single fact that drives every part is that $\mathbf p$ moves
distances by a fixed factor: for all $s,t\in\mathbb R$,
\[
|\mathbf p(s)-\mathbf p(t)|=\bigl|(s-t)(\mathbf b-\mathbf a)\bigr|=|s-t|\,|\mathbf b-\mathbf a|,
\]
so $t_{n}\to t$ in $\mathbb R$ forces $\mathbf p(t_{n})\to\mathbf p(t)$ in $\mathbb R^{k}$.

(a) The claim reduces to one transfer: if $t\in\overline{A_{0}}$ then $\mathbf p(t)\in\overline A$. Take
$t\in\overline{A_{0}}=A_{0}\cup A_{0}'$. If $t\in A_{0}$ then $\mathbf p(t)\in A\subseteq\overline A$. If $t$
is a cluster point of $A_{0}$, some sequence $(t_{n})$ in $A_{0}$ has $t_{n}\to t$
(\factref{cluster-seq}{a cluster point is a sequential limit}); then $\mathbf p(t_{n})\in A$ for every $n$
and $\mathbf p(t_{n})\to\mathbf p(t)$, so $\mathbf p(t)$ is a cluster point of $A$ (or already lies in $A$),
hence $\mathbf p(t)\in\overline A$. The same argument gives $t\in\overline{B_{0}}\Rightarrow\mathbf
p(t)\in\overline B$.

Now suppose $t\in\overline{A_{0}}\cap B_{0}$. Then $\mathbf p(t)\in\overline A$ by the transfer, and
$\mathbf p(t)\in B$ since $t\in B_{0}$, so $\mathbf p(t)\in\overline A\cap B=\varnothing$ as $A,B$ are
separated --- impossible. Hence $\overline{A_{0}}\cap B_{0}=\varnothing$, and interchanging the roles of
$A$ and $B$ gives $A_{0}\cap\overline{B_{0}}=\varnothing$. Thus $A_{0}$ and $B_{0}$ are separated.

(b) Endpoints land where they should: $\mathbf p(0)=\mathbf a\in A$ and $\mathbf p(1)=\mathbf b\in B$, so
$0\in A_{0}$ and $1\in B_{0}$. Suppose, toward a contradiction, that $\mathbf p(t)\in A\cup B$ for every
$t\in[0,1]$. Writing $P=[0,1]\cap A_{0}$ and $Q=[0,1]\cap B_{0}$, this assumption says $[0,1]=P\cup Q$;
moreover $0\in P$ and $1\in Q$, so both are nonempty, they are disjoint, and they are separated, being
subsets of the separated sets $A_{0},B_{0}$ from part (a). I show no such split of $[0,1]$ exists.

Let $c=\sup P$. The set $P$ is nonempty (it contains $0$) and bounded above by $1$, so $c\in[0,1]$ exists
by the least-upper-bound property (\factref{lub}{the supremum property}), and being a supremum, $c\in\overline P$.
Separation forbids $c\in Q$ (else $c\in\overline P\cap Q=\varnothing$), so $c\in P$; in particular $c<1$,
since $1\in Q$. Every $t\in(c,1]$ exceeds $\sup P$, hence lies outside $P$, hence in $Q$; letting $t\downarrow c$
shows $c\in\overline Q$. Then $c\in P\cap\overline Q=\varnothing$, again impossible. This contradiction shows
some $t_{0}\in[0,1]$ has $\mathbf p(t_{0})\notin A\cup B$; and since $\mathbf p(0),\mathbf p(1)$ do lie in
$A\cup B$, that $t_{0}\in(0,1)$.

(c) Let $C\subseteq\mathbb R^{k}$ be convex, meaning $(1-t)\mathbf a+t\mathbf b\in C$ whenever $\mathbf
a,\mathbf b\in C$ and $0\le t\le 1$. Suppose, toward a contradiction, that $C$ is disconnected: $C=A\cup B$
with $A,B$ nonempty, separated (\xrefL{4.4.2}). Pick $\mathbf a\in A$, $\mathbf b\in B$. By convexity
$\mathbf p(t)\in C=A\cup B$ for all $t\in[0,1]$. But part (b), applied to these separated $A,B$, produces a
$t_{0}\in(0,1)$ with $\mathbf p(t_{0})\notin A\cup B=C$ --- contradicting $\mathbf p(t_{0})\in C$. Hence no
such split exists, and $C$ is connected.
\end{SolutionBlue}

\begin{Reflection}
The word that fixes the whole problem is \emph{separated}, and reading it correctly is the exercise.
``Disjoint'' would only forbid a shared point; ``separated'' is the stronger relation $\overline A\cap
B=A\cap\overline B=\varnothing$ that our notes use to define disconnectedness (\xrefL{4.4.2}), and a
relation phrased through closures is one you attack with closures. So the part-(a) target,
``$A_{0},B_{0}$ are separated,'' is really the claim that a closure on the line ($\overline{A_{0}}$) maps
into the matching closure upstairs ($\overline A$) --- which is exactly what the affine map $\mathbf p$,
being distance-scaling, delivers. The second clue is the explicit formula $\mathbf p(t)=(1-t)\mathbf
a+t\mathbf b$: a parametrised segment turns a geometric statement in $\mathbb R^{k}$ into a question about
the interval $[0,1]$ in $\mathbb R^{1}$, the place where ``connected'' is easiest to settle. The third is
the word \emph{convex} in part (c) --- convexity is precisely the promise that this segment never leaves
$C$, which is what lets the line problem report back to the set.

Believe it on a picture first. With $A=[0,1)$ and $B=(1,2]$ on the line, the gap at the coordinate $1$
belongs to neither set, yet any segment running from a point of $A$ to a point of $B$ must pass through
it; so the segment leaves $A\cup B$ at some interior parameter. Convexity forbids the gap from lying
outside the set, so a convex set cannot be cut this way.

The argument runs in three linked stages:
\begin{enumerate}[leftmargin=1.7em,itemsep=3pt,topsep=2pt,label=\textnormal{(\arabic*)}]
  \item part (a) rests on the distance identity $|\mathbf p(s)-\mathbf p(t)|=|s-t|\,|\mathbf b-\mathbf a|$,
        which makes $\mathbf p$ carry convergent sequences to convergent sequences; reading $\overline{A_{0}}$
        as $A_{0}$ together with its cluster points (\xrefL{3.2.1}, \factref{closure-closed}{closure $=$ set $+$
        cluster points}) and pulling a sequence out of any cluster point (\factref{cluster-seq}{cluster points are
        sequential limits}), the image sequence lands in $A$ and converges, so $\mathbf p$ sends
        $\overline{A_{0}}$ into $\overline A$ --- and a point of $\overline{A_{0}}\cap B_{0}$ would put
        $\mathbf p(t)$ in $\overline A\cap B$, which separation has emptied;
  \item part (b) is the connectedness of $[0,1]$ in disguise: a split $[0,1]=P\cup Q$ into nonempty separated
        pieces cannot exist, and the supremum $c=\sup P$ of the piece holding $0$ exhibits the impossibility ---
        $c\in\overline P$ keeps $c$ out of $Q$, yet the values just above $c$ are forced into $Q$ so
        $c\in\overline Q$ keeps it out of $P$, leaving $c$ nowhere to go. This is the one place
        completeness enters, through the least-upper-bound property (\xrefL{1.6.1}, \factref{lub}{the
        supremum property}); the assumption that $\mathbf p$ stays inside $A\cup B$ is exactly what manufactures
        the forbidden split, so denying it yields the escape value $t_{0}$, and the endpoints pin $t_{0}$ to the
        open interval;
  \item part (c) is pure assembly: convexity holds the segment inside $C$, so any separation $C=A\cup B$ would
        feed parts (a) and (b) a segment they force to step outside $A\cup B=C$ --- a contradiction
        (\xrefL{0.5.5}), so $C$ admits no separation and is connected (\xrefL{4.4.2}).
\end{enumerate}

Each hypothesis earns its place, and dropping one breaks the chain. Without \emph{separated} (only
disjoint), part (a)'s closure transfer collapses --- $A=[0,1)$, $B=[1,2]$ are disjoint but $1\in\overline A\cap B$,
and the segment never escapes $A\cup B$, so the conclusion of (b) genuinely fails. Without
\emph{convex} in (c), the set $\{|\mathbf x|<1\}\cup\{|\mathbf x|>2\}$ in $\mathbb R^{k}$ splits into
separated pieces and is disconnected, because the segment between the two shells does leave the set.
Nothing is circular: I never assumed $[0,1]$ connected as a quoted theorem but built its connectedness from
the supremum, the same axiom the notes lean on throughout.

The move to keep is the transport principle. A distance-controlled map carries closures forward, hence
separations backward, so connectedness of a complicated set can be tested along its segments by exporting
the whole question to the one space where connectedness is transparent --- the interval, where the
least-upper-bound property settles it. That is why ``convex $\Rightarrow$ connected'' sits at the end of
the chapter (\xrefL{4.4.2}): it is the first payoff of having both completeness and the topology of
$\mathbb R^{k}$ in hand at once.
\end{Reflection}

% ------------------------------------------------------------
% ------------------------------------------------------------
\prob{22}{Rudin, Chapter 2, Exercise 22}
A metric space is called \emph{separable} if it contains a countable dense subset. Show that $\mathbb R^k$
is separable. \emph{Hint:} consider the set of points which have only rational coordinates.

\begin{Solution}
The set $D=\mathbb Q^{k}=\{(q_1,\dots,q_k):q_j\in\mathbb Q\}$ is a countable dense subset of $\mathbb R^k$,
so $\mathbb R^k$ is separable.

$D$ is countable. It is the $k$-fold product of the countable set $\mathbb Q$, and a finite product of
countable sets is countable (\factref{finite-product}{List of Facts}; $\mathbb Q$ countable by
\xrefL{2.1.7}).

$D$ is dense, meaning every point of $\mathbb R^k$ lies in $D$ or is a cluster point of $D$; equivalently,
every open ball meets $D$. Fix $x=(x_1,\dots,x_k)\in\mathbb R^k$ and $\eps>0$. Because $\mathbb Q$ is dense
in $\mathbb R$ (\factref{q-dense}{List of Facts}; \xrefL{1.7.2}), for each coordinate choose
$q_j\in\mathbb Q$ with $|q_j-x_j|<\eps/\sqrt k$. Then $q=(q_1,\dots,q_k)\in D$ and
\[
|q-x|=\Bigl(\sum_{j=1}^{k}(q_j-x_j)^2\Bigr)^{1/2}
       <\Bigl(\sum_{j=1}^{k}\frac{\eps^2}{k}\Bigr)^{1/2}
       =\Bigl(k\cdot\frac{\eps^2}{k}\Bigr)^{1/2}=\eps,
\]
so $q\in B_\eps(x)$. Hence every ball about every point of $\mathbb R^k$ contains a point of $D$, which is
exactly density (\xrefL{4.4.1}). With $D$ countable and dense, $\mathbb R^k$ is separable.
\end{Solution}

\begin{Reflection}
The definition is handed to you, so the reading is of the word \emph{separable} itself: it asks for a
countable set that approximates every point arbitrarily well, and ``approximates every point'' is the
density condition---closure equal to the whole space (\xrefL{4.4.1})---while ``countable'' is the counting
condition. The hint names the candidate, the rational grid $\mathbb Q^k$, and the only real question is why
it is the right one. It is right because it is built coordinatewise from $\mathbb Q$, which is both
countable (\xrefL{2.1.7}) and dense in $\mathbb R$ (\xrefL{1.7.2}); the two properties demanded of $D$ are
inherited from those two properties of $\mathbb Q$, one for counting and one for approximation.

Picture $k=2$: $\mathbb Q^2$ is the lattice of points with both coordinates rational, a countable web that
nonetheless threads arbitrarily close to every point of the plane. Believe that the web is fine enough to
enter any disc, and the proof is just making ``fine enough'' quantitative.

The verification splits cleanly along the definition:
\begin{enumerate}[leftmargin=1.7em,itemsep=3pt,topsep=2pt,label=\textnormal{(\arabic*)}]
  \item countability is bookkeeping---$\mathbb Q^k$ is a product of $k$ copies of the countable set
        $\mathbb Q$, and a finite product of countable sets is countable
        (\factref{finite-product}{the product fact}, resting on \xrefL{2.1.7}); the finiteness of $k$ is
        what keeps the product countable;
  \item density is where the geometry lives---to land a rational point inside the ball $B_\eps(x)$
        (\xrefL{2.3.1}), approximate each of the $k$ coordinates of $x$ by a rational within
        $\eps/\sqrt k$, which density of $\mathbb Q$ in $\mathbb R$ supplies
        (\factref{q-dense}{$\mathbb Q$ dense}; \xrefL{1.7.2});
  \item the per-coordinate tolerance $\eps/\sqrt k$ is then forced by the Euclidean norm: $k$ squared
        errors each below $\eps^2/k$ sum to less than $\eps^2$, so the distance is below $\eps$ and the
        rational point meets the ball, which is the definition of density (\xrefL{4.4.1}, via cluster
        points \xrefL{3.2.1}).
\end{enumerate}

The estimate is honest because the $\sqrt k$ is not a fudge but exactly the factor the norm contributes
from $k$ coordinates---split the tolerance any more coarsely and the sum of squares could exceed $\eps^2$.
The finiteness of $k$ is load-bearing twice over: it is what makes the product countable, and it is what
lets the single split $\eps/\sqrt k$ work for every coordinate at once. In an infinite-dimensional space
neither step survives unchanged, which is why separability there is a genuine question rather than a
formality.

The transferable move is the recipe for separability of a product: take the product of countable dense
sets in the factors, and split the tolerance across finitely many coordinates so the norm reassembles them
under $\eps$. The same rational grid is the workhorse behind the next exercise---a countable dense set is
precisely what lets you build a countable base by centering balls at its points (\hyperlink{prob:23}{Problem~23}).
\end{Reflection}

% ------------------------------------------------------------
% ------------------------------------------------------------
\prob{23}{Rudin, Chapter 2, Exercise 23}
A collection $\{V_{\alpha}\}$ of open subsets of $X$ is said to be a \emph{base} for $X$ if the
following is true: for every $x\in X$ and every open set $G\subset X$ such that $x\in G$, we have
$x\in V_{\alpha}\subset G$ for some $\alpha$. In other words, every open set in $X$ is the union of a
subcollection of $\{V_{\alpha}\}$.

Prove that every separable metric space has a \emph{countable} base. \emph{Hint:} take all
neighborhoods with rational radius and center in some countable dense subset of $X$.

\begin{Solution}
Let $X$ be separable, fix a countable dense subset $D=\{x_1,x_2,\dots\}$, and set
\[
  \mathcal B=\bigl\{\,B_r(x_i):x_i\in D,\ r\in\mathbb Q,\ r>0\,\bigr\}.
\]
Each member is an open ball, hence open, so $\mathcal B$ is a collection of open sets; the claim is
that it is a countable base.

It is countable. The map $(x_i,r)\mapsto B_r(x_i)$ sends $D\times\bigl(\mathbb Q\cap(0,\infty)\bigr)$
onto $\mathcal B$, and that product of two countable sets is countable, so $\mathcal B$ is at most
countable.

It is a base. Let $G\subseteq X$ be open and $x\in G$. Because $G$ is open, $x$ is an interior point,
so there is $R>0$ with $B_R(x)\subseteq G$. Density of $D$ gives a center $x_i\in D$ with
$d(x,x_i)<R/4$, and density of $\mathbb Q$ in $\mathbb R$ gives a rational radius $r$ with
\[
  d(x,x_i)<r<\tfrac R2 .
\]
Then $x\in B_r(x_i)$, since $d(x,x_i)<r$. Moreover $B_r(x_i)\subseteq B_R(x)$: for any $y\in
B_r(x_i)$ the triangle inequality gives
\[
  d(y,x)\le d(y,x_i)+d(x_i,x)<r+\tfrac R4<\tfrac R2+\tfrac R4=\tfrac{3R}{4}<R,
\]
so $y\in B_R(x)\subseteq G$. Thus $x\in B_r(x_i)\subseteq G$ with $B_r(x_i)\in\mathcal B$.

Since each $x\in G$ sits in some member of $\mathcal B$ contained in $G$, the set $G$ is the union of
those members; equivalently every open $G$ is a union of a subcollection of $\mathcal B$. Hence
$\mathcal B$ is a countable base for $X$.
\end{Solution}

\begin{Reflection}
The defining clue is the word \emph{base}: a supply of open sets out of which every open set can be
rebuilt by taking unions, so that to describe an open $G$ around a point $x$ it is enough to slot one
member of the supply between $x$ and $G$. Built into the very definition of open set
(\xrefL{2.3.3}) is that $x\in G$ open means some ball $B_R(x)\subseteq G$, so the open balls already
form a base---and the real demand here is the adjective \emph{countable}. The other clue is the
hypothesis \emph{separable}, which by definition hands you a countable dense subset $D$
(\hyperlink{prob:22}{Problem~22} is where $\mathbb R^k$ earns that label); ``dense'' is the promise
that points of $D$ sit arbitrarily close to every $x$, which is exactly what lets a ball centered in
$D$ still catch $x$.

Two independent infinities threaten countability---which centers, and which radii---and the cure is
to discretize both. Restrict the centers to the countable set $D$ and the radii to the positive
rationals, so the catalogue $\mathcal B$ is indexed by a product of two countable sets and is
therefore countable (\factref{nzq-countable}{$\mathbb Q$ is countable}, with
\factref{finite-product}{a finite product of countable sets is countable}). Before proving that this
thinned-out catalogue still works, picture why it should: $x$ may be irrational and the enclosing
radius $R$ may be irrational, yet a rational-radius ball pinned to a nearby dense center can be slid
into the gap, because density leaves room on both sides.

The fitting argument is a two-sided $\varepsilon$ estimate, and each move rests on something named:
\begin{enumerate}[leftmargin=1.7em,itemsep=3pt,topsep=2pt,label=\textnormal{(\arabic*)}]
  \item from $x\in G$ open, the definition of interior point (\xrefL{2.3.3}) produces a radius $R>0$
        with $B_R(x)\subseteq G$---the room inside $G$ that everything else must fit within;
  \item denseness of $D$ supplies a center $x_i\in D$ within $R/4$ of $x$, so the base ball can be
        anchored at a legitimate (countable) center near $x$ rather than at $x$ itself;
  \item density of $\mathbb Q$ in $\mathbb R$ (\factref{q-dense}{$\mathbb Q$ is dense}, \xrefL{1.7.2})
        supplies a rational radius $r$ with $d(x,x_i)<r<R/2$: the left inequality forces $x\in
        B_r(x_i)$, so the ball still catches $x$, and the right inequality reserves the slack the next
        step spends;
  \item the triangle inequality bounds any $y\in B_r(x_i)$ by $d(y,x)<r+R/4<3R/4<R$, placing
        $B_r(x_i)\subseteq B_R(x)\subseteq G$, so the chosen ball lies wholly inside $G$;
  \item with one member of $\mathcal B$ wedged between each $x$ and $G$, the open set $G$ is the union
        of those members, which is what ``base'' demands.
\end{enumerate}

The rigor turns on the margins, and it is worth seeing each one carry its weight. Centering at $x_i$
within $R/4$ and capping $r$ below $R/2$ leaves $d(y,x)<R/4+R/2<R$ with room to spare; had the radius
been allowed up to $R$ the containment $B_r(x_i)\subseteq B_R(x)$ could fail. The lower bound
$d(x,x_i)<r$ is equally load-bearing: drop it and the ball might miss $x$, so it would not witness
$x\in B_r(x_i)\subseteq G$. Both halves of ``base''---containing the point and sitting inside the open
set---are checked, and nothing is circular, since the catalogue is built first and only afterward
shown to suffice.

The transferable move is the double discretization: when a continuum of choices threatens countability,
replace each continuous parameter by a countable dense set of values and let a margin-splitting
estimate absorb the error. This is precisely the route from separable to ``second countable,'' and the
countable base it yields is the workhorse behind later reductions of open covers and sequential
arguments.
\end{Reflection}

% ------------------------------------------------------------
% ------------------------------------------------------------
\prob{24}{Rudin, Chapter 2, Exercise 24}
Let $X$ be a metric space in which every infinite subset has a limit point. Prove that $X$ is
separable. \emph{Hint:} Fix $\delta>0$, and pick $x_1\in X$. Having chosen $x_1,\dots,x_j\in X$, choose
$x_{j+1}\in X$, if possible, so that $d(x_i,x_{j+1})\ge\delta$ for $i=1,\dots,j$. Show that this
process must stop after a finite number of steps, and that $X$ can therefore be covered by finitely
many neighborhoods of radius $\delta$. Take $\delta=1/n$ $(n=1,2,3,\dots)$, and consider the centers of
the corresponding neighborhoods.

\begin{SolutionBlue}
Call a finite set $\delta$-separated if any two of its points are at distance $\ge\delta$. The whole
argument rests on one claim: under the hypothesis, for each $\delta>0$ there is no infinite
$\delta$-separated set, so $X$ admits a finite cover by balls of radius $\delta$.

Fix $\delta>0$ and run the greedy process: pick $x_1\in X$, and having chosen $x_1,\dots,x_j$, pick
$x_{j+1}$ with $d(x_i,x_{j+1})\ge\delta$ for all $i\le j$ whenever such a point exists. The points so
chosen form a $\delta$-separated set, and the process must terminate. If it did not, it would produce an
infinite $\delta$-separated set $E=\{x_1,x_2,\dots\}$; by hypothesis $E$ has a limit point $p\in X$, and
every ball about $p$ contains infinitely many points of $E$, in particular two distinct points
$x_i,x_k$ in the ball $B_{\delta/2}(p)$. The triangle inequality then gives
$d(x_i,x_k)\le d(x_i,p)+d(p,x_k)<\delta/2+\delta/2=\delta$, contradicting $\delta$-separation. So the
process halts, say after $x_1,\dots,x_m$.

Once it halts, the balls $B_\delta(x_1),\dots,B_\delta(x_m)$ cover $X$: if some $x\in X$ lay in none of
them, then $d(x_i,x)\ge\delta$ for every $i\le m$, so $x$ would be a legitimate next choice and the
process would not have stopped. Thus for each $\delta>0$ a \emph{finite} set of centers gives a cover of
$X$ by $\delta$-balls.

Apply this with $\delta=1/n$ for every $n\ge 1$: choose a finite set
$S_n=\{x^{(n)}_1,\dots,x^{(n)}_{m_n}\}\subseteq X$ with
$X=\bigcup_{j=1}^{m_n}B_{1/n}(x^{(n)}_j)$, and put $S=\bigcup_{n=1}^{\infty}S_n$. Being a countable union
of finite sets, $S$ is at most countable.

Finally $S$ is dense: given $x\in X$ and $\eps>0$, pick $n$ with $1/n<\eps$; the radius-$1/n$ balls
centered at $S_n$ cover $X$, so $x\in B_{1/n}(s)$ for some $s\in S_n\subseteq S$, whence
$d(x,s)<1/n<\eps$. Every ball about every point of $X$ therefore meets $S$, so $S$ is a countable dense
subset and $X$ is separable.
\end{SolutionBlue}

\begin{Reflection}
The lone hypothesis---``every infinite subset has a limit point''---is the entire engine, and the word
that unlocks it is \emph{limit point}. A cluster point is not merely approached; every ball around it
swallows infinitely many points of the set (\xrefL{3.2.1}; the quotable
\factref{cluster-infinite}{``every ball about a limit point contains infinitely many points''}). That
``infinitely many in an arbitrarily small ball'' is the lever: it says points cannot stay spread out
forever, which is exactly the obstruction the hint's construction is built to expose. Read the goal the
same way---``separable'' means a countable dense subset, and density is a statement about closures
(\xrefL{4.4.1}): a set whose closure is all of $X$, i.e.\ one that comes within every $\eps$ of every
point. So the target is a countable set of points fine enough to approximate all of $X$.

Believe it on a familiar space first. In $\mathbb R$ the rationals are dense and countable; the way you
build such a set by hand is to lay down a grid at each scale---points $1/n$ apart---and let the meshes
shrink. The construction here is that picture made intrinsic: at scale $\delta$ you greedily scatter
points that stay $\delta$ apart, and the hypothesis guarantees you run out of room after finitely many.

The proof is a short chain, each link resting on the one before:
\begin{enumerate}[leftmargin=1.7em,itemsep=3pt,topsep=2pt,label=\textnormal{(\arabic*)}]
  \item the greedy process at scale $\delta$ produces points pairwise at distance $\ge\delta$; if it
        never stopped it would yield an \emph{infinite} such set, and the hypothesis hands that set a
        limit point $p$;
  \item the cluster property (\xrefL{3.2.1}, \factref{cluster-infinite}{infinitely many in every
        ball}) crowds infinitely many of the points into $B_{\delta/2}(p)$, so two distinct ones sit in
        that half-radius ball, and the triangle inequality squeezes them to distance $<\delta$---against
        $\delta$-separation, so the process must halt;
  \item once it halts at $x_1,\dots,x_m$, those centers' $\delta$-balls cover $X$, since any uncovered
        point would itself have been a legal next pick---so finitely many $\delta$-balls cover $X$
        (\xrefL{2.3.1}, the definition of the ball $B_\delta(x_i)=\{x:d(x_i,x)<\delta\}$);
  \item running this at every scale $\delta=1/n$ gives finite center-sets $S_n$, and
        $S=\bigcup_n S_n$ is a countable union of finite sets, hence at most countable
        (\factref{countable-union}{a countable union of countable sets is countable}; \xrefL{2.1.9});
  \item $S$ is dense: for $x\in X$ and $\eps>0$, take $1/n<\eps$, and the cover at scale $1/n$ puts some
        $s\in S_n$ within $1/n<\eps$ of $x$.
\end{enumerate}

The rigor turns on two points worth checking. The radius $\delta/2$ in step (2) is not cosmetic: the
ball must be small enough that \emph{any} two of its points are closer than $\delta$, and only a radius
strictly below $\delta/2$---which $\delta/2$ achieves through the strict triangle inequality on distinct
points---forces the contradiction. And the hypothesis is genuinely load-bearing: drop it and the
conclusion fails, as an uncountable set with the discrete metric shows---every pair sits at distance
$1$, no infinite subset has a limit point, and no countable set can be dense because every ball of
radius $\tfrac12$ holds a single point. There is no circularity: the limit point in step (1) is supplied
by the hypothesis, never assumed of $S$.

The transferable move is the greedy maximal separated set, the metric-space cousin of a packing
argument: to manufacture a countable approximating skeleton, pack in points kept a fixed distance apart,
let a compactness-flavored hypothesis cap each packing at finite size, and shrink the spacing to $0$
along $1/n$. This is the same machine that drives the next problems in the set---Problem~25 reaches a
countable base from compactness by covering at every scale $1/n$, and Problem~26 turns this very
separability back into compactness---so the $1/n$-net is worth keeping in hand.
\end{Reflection}

% ------------------------------------------------------------
% ------------------------------------------------------------
\prob{25}{Rudin, Chapter 2, Exercise 25}
Prove that every compact metric space $K$ has a countable base, and that $K$ is therefore separable.
\emph{Hint:} For every positive integer $n$, there are finitely many neighbourhoods of radius $1/n$
whose union covers $K$.

\begin{Solution}
A \emph{base} for $K$ is a collection $\mathcal B$ of open sets such that every open $G\subseteq K$ is a
union of members of $\mathcal B$; equivalently, for every open $G$ and every $x\in G$ there is some
$B\in\mathcal B$ with $x\in B\subseteq G$. We produce a countable such collection, and a countable dense
set, simultaneously.

Fix $n\in\mathbb N$ with $n\ge 1$. The open balls $\{B_{1/n}(x):x\in K\}$ cover $K$, since each $x$ lies
in its own ball $B_{1/n}(x)$, and each is open (\xrefL{2.3.5}; \factref{balls-open}{open balls are
open}). By compactness this cover has a finite subcover (\xrefL{3.4.1}): there are centres
$x_{n,1},\dots,x_{n,m_n}\in K$ with
\[
  K=\bigcup_{j=1}^{m_n}B_{1/n}(x_{n,j}).
\]
Collect the balls over all scales and all the centres they generate,
\[
  \mathcal B=\bigl\{\,B_{1/n}(x_{n,j}):n\ge 1,\ 1\le j\le m_n\,\bigr\},
  \qquad
  D=\bigl\{\,x_{n,j}:n\ge 1,\ 1\le j\le m_n\,\bigr\}.
\]
Each is a countable union of finite sets, hence countable (\factref{countable-union}{a countable union
of countable sets is countable}).

The collection $\mathcal B$ is a base. Let $G\subseteq K$ be open and $x\in G$; by openness choose
$\eps>0$ with $B_\eps(x)\subseteq G$ (\xrefL{2.3.3}). Pick $n$ with $2/n<\eps$. The radius-$1/n$ balls
cover $K$, so $x\in B_{1/n}(x_{n,j})$ for some $j$. For every $y$ in that ball the triangle inequality
gives
\[
  d(y,x)\le d(y,x_{n,j})+d(x_{n,j},x)<\tfrac1n+\tfrac1n=\tfrac2n<\eps,
\]
so $B_{1/n}(x_{n,j})\subseteq B_\eps(x)\subseteq G$. Thus $x\in B_{1/n}(x_{n,j})\subseteq G$ with
$B_{1/n}(x_{n,j})\in\mathcal B$, and $\mathcal B$ is a countable base.

Finally $D$ is dense: given $x\in K$ and $r>0$, choose $n$ with $1/n<r$; the radius-$1/n$ balls cover
$K$, so $x\in B_{1/n}(x_{n,j})$ for some $j$, that is $d(x,x_{n,j})<1/n<r$. Every ball about every point
of $K$ therefore meets the countable set $D$, so $K$ is separable.
\end{Solution}

\begin{Reflection}
The words that set the strategy are \emph{compact} and \emph{base}. Compactness is a statement about
covers---every open cover admits a finite subcover (\xrefL{3.4.1})---so the way to mine it is to hand it
a cover and pocket the finite list it returns. A \emph{base} is a fixed supply of open sets out of which
every open set is rebuilt by unions (this is exactly the notion of \hyperlink{prob:23}{Problem~23}), and
``countable base'' is the demand that this supply be countable. The hint names the cover to feed in: at
each scale $1/n$, the balls of radius $1/n$. The reason that single family of radii is enough is the part
to internalise---the radii $1/n$ shrink to $0$, so once you can land inside any prescribed ball
$B_\eps(x)$, some scale $1/n$ is fine enough to fit a covering ball between $x$ and $B_\eps(x)$.

Believe it before proving it. Picture $K$ tiled, for each $n$, by finitely many balls of radius $1/n$;
as $n$ grows the tiles get finer, and the centres pile up densely. Any open neighbourhood $B_\eps(x)$,
no matter how small, is eventually coarser than some tiling: pick $n$ so that a radius-$1/n$ tile, being
half the width of $B_\eps(x)$, slides inside it. The countable heap of all these tiles is the base, and
their centres are the dense set.

The construction runs through four linked steps, each resting on the one before:
\begin{enumerate}[leftmargin=1.7em,itemsep=3pt,topsep=2pt,label=\textnormal{(\arabic*)}]
  \item at each scale $n$ the balls $\{B_{1/n}(x):x\in K\}$ form an open cover---open because balls are
        open (\xrefL{2.3.5}; \factref{balls-open}{open balls are open}), a cover because every point
        sits in its own ball---and compactness extracts a finite subcover $B_{1/n}(x_{n,1}),\dots,
        B_{1/n}(x_{n,m_n})$ (\xrefL{3.4.1}); this is the only place compactness is spent;
  \item gathering these finite lists over all $n$ gives the family $\mathcal B$ and the centre set $D$,
        each a countable union of finite sets, hence countable (\factref{countable-union}{a countable
        union of countable sets is countable});
  \item $\mathcal B$ is a base: inside a given $B_\eps(x)\subseteq G$, choose $n$ with $2/n<\eps$, take
        the covering ball $B_{1/n}(x_{n,j})$ that contains $x$, and the triangle inequality keeps that
        whole ball inside $B_\eps(x)$, since two radii $1/n$ sum to less than $\eps$---this is the
        two-ball estimate, and the slack $2/n<\eps$ is exactly what makes it close;
  \item $D$ is dense by the same covering balls read at scale $1/n<r$: the centre nearest $x$ at that
        scale lies within $1/n<r$, so every ball about $x$ meets $D$.
\end{enumerate}

The argument is honest about where each hypothesis lives. Compactness is load-bearing and used once, in
step (1): without it the scale-$1/n$ cover need not reduce to a finite list, the centre set $D$ can be
uncountable, and the base falls apart---an uncountable discrete space (every ball of radius $1/n<1$ is a
single point) is the standard space with no countable base and no countable dense set. The choice of
radius is also genuine, not cosmetic: a covering ball of radius $1/n$ with $1/n<\eps$ would only place
its \emph{centre} near $x$, but its far side could poke outside $B_\eps(x)$; demanding $2/n<\eps$ buys
the margin the triangle inequality needs. Nothing is circular---the base is built from balls produced by
compactness, never assumed---and the two conclusions share one engine, so separability costs no extra
work once the centres are in hand.

The transferable move is to read ``compact'' as ``hand it a cover, keep the finite subcover,'' and to
realise that running this at every scale $1/n$ converts the qualitative finiteness into a single
countable structure that approximates all open sets at once. This is the compact end of a chain the
homework builds deliberately: separable yields a countable base by \hyperlink{prob:23}{Problem~23}, and
here compactness yields the countable dense set outright, so compact $\Rightarrow$ separable
$\Rightarrow$ countable base, with compactness the strongest link that drives the whole thing.
\end{Reflection}

% ------------------------------------------------------------
% ------------------------------------------------------------
\prob{26}{Rudin, Chapter 2, Exercise 26}
Let $X$ be a metric space in which every infinite subset has a limit point. Prove that $X$ is compact.
\emph{Hint:} By Exercises 23 and 24, $X$ has a countable base. It follows that every open cover of $X$
has a \emph{countable} subcover $\{G_n\}$, $n=1,2,3,\ldots$. If no finite subcollection of $\{G_n\}$
covers $X$, then the complement $F_n$ of $G_1\cup\cdots\cup G_n$ is nonempty for each $n$, but $\bigcap
F_n$ is empty. If $E$ is a set which contains a point from each $F_n$, consider a limit point of $E$,
and obtain a contradiction.

\begin{SolutionBlue}
Call a metric space \emph{limit-point compact} if every infinite subset has a cluster point; this is the
standing hypothesis on $X$. The goal is genuine compactness: every open cover admits a finite subcover.

First reduce arbitrary covers to countable ones. By \hyperlink{prob:24}{Problem~24} the hypothesis makes
$X$ separable, and by \hyperlink{prob:23}{Problem~23} a separable metric space has a countable base
$\{V_k\}_{k\ge 1}$. Let $\{U_\alpha\}$ be any open cover of $X$. For each index $k$ such that some
$U_\alpha$ contains $V_k$, fix one such cover member and call it $U_{\alpha(k)}$. The chosen sets cover
$X$: given $x\in X$, pick $\alpha$ with $x\in U_\alpha$; since $\{V_k\}$ is a base there is a $V_k$ with
$x\in V_k\subseteq U_\alpha$, and the member $U_{\alpha(k)}$ chosen for that $k$ then contains $x$.
Indexing by a subset of the positive integers, $\{U_{\alpha(k)}\}$ is countable, so every open cover of
$X$ has a countable subcover. Enumerate one as $G_1,G_2,G_3,\ldots$.

It now suffices to show this countable cover has a finite subcover. Suppose, toward a contradiction, that
it does not. For each $n\ge 1$ set
\[
  F_n=X\setminus(G_1\cup\cdots\cup G_n).
\]
Each $F_n$ is the complement of a finite union of open sets, hence closed (\xrefL{3.2.8}), and each is
nonempty, since no finite collection $G_1,\ldots,G_n$ covers $X$. The sets descend,
$F_1\supseteq F_2\supseteq\cdots$, and
\[
  \bigcap_{n=1}^{\infty}F_n=X\setminus\bigcup_{n=1}^{\infty}G_n=\varnothing,
\]
because the $G_n$ cover $X$.

Choose a point $x_n\in F_n$ for each $n$ and put $E=\{x_n:n\ge 1\}$. The set $E$ is infinite. Were it
finite, some single point $p$ would equal $x_n$ for infinitely many $n$; as the $F_n$ are nested, $p$
would then lie in $F_n$ for arbitrarily large $n$, hence in every $F_n$, contradicting
$\bigcap_n F_n=\varnothing$.

Being infinite, $E$ has a cluster point $x\in X$ by hypothesis. Fix $N$ with $x\in G_N$, which exists
because the $G_n$ cover $X$. Since $G_N$ is open, it is a neighborhood of $x$, and a ball around a
cluster point of $E$ contains infinitely many points of $E$
(\xrefL{3.2.1}; \factref{cluster-infinite}{List of Facts}). In particular $G_N$ contains some $x_m$ with
$m\ge N$. But $x_m\in F_m\subseteq F_N=X\setminus(G_1\cup\cdots\cup G_N)$, so $x_m\notin G_N$---a
contradiction.

Therefore the countable cover $\{G_n\}$ has a finite subcover. Since every open cover of $X$ reduces to
such a countable cover, every open cover of $X$ has a finite subcover, and $X$ is compact (\xrefL{3.4.1}).
\end{SolutionBlue}

\begin{Reflection}
The phrase ``every infinite subset has a limit point'' is the diagnostic clue: it is the conclusion of
the Bolzano--Weierstrass theorem, \factref{infinite-cluster}{``an infinite subset of a compact space has
a cluster point''} (\xrefL{4.3.3}). The problem hands you that consequence as a hypothesis and asks you to
recover compactness, so it is asking whether the implication reverses. In $\mathbb R^k$ the answer is yes
by Heine--Borel, but here $X$ is an arbitrary metric space, so you cannot lean on boundedness or cells;
you must build a finite subcover out of nothing but the limit-point property itself.

The hint is really a map of the whole proof, and reading it teaches the strategy. ``$X$ has a countable
base'' points back through \hyperlink{prob:24}{Problem~24} (the same hypothesis yields separability) and
\hyperlink{prob:23}{Problem~23} (separable metric spaces are second countable); that base is what shrinks
any open cover to a countable one. The rest of the hint is the signature of a nested-closed-sets
argument---nonempty descending closed sets $F_n$ with empty intersection, one point chosen from each, and
a cluster point used to spring the trap. The cluster point is where the lone hypothesis finally gets spent.

Believe the shape before checking it. The $F_n$ are the parts of $X$ still uncovered after $n$ of the
$G_i$; if no finite stage ever covers everything, each $F_n$ has a leftover point, and stringing those
leftovers together produces an infinite set $E$ marching out toward the part of $X$ that the cover is
failing to reach. A cluster point of $E$ is a point the cover \emph{does} reach---and that is the
collision.

Each step rests on a named support:
\begin{enumerate}[leftmargin=1.7em,itemsep=3pt,topsep=2pt,label=\textnormal{(\arabic*)}]
  \item the countable base from \hyperlink{prob:23}{Problem~23} lets you pick, for each basic set that
        fits inside some cover member, a single such member; these cover $X$ exactly because a base
        rebuilds every open set, so an arbitrary cover collapses to a countable subcover $\{G_n\}$;
  \item assuming no finite subcover, the sets $F_n=X\setminus(G_1\cup\cdots\cup G_n)$ are nonempty (the
        failure hypothesis), closed as complements of finite unions of open sets (\xrefL{3.2.8}), nested,
        and with empty intersection precisely because the $G_n$ cover $X$;
  \item a transversal $E$ choosing one $x_n\in F_n$ is infinite, since a repeated value would, by
        nesting, sit in every $F_n$ and violate $\bigcap_n F_n=\varnothing$;
  \item the standing hypothesis gives $E$ a cluster point $x$, which lies in some $G_N$; openness makes
        $G_N$ a neighborhood, and a neighborhood of a cluster point holds infinitely many points of $E$
        (\xrefL{3.2.1}; \factref{cluster-infinite}{List of Facts}), so it catches some $x_m$ with
        $m\ge N$ even though $x_m\in F_m\subseteq F_N$ was built to avoid $G_N$.
\end{enumerate}

The rigor turns on two places where a weaker reading would fail. The contradiction needs $x_m$ with an
index $m\ge N$, not merely \emph{some} $x_m\in G_N$; that is exactly why ``infinitely many points of
$E$'' is the right tool rather than ``some point''---finitely many points could all carry indices below
$N$ and slip through. And the chosen $E$ must be genuinely infinite, or it would have no cluster point to
exploit; the nesting of the $F_n$ is what guarantees infinitude, and it is the same nesting that makes
$x_m\in F_m\subseteq F_N$ land where the contradiction needs it. The limit-point hypothesis is load-bearing
once: drop it and the construction still produces nested nonempty closed $F_n$ with empty intersection,
which is precisely the picture of a non-compact space such as $\mathbb Q$ or an open interval, so without
it nothing forces a finite subcover.

The transferable move is the conversion itself. To turn a sequential or limit-point property into the
covering definition of compactness, first buy a countable base so that ``cover'' means ``countable
cover,'' then run the nested-closed-sets contradiction: the tails of an unending cover leave a descending
chain of nonempty closed sets, a transversal through them is infinite, and its cluster point must lie in a
cover member that the transversal was constructed to avoid. It is the metric-space twin of the
finite-intersection characterization at \xrefL{4.2.1}, read in reverse---there nested compact sets are
forced to meet, here a would-be-non-compact space is forced to meet itself and contradict.
\end{Reflection}

% ------------------------------------------------------------
% ------------------------------------------------------------
\prob{27}{Rudin, Chapter 2, Exercise 27}
Define a point $p$ in a metric space $X$ to be a \emph{condensation point} of a set $E\subset X$ if
every neighborhood of $p$ contains uncountably many points of $E$.

Suppose $E\subset\mathbb R^{k}$, $E$ is uncountable, and let $P$ be the set of all condensation points
of $E$. Prove that $P$ is perfect and that at most countably many points of $E$ are not in $P$. In
other words, show that $P^{c}\cap E$ is at most countable. \emph{Hint:} Let $\{V_{n}\}$ be a countable
base of $\mathbb R^{k}$, let $W$ be the union of those $V_{n}$ for which $E\cap V_{n}$ is at most
countable, and show that $P=W^{c}$.

\begin{SolutionBlue}
First fix a countable base for $\mathbb R^{k}$: let
\[
\{V_n\}_{n\in\mathbb N}=\bigl\{\,B(q,r):q\in\mathbb Q^{k},\ r\in\mathbb Q,\ r>0\,\bigr\},
\]
the family of open balls with rational centre and rational radius. This is a countable family, being
indexed by $\mathbb Q^{k}\times\mathbb Q$, a finite product of countable sets. It is a \emph{base}: if
$U$ is open and $x\in U$, choose $\rho>0$ with $B(x,\rho)\subset U$, then a rational point $q$ with
$d(x,q)<\rho/2$ and a rational $r$ with $d(x,q)<r<\rho/2$; then $x\in B(q,r)$, and for any
$y\in B(q,r)$ one has $d(x,y)\le d(x,q)+d(q,y)<\rho/2+\rho/2=\rho$, so $B(q,r)\subset B(x,\rho)\subset U$.
Thus every open set is a union of members of $\{V_n\}$.

Let $W=\bigcup\{\,V_n:E\cap V_n\text{ is at most countable}\,\}$, the union of the ``small'' basic sets,
and we show $P=W^{c}$.

If $p\in W$, then $p\in V_n$ for some $n$ with $E\cap V_n$ at most countable. Since $V_n$ is open it is
a neighborhood of $p$ meeting $E$ in at most countably many points, so $p$ is not a condensation point;
that is, $p\notin P$. Hence $W\subset P^{c}$.

If $p\notin P$, some neighborhood of $p$ contains at most countably many points of $E$; shrinking, there
is $\rho>0$ with $E\cap B(p,\rho)$ at most countable. By the base property there is a basic set $V_n$
with $p\in V_n\subset B(p,\rho)$, and then $E\cap V_n\subset E\cap B(p,\rho)$ is at most countable, so
$V_n$ is one of the sets defining $W$ and $p\in V_n\subset W$. Hence $P^{c}\subset W$. Combining the two
inclusions, $P=W^{c}$.

Now $W$ is a union of open sets, hence open, so $P=W^{c}$ is closed. Moreover
\[
E\cap W=\bigcup\{\,E\cap V_n:E\cap V_n\text{ at most countable}\,\}
\]
is a union of at most countably many at-most-countable sets (there are only countably many basic sets in
all), hence at most countable. Since
\[
P^{c}\cap E=W\cap E
\]
this says exactly that at most countably many points of $E$ lie outside $P$.

It remains to see that every point of $P$ is a limit point of $P$. Fix $p\in P$ and any $\rho>0$. As $p$
is a condensation point, $E\cap B(p,\rho)$ is uncountable. Removing the at-most-countable set $E\cap W$
cannot exhaust an uncountable set, so
\[
E\cap B(p,\rho)\cap W^{c}=E\cap B(p,\rho)\cap P
\]
is uncountable; in particular it contains a point of $P$ other than $p$. Thus every ball about $p$
contains a point of $P$ distinct from $p$, so $p$ is a limit point of $P$. Being closed with no isolated
points, $P$ is perfect.
\end{SolutionBlue}

\begin{Reflection}
The statement defines its own vocabulary, and the first job is to translate ``condensation point'' into
a familiar shape. It is the cluster-point definition (\xrefL{3.2.1}) with ``infinitely many'' upgraded to
``uncountably many''---every neighbourhood of $p$ meets $E$ in an uncountable set. The conclusion asks
for two things at once, ``$P$ is perfect'' and ``$P^{c}\cap E$ is at most countable,'' and the second
clause is the louder clue: a countability bound on a topologically defined set is the signature of a
\emph{countable base}. The reason a base helps is that it converts the uncountably many neighbourhoods
hiding in the definition into a countable checklist, and the hint hands you exactly this device.

Believe the mechanism on a small picture before proving it. Take $E=[0,1]\cup\{2\}$ in $\mathbb R$. Every
point of $[0,1]$ is a condensation point, since any interval around it catches a whole subinterval, which
is uncountable; but the isolated point $2$ is not, as the neighbourhood $(1.5,2.5)$ meets $E$ only in
$\{2\}$. So $P=[0,1]$, which is perfect, and the single stray point $2$ is the entire countable
remainder $P^{c}\cap E$. The proof simply makes ``stray points are rare'' precise by collecting them
inside basic neighbourhoods that are individually thin.

The argument runs in stages, each resting on a named result:
\begin{enumerate}[leftmargin=1.7em,itemsep=3pt,topsep=2pt,label=\textnormal{(\arabic*)}]
  \item manufacture the countable base from rational-centre, rational-radius balls; it is countable
        because $\mathbb Q^{k}\times\mathbb Q$ is a finite product of countable sets
        (\factref{nzq-countable}{$\mathbb Q$ is countable}, \factref{finite-product}{finite products of
        countable sets are countable}), and it is a base because $\mathbb Q$ is dense
        (\factref{q-dense}{$\mathbb Q$ dense in $\mathbb R$}, \xrefL{1.7.2}) lets you slip such a ball
        between any point and any open neighbourhood, the same balls-are-open fact behind
        \factref{balls-open}{open balls are open} (\xrefL{2.3.5}) and the description of open sets as
        unions of balls (\xrefL{3.1.5});
  \item with $W$ the union of the basic sets meeting $E$ countably, identify $P=W^{c}$: a point in $W$
        sits in a thin basic neighbourhood, so it fails the condensation condition, while a point failing
        condensation has \emph{some} thin neighbourhood and the base supplies a thin basic set inside it;
  \item read off that $P$ is closed because $W$, a union of open sets, is open, so its complement is
        closed---the open/closed duality \factref{open-iff-closed}{open iff complement closed}
        (\xrefL{3.2.8}, \xrefL{3.2.4});
  \item bound $P^{c}\cap E=W\cap E$: it is the union of the countably many sets $E\cap V_n$ that defined
        $W$, each at most countable, so a countable union of countable sets is at most countable
        (\factref{countable-union}{countable union of countable sets}, \factref{subset-countable}{subsets
        of countable sets});
  \item prove every $p\in P$ is a limit point of $P$: each ball $B(p,\rho)$ meets $E$ uncountably, and
        deleting the at-most-countable $E\cap W$ leaves an uncountable---hence nonempty, and in fact
        infinite---set of points of $E\cap P$ near $p$, forcing a point of $P$ other than $p$ into the
        ball.
\end{enumerate}

The one inequality doing all the work in step (5) is that an uncountable set minus a countable set is
still uncountable; this is where ``$E$ uncountable'' is finally spent, and dropping it breaks everything
($P$ would be empty). The other hypotheses are equally load-bearing: without the countable base, step (4)
has no reason to produce a countable total---an uncountable union of thin sets need not be thin---which is
precisely why the construction lives in $\mathbb R^{k}$, where a countable base exists. The argument is
not circular: $P$ is built from $W$, and $W$ from the fixed base, with the condensation property only
read off afterward, never assumed.

The transferable move is the one the hint enforces: when a property is quantified over \emph{all}
neighbourhoods, replace ``all neighbourhoods'' by ``all basic neighbourhoods,'' and if the base is
countable an uncountable quantifier collapses to a countable bookkeeping problem. The same trick---small
basic sets union to a controllably small set---is what makes second-countable spaces so well-behaved, and
it is exactly the strengthening of the cluster-point idea of \xrefL{3.2.1} from ``infinitely many'' to
``uncountably many.''
\end{Reflection}

% ------------------------------------------------------------
% ------------------------------------------------------------
\prob{28}{Rudin, Chapter 2, Exercise 28}
Prove that every closed set in a separable metric space is the union of a (possibly empty) perfect set
and a set which is at most countable. (\emph{Corollary:} Every countable closed set in $\mathbb R^k$ has
isolated points.) \emph{Hint:} Use Exercise 27.

\begin{SolutionBlue}
Call a point $p$ a \emph{condensation point} of a set $E$ if every neighborhood of $p$ contains
uncountably many points of $E$; call a closed set \emph{perfect} if every one of its points is a cluster
point of it. Let $F$ be a closed subset of a separable metric space $X$, let $P$ be the set of
condensation points of $F$, and write $F\setminus P$ for the rest. The claim is that $P$ is perfect
(possibly empty), that $F\setminus P$ is at most countable, and that $F=P\cup(F\setminus P)$.

By \hyperlink{prob:23}{Problem~23} a separable metric space has a countable base $\{V_n\}_{n\ge 1}$. Let
$W$ be the union of those basic sets $V_n$ for which $F\cap V_n$ is at most countable. The first task,
exactly as in \hyperlink{prob:27}{Problem~27}, is to identify the leftover set with $W$:
\[
  P^c=W,\qquad\text{equivalently}\qquad P=W^c .
\]
If $p\in W$, then $p\in V_n$ for a basic set with $F\cap V_n$ at most countable, so $p$ has a neighborhood
meeting $F$ in at most countably many points and is not a condensation point. Conversely, if $p\notin P$,
some neighborhood $U$ of $p$ has $F\cap U$ at most countable; since $\{V_n\}$ is a base there is a $V_n$
with $p\in V_n\subseteq U$, whence $F\cap V_n\subseteq F\cap U$ is at most countable and $V_n$ is one of
the sets forming $W$, so $p\in W$.

$P$ is closed, being the complement of the open set $W$ (\xrefL{3.2.8};
\factref{open-iff-closed}{List of Facts}). The leftover $F\setminus P=F\cap W$ is at most countable: it is
contained in the union of the sets $F\cap V_n$ used to build $W$, a countable union of at-most-countable
sets, hence at most countable (\factref{countable-union}{List of Facts}, with
\factref{subset-countable}{``a subset of a countable set is finite or countable''}, \xrefL{2.1.9}).

The perfect part lies inside $F$. If $p\notin F$, then because $F$ is closed its complement $X\setminus F$
is open (\xrefL{3.2.8}) and is a neighborhood of $p$ disjoint from $F$, so $p$ meets $F$ in no points at
all and is not a condensation point. Thus $P\subseteq F$, and here the closedness of $F$ is spent.

$P$ is perfect. It is closed, so only the cluster-point condition remains. Fix $p\in P$ and any
neighborhood $U$ of $p$. Then $U\cap F$ is uncountable, since $p$ is a condensation point. Removing the at
most countable set $F\setminus P$ leaves
\[
  U\cap P=(U\cap F)\setminus(F\setminus P)
\]
still uncountable, hence infinite, hence containing some point of $P$ other than $p$. Every neighborhood of
$p$ therefore contains a point of $P$ distinct from $p$, so $p$ is a cluster point of $P$
(\xrefL{3.2.1}). As $p\in P$ was arbitrary, every point of $P$ is a cluster point of $P$, and $P$ is
perfect. Combining,
\[
  F=P\cup(F\setminus P),
\]
the union of a (possibly empty) perfect set and an at most countable set.

For the corollary, let $F\subseteq\mathbb R^k$ be nonempty, countable, and closed, and suppose toward a
contradiction that $F$ has no isolated points. Then every point of $F$ is a cluster point of $F$, and
since $F$ is closed it is therefore a nonempty perfect subset of $\mathbb R^k$. But a nonempty perfect
subset of $\mathbb R^k$ is uncountable. To see this, suppose such a set $P$ were the countable list
$\{p_1,p_2,\dots\}$. Pick a closed ball $\overline B_1$ centered at a point of $P$ with $p_1\notin\overline
B_1$ and $\overline B_1\cap P\neq\varnothing$; this is possible because $p_2$ (say) is a cluster point of
$P$, so points of $P$ accumulate near it and a small enough ball excludes the single point $p_1$.
Inductively, given a closed ball $\overline B_n$ with $\overline B_n\cap P\neq\varnothing$, its center is a
cluster point of $P$, so $\overline B_n$ contains a point $q$ of $P$ with $q\neq p_{n+1}$; choose a closed
ball $\overline B_{n+1}\subseteq \overline B_n$ centered at $q$, of radius small enough that
$p_{n+1}\notin\overline B_{n+1}$, with $\overline B_{n+1}\cap P\neq\varnothing$. The sets
$K_n=\overline B_n\cap P$ are nonempty; each is closed and bounded in $\mathbb R^k$, hence compact
(\xrefL{4.3.5}; \factref{heine-borel}{List of Facts}), and they are nested,
$K_1\supseteq K_2\supseteq\cdots$. By the nested-compact-sets property their intersection is nonempty
(\xrefL{4.2.2}; \factref{nested-compact}{List of Facts}), so some $q^\ast\in\bigcap_n K_n\subseteq P$.
But $q^\ast=p_m$ for some $m$, while $q^\ast\in K_m\subseteq\overline B_m$ was built to exclude $p_m$---a
contradiction. Hence $P$ is uncountable, contradicting that $F$ is countable. Therefore a nonempty
countable closed set in $\mathbb R^k$ must have an isolated point.
\end{SolutionBlue}

\begin{Reflection}
The clue is the single word \emph{perfect}, paired with the instruction to ``use Exercise~27.'' A perfect
set is a closed set all of whose points are cluster points (\xrefL{3.2.1}, \xrefL{3.2.4}); to split an
arbitrary closed $F$ into a perfect piece and a small remainder you need a canonical rule for deciding
which points of $F$ are ``robustly surrounded'' by $F$ and which are not. \hyperlink{prob:27}{Problem~27}
built exactly that rule---the condensation points, the points where $F$ stays \emph{uncountably} large in
every neighborhood---and showed they form a closed set $P$ whose complement captures all but at most
countably many points of $E$. Reading ``perfect set $+$ countable set'' as ``condensation points $+$ the
rest'' is the whole move; the present problem is \hyperlink{prob:27}{Problem~27} re-run on a closed set,
with one extra observation that closedness contributes.

The hidden hinge is the word \emph{separable}. Without a countable base there is no way to assemble the
small leftover, because the bound on $F\setminus P$ comes from there being only countably many basic sets
to fail on. \hyperlink{prob:23}{Problem~23} supplies the base $\{V_n\}$, and the union $W$ of the basic
sets that meet $F$ countably is the engine: it is open by construction, so its complement $P$ is closed
(\xrefL{3.2.8}; \factref{open-iff-closed}{List of Facts}), and $F\cap W$ is swept up inside a countable
union of countable sets (\factref{countable-union}{List of Facts}).

Believe the shape on a familiar set before checking it. Take $F=\{0\}\cup\{1/n:n\ge 1\}$ in $\mathbb R$, a
countable closed set. No point is a condensation point, so $P=\varnothing$ and $F=\varnothing\cup F$ is
the whole of it as the countable remainder---and the corollary is visible too, since $1$ is an isolated
point. For an uncountable example, $F=[0,1]$ has $P=[0,1]$ and the remainder empty. The split always sends
the ``thick'' part to $P$ and the ``thin'' part to the countable scrap.

Each step rests on a named support:
\begin{enumerate}[leftmargin=1.7em,itemsep=3pt,topsep=2pt,label=\textnormal{(\arabic*)}]
  \item the countable base from \hyperlink{prob:23}{Problem~23} lets you define $W$ as the union of the
        basic sets meeting $F$ in at most countably many points; the base property gives $P=W^c$, exactly
        as in \hyperlink{prob:27}{Problem~27}, because every neighborhood with countable $F$-trace contains
        a basic such neighborhood;
  \item $P$ is closed as the complement of the open $W$ (\xrefL{3.2.8};
        \factref{open-iff-closed}{List of Facts}), and $F\setminus P=F\cap W$ is at most countable, lying
        in a countable union of at-most-countable sets (\factref{countable-union}{List of Facts},
        \factref{subset-countable}{subset of a countable set}, \xrefL{2.1.9});
  \item closedness of $F$ forces $P\subseteq F$: a point outside the closed $F$ has the open complement
        $X\setminus F$ (\xrefL{3.2.8}) as a $F$-free neighborhood, so it cannot be a condensation point---
        this is the lone place ``$F$ closed'' is used, and it is what keeps the perfect piece honestly
        inside $F$;
  \item $P$ is perfect because near any $p\in P$ the neighborhood meets $F$ uncountably, and subtracting
        the at most countable $F\setminus P$ still leaves uncountably many points of $P$, in particular one
        $\neq p$, so $p$ is a cluster point of $P$ (\xrefL{3.2.1}).
\end{enumerate}

The rigor turns on the arithmetic of cardinalities in step (4): \emph{uncountable minus at-most-countable
is uncountable}. That is why \hyperlink{prob:27}{Problem~27} insists on \emph{condensation} points rather
than mere cluster points---an ordinary cluster point only guarantees \emph{infinitely} many nearby points
of $F$ (\factref{cluster-infinite}{List of Facts}, \xrefL{3.2.1}), and ``infinite minus countable'' can be
empty, so the leftover could swallow every nearby point and break the cluster-point check. The uncountable
local mass is exactly the surplus that survives the subtraction. Closedness of $F$ is load-bearing in step
(3): drop it and a condensation point of $F$ could sit on the boundary outside $F$, so $P$ would not lie
in $F$ and the decomposition $F=P\cup(F\setminus P)$ would fail. Nothing is circular: $P$ is defined from
$F$ and the base, never assumed perfect in advance.

The corollary is the payoff and shows why the ambient space matters. A nonempty countable closed set in
$\mathbb R^k$ cannot be perfect, because a nonempty perfect subset of $\mathbb R^k$ is uncountable---the
classical fact proved here by trapping a point of $P$ inside a nested chain of compact pieces
$K_n=\overline B_n\cap P$, each nonempty (the centers are cluster points), each closed and bounded hence
compact (\xrefL{4.3.5}; \factref{heine-borel}{List of Facts}), with nonempty intersection
(\xrefL{4.2.2}; \factref{nested-compact}{List of Facts}) engineered to dodge the $n$-th listed point. So
a nonempty countable closed set, having $P=\varnothing$, equals its countable remainder, and a set all of
whose points were cluster points would be perfect and so uncountable---hence some point must be isolated.

The transferable idea is the cleavage itself, the Cantor--Bendixson move: to peel a closed set into a
``thick'' core and a ``thin'' shell, define the core by a local-largeness condition stable under shrinking
neighborhoods (here, uncountably many nearby points), prove the core closed and the shell small with a
countable base, and then exploit that the largeness condition survives removing the shell. The same template
recovers the perfect kernel of any closed set and, in $\mathbb R^k$, separates the countable closed sets
(all isolated debris, no kernel) from the uncountable ones.
\end{Reflection}

% ------------------------------------------------------------
% ------------------------------------------------------------
\prob{29}{Rudin, Chapter 2, Exercise 29}
Prove that every open set in $\mathbb R^{1}$ is the union of an at most countable collection of disjoint
segments. \emph{Hint:} Use Exercise~22.

\begin{Solution}
Let $E\subseteq\mathbb R$ be open. For each $x\in E$ define
\[
  I_x=\bigcup\{\,J:\ J\text{ is an open interval with }x\in J\subseteq E\,\}.
\]
Because $E$ is open, at least one such $J$ exists, so $I_x$ is nonempty and contains $x$.

First, each $I_x$ is itself an open interval contained in $E$. It is open as a union of open sets, and
it lies in $E$ since every $J$ in the union does. To see it is an interval, take $u<v$ in $I_x$ and $w$
with $u<w<v$; say $u\in J_1$ and $v\in J_2$, where $J_1,J_2$ are intervals from the union, both
containing $x$. If $w\le x$ then $u<w\le x$ forces $w\in J_1$ (an interval containing both $u$ and $x$);
if $w\ge x$ then $x\le w<v$ forces $w\in J_2$. Either way $w\in I_x$, so $I_x$ is an interval. By
construction it is the \emph{largest} open interval with $x\in I_x\subseteq E$ (an infinite endpoint is
allowed).

Second, any two of these intervals are equal or disjoint. If $I_x\cap I_y\neq\varnothing$, then
$I_x\cup I_y$ is again an open interval contained in $E$ (two intervals that meet have an interval union),
and it contains $x$, so maximality of $I_x$ gives $I_x\cup I_y\subseteq I_x$, i.e.\ $I_y\subseteq I_x$;
symmetrically $I_x\subseteq I_y$, hence $I_x=I_y$. The distinct sets among the $I_x$ are therefore
pairwise disjoint, and since each $x\in E$ lies in its own $I_x$, their union is exactly $E$.

It remains to count them. List the rationals as a sequence $r_1,r_2,r_3,\dots$, which is possible because
$\mathbb Q$ is countable (\xrefL{2.1.7}; \factref{nzq-countable}{List of Facts}). Each of these intervals
is a nonempty open interval, so it contains a rational by density of $\mathbb Q$ in $\mathbb R$
(\xrefL{1.7.2}; \factref{q-dense}{List of Facts}); assign to each distinct interval $I$ the rational
$r_I$ of least index lying in $I$. Distinct intervals are disjoint, so they receive distinct rationals,
and $I\mapsto r_I$ is an injection from the collection of intervals into $\mathbb Q$. A set that injects
into the countable set $\mathbb Q$ is finite or countable (\factref{subset-countable}{a subset of a
countable set is finite or countable}). Thus $E$ is the union of an at most countable collection of
pairwise disjoint segments.
\end{Solution}

\begin{Reflection}
The phrase ``disjoint segments'' is the whole map. On the line a connected open set is precisely an open
interval, so ``write $E$ as disjoint open intervals'' is the same request as ``split $E$ into its
connected components'' (\xrefL{4.4.2})---and that single reframing tells you to manufacture, around each
point, the biggest interval-shaped chunk of $E$ that holds it. The second clue is ``at most countable,''
the standard countability verdict, which forces the question of \emph{how} to count an a-priori-unknown
family; and the hint, ``use Exercise~22,'' points at the device, since Exercise~22 is exactly where the
density of $\mathbb Q$ is harnessed to count---there to find a countable dense $\mathbb Q^k$, here to tag
each interval with a rational.

Believe the statement on a picture before proving it: $E=(0,1)\cup(2,5)\cup(7,\infty)$ is already three
disjoint segments, and any open $E$ you draw breaks into such pieces with gaps (points not in $E$)
between them. The content is that this always happens and that the pieces are never more than countably
many.

The proof has three movements, each resting on the one before or on a named result:
\begin{enumerate}[leftmargin=1.7em,itemsep=3pt,topsep=2pt,label=\textnormal{(\arabic*)}]
  \item around each $x\in E$ form $I_x$, the union of \emph{all} open intervals trapped between $x$ and
        $E$; it is open as a union of opens (\factref{union-open}{any union of open sets is open}), it is
        an interval because all the pieces share the anchor $x$ (splitting on whether the in-between point
        $w$ sits left or right of $x$, a clean two-case check, \xrefL{0.5.7}), and by construction it is
        the maximal such interval---this is the component of $x$;
  \item maximality forces a dichotomy: if $I_x$ and $I_y$ overlap, their union is a single larger
        interval inside $E$ containing $x$, which maximality cannot allow unless $I_x=I_y$, so distinct
        components are disjoint and, as each point lies in its own, they cover $E$;
  \item to count, send each distinct interval to the least-indexed rational it contains; density of
        $\mathbb Q$ (\factref{q-dense}{$\mathbb Q$ is dense in $\mathbb R$}) guarantees there is one, and
        disjointness makes the assignment injective, so the family injects into the countable $\mathbb Q$
        and is at most countable (\factref{subset-countable}{a subset of a countable set is finite or
        countable}).
\end{enumerate}

Each hypothesis earns its place. Openness of $E$ is what lets step~(1) start---without it a point of $E$
need sit in no interval inside $E$, and $I_x$ could be empty or a single point, not a segment. The
maximality built in step~(1) is the sole engine of the disjointness in step~(2): drop it and overlapping
intervals would double-count the same stretch of $E$, so the family would neither be disjoint nor inject
into $\mathbb Q$. And the counting in step~(3) leans on disjointness, not on the intervals being long; an
interval of any positive length still catches a rational, so even a swarm of tiny components stays
countable. The argument is not circular: the components are constructed from $E$ outward, and only after
they are pinned down are they counted.

The transferable move is the rational-tagging trick, and it is worth abstracting past intervals: any
family of pairwise disjoint nonempty open subsets of $\mathbb R$ (and of $\mathbb R^k$ too, using
$\mathbb Q^k$ as in Exercise~22) is at most countable, because each open set swallows a point of the
countable dense set and disjointness keeps those points distinct. Whenever a problem hands you a disjoint
family of fat sets and asks ``how many?'', reach for a countable dense set and let disjointness do the
injecting.
\end{Reflection}

% ------------------------------------------------------------
% ------------------------------------------------------------
\prob{30}{Rudin, Chapter 2, Exercise 30}
Imitate the proof of Theorem 2.43 to obtain the following result: if $\mathbb R^{k}=\bigcup_{n=1}^{\infty}F_{n}$, where each $F_{n}$ is a closed subset of $\mathbb R^{k}$, then at least one $F_{n}$ has a nonempty interior. \emph{Equivalent statement:} if $G_{n}$ is a dense open subset of $\mathbb R^{k}$ for $n=1,2,3,\ldots$, then $\bigcap_{n=1}^{\infty}G_{n}$ is not empty (in fact, it is dense in $\mathbb R^{k}$). \emph{(This is a special case of Baire's theorem.)}

\begin{SolutionBlue}
The two formulations are complementary, and it is the dense-open form that carries the construction; the closed-set form follows by taking complements. So let $G_1,G_2,\dots$ be dense open subsets of $\mathbb R^k$.

It suffices to show that $\bigcap_{n\ge 1}G_n$ meets every nonempty open set, which is exactly density (\xrefL{4.4.1}) and in particular forces nonemptiness. Fix a nonempty open $U\subseteq\mathbb R^k$. The plan is to manufacture a nested sequence of closed balls, each one trapped inside both the previous ball and the next $G_n$, and then catch a common point with the nested-compact-set property.

Since $G_1$ is dense, $U\cap G_1$ is nonempty, and since $G_1$ is open it is open; pick a point in it, and because it is open choose a radius small enough that the \emph{closed} ball about that point still lies inside, with radius less than $1$:
\[
  K_1=\overline B(x_1,r_1)\subseteq U\cap G_1,\qquad 0<r_1<1 .
\]
(One can always shrink: if $B(p,\rho)\subseteq W$ with $W$ open, then $\overline B(p,\rho/2)\subseteq B(p,\rho)\subseteq W$.)

Suppose $K_{n-1}=\overline B(x_{n-1},r_{n-1})$ has been chosen. The open ball $B(x_{n-1},r_{n-1})$ is nonempty and open (\xrefL{2.3.1}), and $G_n$ is dense and open, so $B(x_{n-1},r_{n-1})\cap G_n$ is nonempty and open. Shrinking a closed ball into it as before, choose
\[
  K_n=\overline B(x_n,r_n)\subseteq B(x_{n-1},r_{n-1})\cap G_n,\qquad 0<r_n<\tfrac1n .
\]
Each $K_n$ is closed and bounded in $\mathbb R^k$, hence compact by Heine--Borel (\xrefL{4.3.5}; \factref{heine-borel}{List of Facts}), and each is nonempty. Because $K_n\subseteq B(x_{n-1},r_{n-1})\subseteq K_{n-1}$, the balls are nested:
\[
  K_1\supseteq K_2\supseteq K_3\supseteq\cdots .
\]
A nested sequence of nonempty compact sets has nonempty intersection (\xrefL{4.2.2}; \factref{nested-compact}{List of Facts}), so there is a point $x\in\bigcap_{n\ge 1}K_n$. From $K_1\subseteq U$ we get $x\in U$, and from $K_n\subseteq G_n$ for every $n$ we get $x\in G_n$ for every $n$; hence
\[
  x\in U\cap\bigcap_{n=1}^{\infty}G_n .
\]
Thus $\bigcap_{n\ge 1}G_n$ meets every nonempty open $U$, so it is dense in $\mathbb R^k$, and being dense it is nonempty.

Now the closed-set form. Suppose $\mathbb R^k=\bigcup_{n\ge 1}F_n$ with each $F_n$ closed, and suppose toward a contradiction that every $F_n$ has empty interior. Put $G_n=\mathbb R^k\setminus F_n$. Each $G_n$ is open because $F_n$ is closed (\xrefL{3.2.8}; \factref{open-iff-closed}{List of Facts}), and $G_n$ is dense: a point $p$ with some ball $B(p,\rho)\subseteq F_n$ would be an interior point of $F_n$, against the assumption, so every ball about every point meets $G_n$. By the result just proved $\bigcap_{n\ge 1}G_n\neq\varnothing$. But
\[
  \bigcap_{n=1}^{\infty}G_n=\mathbb R^k\setminus\bigcup_{n=1}^{\infty}F_n=\mathbb R^k\setminus\mathbb R^k=\varnothing ,
\]
a contradiction. Hence at least one $F_n$ has nonempty interior.
\end{SolutionBlue}

\begin{Reflection}
The instruction ``imitate the proof of Theorem 2.43'' is the loudest clue in the problem: that theorem (perfect sets in $\mathbb R^k$ are uncountable) runs on a single engine, a shrinking tower of closed balls pinned down at the end by the fact that nested nonempty compact sets meet (\xrefL{4.2.2}; \factref{nested-compact}{List of Facts}). So before reading another word you should expect to build $K_1\supseteq K_2\supseteq\cdots$ and harvest a point from the intersection. The only question is what each ball must dodge, and the two hypotheses answer it---``open'' is what lets you fit a ball inside at all, and ``dense'' is what guarantees the next $G_n$ reaches into whatever ball you have so far.

The other half of reading the problem is recognising that ``$F_n$ closed with empty interior'' and ``$G_n=F_n^c$ dense and open'' are the same statement seen from opposite sides. Complementation turns closed into open (\xrefL{3.2.8}; \factref{open-iff-closed}{List of Facts}), turns the union $\bigcup F_n$ into the intersection $\bigcap G_n$, and---this is the part worth pausing on---turns ``$F_n$ has empty interior'' into ``$G_n$ is dense,'' because a point fails to be in the closure of $F_n^c$ exactly when some ball around it avoids $F_n^c$, i.e.\ sits inside $F_n$ as an interior point (\xrefL{4.4.1}). Once you see this dictionary, you prove whichever form is easier and read the other off for free.

The construction itself is a short loop, each step leaning on the one before:
\begin{enumerate}[leftmargin=1.7em,itemsep=3pt,topsep=2pt,label=\textnormal{(\arabic*)}]
  \item density of $G_1$ makes $U\cap G_1$ nonempty and openness keeps it open, so it contains a point with a whole ball around it; passing to a closed ball of half the radius gives $K_1\subseteq U\cap G_1$, the base of the tower;
  \item at stage $n$ the \emph{open} ball $B(x_{n-1},r_{n-1})$ is a legitimate nonempty open target (\xrefL{2.3.1}), so density of $G_n$ lets it reach in and openness lets you again shrink to a closed ball $K_n\subseteq B(x_{n-1},r_{n-1})\cap G_n$---and one shrinks into the open ball rather than the closed ball $K_{n-1}$ precisely because the shrinking step needs an \emph{open} set to seat a closed ball inside, $B(x_{n-1},r_{n-1})\cap G_n$ being open while $K_{n-1}\cap G_n$ need not be; nesting $K_n\subseteq K_{n-1}$ then comes for free from $B(x_{n-1},r_{n-1})\subseteq\overline B(x_{n-1},r_{n-1})=K_{n-1}$;
  \item each $K_n$ is closed and bounded, hence compact in $\mathbb R^k$ by Heine--Borel (\xrefL{4.3.5}; \factref{heine-borel}{List of Facts})---this is where ``$\mathbb R^k$'' is spent, since a closed ball need not be compact in a general metric space;
  \item nested nonempty compacts meet (\xrefL{4.2.2}; \factref{nested-compact}{List of Facts}), producing $x$; then $K_1\subseteq U$ and $K_n\subseteq G_n$ place $x$ in $U\cap\bigcap_n G_n$, so the intersection meets the arbitrary open set $U$ and is therefore dense (\xrefL{4.4.1}).
\end{enumerate}

It is worth checking that the radii $r_n<1/n$ are decoration here, not load: nested nonempty compact sets already meet without any control on size, so the shrinking radii buy nothing for this problem. They are kept only to mirror Theorem 2.43, where the radii must vanish so that the intersection is a single point and the branching produces distinct limits. The genuinely load-bearing hypotheses are the three the construction actually consumes---openness (to fit a ball), density (to reach the next set), and the completeness of $\mathbb R^k$ packaged as Heine--Borel and the nested-compact property. Drop completeness and the theorem fails: $\mathbb Q$ is a countable union of its (closed, empty-interior) singletons, so the closed-set form is false there, and the matching dense open sets $\mathbb Q\setminus\{q\}$ have empty intersection.

The transferable move is the one the hint is really teaching: when you must produce a point lying in infinitely many sets at once, do not seek it directly---trap it. Nest a tower of compact sets so that membership in the $n$-th set is forced by the $n$-th ball, and let the nested-compact-set property (\xrefL{4.2.2}) hand you the common point. This is the skeleton of every Baire-category argument, and it is the same skeleton our notes use to extract an uncountable perfect set.
\end{Reflection}

\end{document}
